%%============================================================================
% APPENDIX: Chapter 3 --- Research Methods Supplementary Material
%%============================================================================
\setchaptercolor{3}

\chapter*{Appendix: Chapter 3 --- Research Methods Supplementary Material}
%% TOC entry is in main.tex (Appendix C label)
\label{chap:appendix_ch3}
\phantomsection\label{app:ch3_support_appx}
\normalsize\normalfont\color{black}
\phantomsection\label{app:ch3_support}

% Global overflow prevention for appendix tables: allow hyphenation in p{} columns
\renewcommand{\tabfixsetup}{\sloppy\hyphenpenalty=0\exhyphenpenalty=0\tolerance=9999\emergencystretch=8em\hfuzz=100pt\hbadness=9999}
% Stub labels for suppressed content references (phantom — resolve to appendix start)
\label{fig:pii_flow}
\label{tab:mcp_capabilities}

\noindent\textbf{Appendix C Scope Contract --- Conceptual Methodological Supplement Only.}
\label{para:appc_scope_contract}
This appendix supplements Chapter~3 methodology with conceptual artefacts (multi-agent architecture references, governance-gate references, MCP capability tables, sample workflow diagrams). \textbf{It is NOT}: a deployed system specification; a regulatory submission artefact; a clinical-implementation guide; or a future-operational-validation roadmap. Where the appendix references production monitoring, governance gates, deployment architecture, or wearable-monitoring components, these are \emph{conceptual reference frameworks supporting the methodological narrative} --- not validated implementations. Any future operational realisation would require separate IRB review, prospective validation, and regulatory assessment --- explicitly out of scope. \textbf{Prior IRB approval references} for the $n{=}131$ PLS-SEM and $n{=}12$ expert-panel data streams cited from Chapters~3 and~4 are consolidated in Chapter~\ref{chap:methodology} Section~\ref{para:prior_irb_approval_trail} (Prior IRB Approval Trail).

This appendix contains supplementary tables, figures, and detailed analyses that support Chapter~3 but were moved from the main text to maintain narrative focus. Each item is cross-referenced from the main chapter.

%% --- Task Flow Diagram ---
\needspace{40\baselineskip}
\begin{figure}[!htbp]
\centering
\resizebox{0.97\textwidth}{0.75\textheight}{%
\begin{tikzpicture}[
    node distance=0.8cm,
    item/.style={rectangle, draw=ggunavy, fill=ggunavy!5, rounded corners=2pt,
        text width=13.5cm, minimum height=0.6cm, font=\small, align=left, inner sep=4pt, thick},
    arrow/.style={-{Stealth[length=3mm, width=2mm]}, thick, ggunavy}
]
\node[item, fill=lightblue] (s1) {\textbf{S1.} Redundant EEG Signal Processing Pipeline (fig:eeg\_\allowbreak signal\_\allowbreak pipeline, tab:eeg\_\allowbreak filter\_\allowbreak p...};
\node[item, fill=lightorange, below=of s1] (s2) {\textbf{S2.} Archived non-ASD GAN/PBS material retained only as a placeholder note in the ASD build};
\node[item, fill=lightteal, below=of s2] (s3) {\textbf{S3.} Archived comparative dual-track preprocessing material replaced with an ASD-build note};
\node[item, fill=lightpurple, below=of s3] (s4) {\textbf{S4.} Redundant EEG Feature Taxonomy table (tab:eeg\_\allowbreak feature\_\allowbreak categories) — overlaps with tab...};
\node[item, fill=lightgreen, below=of s4] (s5) {\textbf{S5.} Wearable Sensor Integration (tab:wearable\_\allowbreak sensors, fig:edge\_\allowbreak deployment) — no wearable...};
\node[item, fill=lightyellow, below=of s5] (s6) {\textbf{S6.} Data Risk Assessment Matrix (fig:data\_\allowbreak risk\_\allowbreak matrix) — redundant with heatmap (fig:data...};
\node[item, fill=lightpink, below=of s6] (s7) {\textbf{S7.} MCP Governance Capabilities table (tab:mcp\_\allowbreak capabilities) and Flow diagram (fig:mcp\_\allowbreak flow)};
\node[item, fill=lightcoral, below=of s7] (s8) {\textbf{S8.} NeuroMCP Multi-Agent Architecture diagram (fig:neuromcp\_\allowbreak layers) and Tool Registry tabl...};
\node[item, fill=lightsky, below=of s8] (s9) {\textbf{S9.} Guardrail AI input/output safety rules tables (tab:ch3\_\allowbreak input\_\allowbreak guardrails, tab:ch3\_\allowbreak out...};
\node[item, fill=lightmint, below=of s9] (s10) {\textbf{S10.} PII Protection Pipeline (tab:pii\_\allowbreak pipeline) and Security Alignment (tab:security\_\allowbreak align...};
\node[item, fill=lightlavender, below=of s10] (s11) {\textbf{S11.} 13-Phase GenAI QA Framework overview table (tab:13phase\_\allowbreak overview) — audit checklist de...};
\node[item, fill=lightblue!50, below=of s11] (s12) {\textbf{S12.} 8-Phase DL Analysis Framework overview table (tab:8phase\_\allowbreak overview) — audit checklis...};
\node[item, fill=lightorange!50, below=of s12] (s13) {\textbf{S13.} 10 Pillars of Cross-Cutting AI Governance overview table (tab:10pillars) — audit checkl...};
\node[item, fill=lightteal!50, below=of s13] (s14) {\textbf{S14.} Responsible AI Governance Summary table (tab:responsible\_\allowbreak ai\_\allowbreak summary) — audit-level sp...};
\draw[arrow] (s1) -- (s2);
\draw[arrow] (s2) -- (s3);
\draw[arrow] (s3) -- (s4);
\draw[arrow] (s4) -- (s5);
\draw[arrow] (s5) -- (s6);
\draw[arrow] (s6) -- (s7);
\draw[arrow] (s7) -- (s8);
\draw[arrow] (s8) -- (s9);
\draw[arrow] (s9) -- (s10);
\draw[arrow] (s10) -- (s11);
\draw[arrow] (s11) -- (s12);
\draw[arrow] (s12) -- (s13);
\draw[arrow] (s13) -- (s14);
\end{tikzpicture}%
}%
\caption[Appendix Task Flow]{Appendix Task Flow: Supplementary Material for Chapter~3 (Research Methods)}
\label{fig:appendix_ch3_taskflow}
\end{figure}

%%----------------------------------------------------------------------------
% S1: Redundant EEG Signal Processing Pipeline (fig:eeg_signal_pipeline, tab:eeg_filter_params) — overlaps with fig:eeg_preproc
%%----------------------------------------------------------------------------
% [SPACE-FIX] clearpage removed to eliminate empty page


\subsubsection{EEG Signal Processing Pipeline}
\label{sec:eeg_signal_pipeline}

The EEG-RAG framework~\cite{asthana2025eegmaster} implements a seven-stage signal processing pipeline that transforms raw multi-channel EEG recordings into classification-ready feature vectors. Figure~\ref{fig:eeg_signal_pipeline} illustrates this pipeline, and Table~\ref{tab:eeg_filter_params} documents the filtering parameters with their neurophysiological justification. The pipeline design reflects two constraints:

\begin{enumerate}[leftmargin=*, itemsep=3pt]
    \item \textbf{Preserve clinically relevant content:} Retain oscillatory content in the 0.5--45~Hz range while removing non-neural contamination such as DC drift, muscle artifacts, and power-line interference.
    \item \textbf{Produce classification-ready segments:} Generate fixed-length, normalised segments suitable for both classical feature extraction and deep learning input.
\end{enumerate}
\needspace{24\baselineskip}
\begin{figure}[!htbp]
\centering
\resizebox{0.97\textwidth}{!}{%
\begin{tikzpicture}[
  pnode/.style={rectangle, draw=ggunavy, fill=ggunavy!8, rounded corners=3pt,
    text width=3.8cm, minimum height=1.4cm, align=center, font=\small, inner sep=5pt},
  arr/.style={-{Stealth[length=4mm, width=3mm]}, very thick, ggunavy},
  note/.style={font=\small\itshape, text=gray, align=center},
  node distance=0.8cm
]
% Row 1
\node[pnode, fill=lightblue] (raw) {\textbf{\color{ggunavy}Raw EEG}\\[2pt]Multi-channel\\128--512 Hz};
\node[pnode, fill=lightorange, right=of raw] (bp) {\textbf{\color{ggunavy}Band-pass Filter}\\[2pt]0.5--45 Hz\\Butterworth 4th order};
\node[pnode, fill=lightteal, right=of bp] (notch) {\textbf{\color{ggunavy}Notch Filter}\\[2pt]50/60 Hz\\IIR notch};
\node[pnode, fill=lightpurple, right=of notch] (ica) {\textbf{\color{ggunavy}ICA Artifact}\\[2pt]FastICA\\Ocular + muscular};

% Row 2
\node[pnode, fill=lightgreen, below=0.6cm of ica] (norm) {\textbf{\color{ggunavy}Z-score}\\[2pt]Per-channel\\normalisation};
\node[pnode, fill=lightyellow, left=of norm] (win) {\textbf{\color{ggunavy}Windowing}\\[2pt]4 s epochs\\50\% overlap};
\node[pnode, fill=lightpink, left=of win] (feat) {\textbf{\color{ggunavy}Feature Extraction}\\[2pt]127 features\\per channel};

% Arrows
\draw[arr] (raw) -- (bp);
\draw[arr] (bp) -- (notch);
\draw[arr] (notch) -- (ica);
\draw[arr] (ica.south) -- ++(0,-0.15) -| (norm.north);
\draw[arr] (norm) -- (win);
\draw[arr] (win) -- (feat);

% Notes
\node[note, below=0.08cm of raw] {14--64 channels};
\node[note, below=0.08cm of bp] {Remove DC drift};
\node[note, below=0.08cm of notch] {Remove line noise};
\node[note, below=0.08cm of ica] {Remove blinks};
\end{tikzpicture}%
}
\caption[Seven-Stage EEG Signal Processing Pipeline]{Seven-Stage EEG Signal Processing Pipeline from the EEG-RAG Framework}
\label{fig:eeg_signal_pipeline}
\vspace{0.5em}\noindent\textit{Note.} ICA = independent component analysis; IIR = infinite impulse response; Hz = hertz; s = seconds. The pipeline produces fixed-length, normalised feature vectors suitable for both VotingClassifier and deep learning architectures. \textit{Source:}~\cite{asthana2025eegmaster}.
\end{figure}

\paragraph{EEG Filtering Parameters and.}



\noindent This table lists each \textit{filter} across type, parameters, and purpose, so examiners can trace what was assessed and what evidence supports each row.

{\tablestripes\tablefontsize
\needspace{20\baselineskip}
\begin{xltabular}{\textwidth}{@{} >{\raggedright\arraybackslash}p{1.8cm} >{\raggedright\arraybackslash}p{2.8cm} >{\raggedright\arraybackslash}p{3.6cm} L @{}}
\caption[EEG Filtering Parameters and Justification]{EEG Filtering Parameters: Type, Specifications, and Neurophysiological Justification}\label{tab:eeg_filter_params} \\
\toprule
\thead{Filter} & \thead{Type} & \thead{Parameters} & \thead{Purpose} \\
\midrule
\endfirsthead
\endhead
\bottomrule
\endlastfoot
High-pass & Butterworth 4th order & $f_c = 0.5$ Hz & Remove DC drift and slow electrode potential changes that obscure neural oscillations \\
\addlinespace
Low-pass & Butterworth 4th order & $f_c = 45$ Hz & Remove electromyographic (muscle) artifacts above the gamma band upper limit \\
\addlinespace
Notch & IIR & $f_0 = 50/60$ Hz; $Q = 30$ & Remove power-line interference (50 Hz in Europe/Asia; 60 Hz in North America) \\
\addlinespace
Band-pass & Butterworth 4th order & $f_L = 0.5$ Hz; $f_H = 45$ Hz & Combined high-pass and low-pass; retains all five canonical EEG bands ($\delta$ through $\gamma$) \\
\end{xltabular}}
\vspace{0.5em}\noindent\textit{Note.} $f_c$ = cutoff frequency; $f_0$ = centre frequency; $Q$ = quality factor; $f_L$ = lower cutoff; $f_H$ = upper cutoff; IIR = infinite impulse response. Butterworth design chosen for maximally flat passband response; 4th-order provides 24 dB/octave roll-off. Parameters follow~\cite{asthana2025eegmaster}.

%%----------------------------------------------------------------------------
% S2: GAN Architecture diagram (fig:gan_pipeline) and model params table (tab:gan_model_params) — implementation detail
%%----------------------------------------------------------------------------
% MOVED TO CHAPTER 3 MAIN TEXT (Mar 2026)
\iffalse
\clearpage


\subsubsection{GAN Architecture for PBS Cancer Classification}
\label{sec:gan_cancer_architecture}

The cancer detection pipeline~\cite{asthana2025cancer} employs a generative adversarial network (GAN) architecture coupled with a ResNet-18 classifier for four-class acute lymphoblastic leukaemia (ALL) classification. Figure~\ref{fig:gan_pipeline} illustrates the end-to-end pipeline from input PBS images through GAN-based augmentation to the final four-class output. The GAN serves a dual purpose:

\begin{enumerate}[leftmargin=*, itemsep=3pt]
    \item \textbf{Data augmentation:} Address class imbalance (the Pro-B class has 6\% more samples than Early Pre-B) by synthesising realistic cell images via the generator network.
    \item \textbf{Auxiliary classification signals:} The discriminator's learned features provide supplementary classification signals that improve the final classifier's decision boundary.
\end{enumerate}

\needspace{24\baselineskip}
\begin{figure}[!htbp]
\centering
\resizebox{0.97\textwidth}{!}{%
\begin{tikzpicture}[
  pbox/.style={rectangle, draw=ggunavy, fill=ggunavy!8, rounded corners=3pt,
    text width=2.8cm, minimum height=1.6cm, align=center, font=\small, inner sep=5pt},
  gbox/.style={rectangle, draw=ggunavy, fill=ggunavy!18, rounded corners=3pt,
    text width=3.2cm, minimum height=1.6cm, align=center, font=\small, inner sep=5pt},
  obox/.style={rectangle, draw=ggunavy, fill=ggunavy!30, rounded corners=3pt,
    text width=2.6cm, minimum height=1.6cm, align=center, font=\small\bfseries, inner sep=5pt},
  arr/.style={-{Stealth[length=4mm, width=3mm]}, very thick, ggunavy},
  node distance=0.8cm
]
% Row 1: Input processing
\node[pbox, fill=lightblue] (input) {\textbf{\color{ggunavy}Input PBS}\\[2pt]20,000 images\\512$\times$512 RGB};
\node[pbox, fill=lightorange, right=of input] (resize) {\textbf{\color{ggunavy}Resize}\\[2pt]512$\times$512\\standardise};
\node[pbox, fill=lightteal, right=of resize] (norm) {\textbf{\color{ggunavy}Normalise}\\[2pt]$[0,1]$ range\\HSV threshold};

% Row 1: GAN
\node[gbox, fill=lightpurple, right=of norm] (gan) {\textbf{\color{ggunavy}GAN}\\[2pt]Generator:\\151,808 params\\Discriminator:\\262,297 params};

% Row 2: Classification
\node[pbox, fill=lightgreen, right=of gan] (resnet) {\textbf{\color{ggunavy}ResNet-18}\\[2pt]Classifier\\6,149,220 params};
\node[obox, fill=lightyellow, right=of resnet] (output) {\textbf{\color{ggunavy}4-Class Output}\\[2pt]Benign\\Early Pre-B\\Pre-B, Pro-B};

% Arrows
\draw[arr] (input) -- (resize);
\draw[arr] (resize) -- (norm);
\draw[arr] (norm) -- (gan);
\draw[arr] (gan) -- (resnet);
\draw[arr] (resnet) -- (output);
\end{tikzpicture}%
}
\caption{GAN-Based PBS Cancer Classification Pipeline}
\label{fig:gan_pipeline}
\vspace{0.5em}\noindent\textit{Note.} PBS = peripheral blood smear; GAN = generative adversarial network; RGB = red--green--blue; HSV = hue--saturation--value. The generator uses ConvTranspose2d layers to synthesise realistic cell images; the discriminator provides auxiliary classification features. Total pipeline: 6,563,325 trainable parameters. \textit{Source:}~\cite{asthana2025cancer}.
\end{figure}

Table~\ref{tab:app_gan_model_params} documents the parameter counts for each pipeline component. The ResNet-18 classifier accounts for 93.7\% of the total parameters, reflecting the transfer learning strategy where ImageNet-pretrained convolutional features are fine-tuned on PBS microscopy images. The GAN generator and discriminator together comprise only 414,105 parameters (6.3\%), making the augmentation component computationally lightweight.

\noindent Table~\ref{tab:app_gan_model_params} gan cancer classification pipeline: model component parameters.

{\tablestripes\tablefontsize
\needspace{20\baselineskip}
\begin{xltabular}{\textwidth}{@{} >{\raggedright\arraybackslash}p{2.9cm} p{2.5cm} >{\raggedright\arraybackslash}p{2.0cm} p{6.8cm} @{}}
\caption[GAN Cancer Classification Pipeline]{GAN Cancer Classification Pipeline: Model Component Parameters}\label{tab:app_gan_model_params} \\
\toprule
\thead{Component} & \thead{Parameters} & \thead{Share (\%)} & \thead{Function} \\
\midrule
\endfirsthead
\endhead
\bottomrule
\endlastfoot
Generator & 151,808 & 2.3 & Synthesise realistic PBS cell images via ConvTranspose2d upsampling layers to augment training data \\
\addlinespace
Discriminator & 262,297 & 4.0 & Distinguish real from generated images; provide auxiliary classification features for the downstream classifier \\
\addlinespace
ResNet-18 classifier & 6,149,220 & 93.7 & Four-class ALL classification using ImageNet-pretrained features fine-tuned on PBS microscopy \\
\midrule
\textbf{Total} & \textbf{6,563,325} & \textbf{100.0} & --- \\
\end{xltabular}}
\vspace{0.5em}\noindent\textit{Note.} ALL = acute lymphoblastic leukaemia; PBS = peripheral blood smear; GAN = generative adversarial network. Parameter counts from~\cite{asthana2025cancer}. The generator uses five ConvTranspose2d layers (64$\to$512 channels); the discriminator mirrors this with five Conv2d layers. ResNet-18 is initialised with ImageNet weights and fine-tuned with a learning rate of $10^{-4}$.
\fi

%%----------------------------------------------------------------------------
% S3: Redundant Dual-Track Preprocessing Pipeline diagram (fig:preproc_pipeline) — overlaps with fig:preproc_comparison
%%----------------------------------------------------------------------------
% [SPACE-FIX] clearpage removed to eliminate empty page



\noindent\textbf{Appendix C Objective.} Provide detailed survey instruments, EEG pipeline flowcharts, clinical assessment forms, and methodology tables that support Chapter~3 without cluttering the main text.

\needspace{24\baselineskip}
\noindent The table below presents the appendix C Roadmap: Contents and Purpose.

\begin{table}[!htbp]
\centering
\caption{Appendix C Roadmap: Contents and Purpose.}
\tablestripes\tablefontsize
\begin{tabularx}{\textwidth}{@{} >{\raggedright\arraybackslash}p{3.0cm} L @{}}
\toprule
\thead{Content} & \thead{What It Contains} \\
\midrule
Survey instruments & S1--S4 full item tables (20 items each) \\\\\addlinespace
Clinical forms & ASD, PD, Epilepsy assessment forms \\\\\addlinespace
EEG flowcharts & Phase 1--5 pipeline diagrams \\\\\addlinespace
Variable tables & Questionnaire variable mapping, research variables \\\\\addlinespace
Planning tables & Chapter alignment, SMART matrix, decision flow \\\\\addlinespace
\bottomrule
\end{tabularx}
\end{table}


\subsection{Archived Comparative Preprocessing Note}

\phantomsection\label{fig:preproc_pipeline}
The earlier dual-track preprocessing diagram combined the active ASD EEG workflow with an archived non-ASD imaging branch. That comparative figure is not used in the ASD-only thesis build and is therefore omitted here. The active preprocessing logic retained in this thesis is the EEG-based ASD pipeline described in Chapter~3, including signal quality checks, band-pass filtering, ICA-based artifact removal, epoching, normalisation, and feature extraction.

\noindent\textit{Note.} The omitted comparative branch previously illustrated non-ASD image processing and is intentionally excluded from the ASD build to avoid ASD-focused ambiguity.

%%----------------------------------------------------------------------------
% S4: Redundant EEG Feature Taxonomy table (tab:eeg_feature_categories) — overlaps with tab:feature_127 and tab:feature_eng
%%----------------------------------------------------------------------------
% [SPACE-FIX] clearpage removed to eliminate empty page


\subsubsection{Comprehensive EEG Feature Extraction Taxonomy}
\label{sec:eeg_feature_taxonomy}

Table~\ref{tab:eeg_feature_categories} presents the five-category feature extraction taxonomy used in the ASD EEG pipeline. Each category captures a distinct aspect of neural dynamics relevant to ASD screening and interpretability:

\begin{itemize}[leftmargin=*, itemsep=3pt]
    \item \textbf{Time-domain features:} Characterise signal morphology through statistical descriptors and Hjorth parameters.
    \item \textbf{Frequency-domain features:} Capture oscillatory power distributions across the five canonical EEG bands.
    \item \textbf{Time-frequency features:} Reveal non-stationary spectral changes via wavelet and short-time Fourier transforms.
    \item \textbf{Connectivity features:} Quantify inter-channel synchronisation using coherence and phase-locking metrics.
    \item \textbf{Nonlinear features:} Measure signal complexity and chaotic behaviour through entropy and fractal dynamics.
\end{itemize}

\noindent Together, these five categories produce the 127 features per channel documented in Table~\ref{tab:feature_127}.
\needspace{24\baselineskip}
\noindent Table~\ref{tab:eeg_feature_categories} categorises comprehensive EEG Feature Extraction Categories: Methods, Examples, and Clinical Relevance.

\paragraph{EEG Feature Extraction Categories.}




{\tablestripes\tablefontsize

\needspace{20\baselineskip}
\begin{xltabular}{\textwidth}{@{} >{\raggedright\arraybackslash}p{2.0cm} >{\raggedright\arraybackslash}p{2.5cm} p{5.2cm} p{4.5cm} @{}}
\caption[EEG Feature Extraction Categories]{Comprehensive EEG Feature Extraction Categories: Methods, Examples, and Clinical Relevance}\label{tab:eeg_feature_categories} \\
\toprule
\thead{Category} & \thead{Feature Family} & \thead{Specific Features} & \thead{Clinical Relevance} \\
\midrule
\endfirsthead
\endhead
\bottomrule
\endlastfoot
Time domain & Statistical, Hjorth params & Mean, variance, skewness, kurtosis, zero-crossing rate, line length, RMS, peak-to-peak; activity, mobility, complexity & Signal morphology and behavioural-state variation relevant to ASD screening \\
\addlinespace
Frequency domain & Band powers, spectral ratios & Delta (0.5--4\,Hz), theta (4--8\,Hz), alpha (8--13\,Hz), beta (13--30\,Hz), gamma (30--45\,Hz); relative powers; $\theta$/$\beta$, $\alpha$/$\beta$ ratios & Oscillatory-state characterisation and ASD-relevant band-power differences \\
\addlinespace
Time--frequency & Wavelets, STFT & CWT scalograms (Morlet wavelet, 64 frequency bins), STFT spectrograms (Hann window, 256 samples) & Non-stationary spectral dynamics and short-lived ASD-relevant shifts in EEG structure \\
\addlinespace
Connectivity & Coherence, phase synchrony & Inter-channel coherence, phase-locking value (PLV), weighted phase-lag index (wPLI), mutual information & Inter-regional neural synchronisation and ASD-relevant connectivity differences \\
\addlinespace
Nonlinear & Entropy, fractal dynamics & Shannon entropy, sample entropy ($m$=2, $r$=0.2$\sigma$), Lyapunov exponent, Hurst exponent, correlation dimension & Signal complexity, regularity, and developmental EEG variability relevant to ASD screening \\
\end{xltabular}}
\vspace{0.5em}\noindent\textit{Note.} CWT = continuous wavelet transform; STFT = short-time Fourier transform; PLV = phase-locking value; wPLI = weighted phase-lag index; RMS = root mean square; PSD = power spectral density; ASD = autism spectrum disorder. Feature taxonomy follows the ASD EEG implementation documented in this thesis.

%%----------------------------------------------------------------------------
% S5: Wearable Sensor Integration (tab:wearable_sensors, fig:edge_deployment) — no wearable data used in experiments
%%----------------------------------------------------------------------------
% [SPACE-FIX] clearpage removed to eliminate empty page


\subsection{Archived Wearable Extension Note}
\label{sec:wearable_sensors}

\phantomsection\label{tab:wearable_sensors}
\phantomsection\label{fig:edge_deployment}
Earlier supplementary drafts included a wearable multi-sensor extension and edge-deployment pathway. That material is not part of the ASD-only thesis build because the evaluated artifact and evidence base in this version are limited to the ASD EEG workflow and its associated questionnaire, governance, and explainability layers.

\noindent\textit{Note.} The archived wearable extension is intentionally omitted here to keep Appendix C aligned with the ASD-only scope of the thesis.

%%----------------------------------------------------------------------------
% S6: Data Risk Assessment Matrix (fig:data_risk_matrix) — redundant with heatmap (fig:data_risk_heatmap) below
%%----------------------------------------------------------------------------
% [SPACE-FIX] clearpage removed to eliminate empty page


\subsection{Data Risk Assessment Matrix}
\label{sec:data_risk_matrix}

Figure~\ref{fig:data_risk_matrix} presents a risk assessment matrix for the nine data-related threats identified during the methodological design. Each risk is scored on two dimensions---likelihood (how probable the risk is given the study's design safeguards) and impact (the severity of consequence if the risk materialises)---producing a composite risk level that informs the prioritisation of mitigation efforts.
\needspace{24\baselineskip}
\begin{figure}[!htbp]
\centering
\resizebox{0.97\textwidth}{!}{%
\begin{tikzpicture}[
  rbox/.style={rectangle, draw=ggunavy, rounded corners=3pt,
    minimum height=1.1cm, text width=4.6cm, align=left, inner sep=5pt, font=\small},
  rlow/.style={rbox, fill=ggunavy!5},
  rmed/.style={rbox, fill=ggunavy!12},
  rhigh/.style={rbox, fill=ggunavy!22},
  rlbl/.style={rectangle, draw=ggunavy, fill=ggunavy, text=white, rounded corners=3pt,
    minimum height=1.0cm, text width=4.6cm, align=center, inner sep=4pt, font=\small\bfseries},
  node distance=0.8cm
]
% Column headers
\node[rlbl] (h1) at (0,0) {LOW RISK (2 risks)};
\node[rlbl] (h2) at (5.5,0) {MEDIUM RISK (4 risks)};
\node[rlbl] (h3) at (11,0) {HIGH RISK (3 risks)};

% Low risk column
\node[rlow, below=0.3cm of h1] (l1) {\textbf{R5} Outlier-driven accuracy inflation\\Likelihood: Low; Impact: Medium\\ICA rejection + z-score thresholds};
\node[rlow, below=0.6cm of l1] (l2) {\textbf{R8} Feature multicollinearity\\Likelihood: Medium; Impact: Low\\Ablation study isolates contributions};

% Medium risk column
\node[rmed, below=0.3cm of h2] (m1) {\textbf{R1} Data leakage (test contamination)\\Likelihood: Low; Impact: Critical\\Subject-level splits};
\node[rmed, below=0.6cm of m1] (m2) {\textbf{R3} Preprocessing param leakage\\Likelihood: Low; Impact: High\\Train-only fitting};
\node[rmed, below=0.6cm of m2] (m3) {\textbf{R6} Sampling rate heterogeneity\\Likelihood: Medium; Impact: Medium\\Standardisation to 256 Hz};
\node[rmed, below=0.6cm of m3] (m4) {\textbf{R7} Label noise (ground truth error)\\Likelihood: Low; Impact: Critical\\Peer-reviewed expert labels};

% High risk column
\node[rhigh, below=0.3cm of h3] (r1) {\textbf{R2} Class imbalance bias\\Likelihood: Medium; Impact: High\\SMOTE + stratified CV};
\node[rhigh, below=0.6cm of r1] (r2) {\textbf{R4} Dataset non-representativeness\\Likelihood: Medium; Impact: High\\Multi-site benchmarks; limitation declared};
\node[rhigh, below=0.6cm of r2] (r3) {\textbf{R9} Temporal autocorrelation in EEG\\Likelihood: Medium; Impact: High\\Subject-level (not epoch-level) splitting};
\end{tikzpicture}%
}
\caption{Data Risk Assessment Matrix: Nine Risks by Composite Level}
\label{fig:data_risk_matrix}
\vspace{0.5em}\noindent\textit{Note.} Composite level = likelihood $\times$ impact. Three risks (R2, R4, R9) reside in the high-risk zone; none reach critical. SMOTE = Synthetic Minority Over-sampling Technique; CV = cross-validation; ICA = independent component analysis.
\end{figure}

%%----------------------------------------------------------------------------
% S7: MCP Governance Capabilities table (tab:mcp_capabilities) and Flow diagram (fig:mcp_flow)
%%----------------------------------------------------------------------------
% MOVED TO CHAPTER 3 MAIN TEXT (Mar 2026)
\iffalse
\clearpage
\section{MCP Governance Capabilities table and Flow diagram}
\label{sec:appendix_ch3_s7}

\noindent Table~\ref{tab:app_mcp_capabilities} outlines model Control Protocol: Six Governance Capabilities.

\noindent Capabilities below indicate which agent owns each operation and how the result lands in the audit ledger.
{\tablestripes\tablefontsize
\needspace{20\baselineskip}
\begin{xltabular}{\textwidth}{@{} >{\raggedright\arraybackslash}p{2cm} p{4.7cm} p{4.0cm} >{\raggedright\arraybackslash}p{3.5cm} @{}}
\caption[Model Control Protocol]{Model Control Protocol: Six Governance Capabilities}\label{tab:app_mcp_capabilities} \\
\toprule
\thead{Capability} & \thead{Function} & \thead{Implementation} & \thead{Enforcement} \\
\midrule
\endfirsthead
\endhead
\bottomrule
\endlastfoot
Model registry & Central catalogue of all deployed models with version, hash, and metadata & MLflow model registry; SHA-256 integrity hash per model artefact & No model loads without registry entry \\
\addlinespace
Version control & Track model lineage: data version $\to$ feature version $\to$ model version $\to$ evaluation & Git-based versioning; DVC for data and model artefacts & Evaluation blocked if version chain incomplete \\
\addlinespace
Audit logs & Immutable record of all model predictions, inputs, outputs, and configuration & Append-only SQLite log; daily rotation; 90-day retention & Every inference call logged; no silent execution \\
\addlinespace
Access control & Role-based permissions: who can train, deploy, query, or rollback models & RBAC policy engine; three roles (admin, clinician, patient) & API key + role verification per request \\
\addlinespace
Deployment gates & Pre-deployment quality checks: accuracy $\geq$85\%, fairness variance $\leq$2\%, latency $<$200\,ms & Automated gate in CI/CD pipeline; all gates must pass & Failed gate blocks deployment; alerts to admin \\
\addlinespace
Rollback & Instant reversion to previous model version if post-future-validation metrics degrade & Blue-green deployment; previous version kept hot & Auto-rollback if accuracy drops $>$3\% in 24h \\
\end{xltabular}}
\vspace{0.5em}\noindent\textit{Note.} MCP = Model Control Protocol; RBAC = role-based access control; DVC = Data Version Control; CI/CD = continuous integration/continuous deployment. MCP operates at the orchestration layer, governing all seven agents (Figure~\ref{fig:agent_responsibilities}).

\needspace{24\baselineskip}
\begin{figure}[!htbp]
\centering
\begin{tikzpicture}[
  node distance=0.8cm and 0.15cm,
  block/.style={rectangle, draw=ggunavy, fill=ggunavy!8, rounded corners=3pt,
    minimum height=0.9cm, minimum width=2.0cm, align=center, font=\small},
  policyblock/.style={block, fill=ggunavy!20},
  arrow/.style={-{Stealth[length=5pt]}, thick, ggunavy},
  redarrow/.style={-{Stealth[length=5pt]}, thick, red!60!black, dashed}
]
\node[block, fill=lightblue] (req) {Agent\\Request};
\node[policyblock, fill=lightorange, right=of req] (policy) {MCP Policy\\Engine};
\node[policyblock, fill=lightteal, right=of policy] (auth) {Access\\Control};
\node[block, fill=lightpurple, right=of auth] (tool) {Tool / Model\\Execution};
\node[policyblock, fill=lightgreen, right=of tool] (audit) {Audit\\Logger};
\node[block, fill=lightyellow, right=of audit] (resp) {Response\\to Agent};
% Rejection path
\node[block, below=0.6cm of auth, fill=red!10, draw=red!60!black] (deny) {Deny +\\Alert Admin};
\draw[arrow] (req) -- (policy);
\draw[arrow] (policy) -- (auth);
\draw[arrow] (auth) -- (tool);
\draw[arrow] (tool) -- (audit);
\draw[arrow] (audit) -- (resp);
\draw[redarrow] (auth) -- (deny);
\node[below=0.08cm of policy, font=\small\itshape, text=gray] {6 capabilities};
\node[below=0.08cm of tool, font=\small\itshape, text=gray] {sandboxed};
\node[below=0.08cm of audit, font=\small\itshape, text=gray] {immutable log};
\end{tikzpicture}
\caption{Model Control Protocol (MCP) Governance Flow}
\label{fig:mcp_flow}
\vspace{0.5em}\noindent\textit{Note.} MCP = Model Control Protocol. Every agent request passes through the policy engine, access control, and audit logger before reaching the target tool or model. Failed access control checks trigger denial with admin notification. The audit logger creates an immutable record of all operations for regulatory compliance.
\end{figure}
\fi

%%----------------------------------------------------------------------------
% S8: NeuroMCP Multi-Agent Architecture diagram (fig:neuromcp_layers) and Tool Registry table (tab:mcp_tools)
%%----------------------------------------------------------------------------
% MOVED TO CHAPTER 3 MAIN TEXT (Mar 2026) — fig:neuromcp_layers only; tab:mcp_tools remains here
\iffalse
\clearpage


\subsubsection{NeuroMCP Multi-Agent Architecture}
\label{sec:neuromcp_architecture}

The NeuroMCP system~\cite{asthana2025eegmaster} extends the governance-level MCP (Table~\ref{tab:mcp_capabilities}) into a four-layer multi-agent architecture where disease-specific agents operate autonomously under centralised orchestration. Figure~\ref{fig:neuromcp_layers} illustrates the four architectural layers. This architecture addresses a key limitation of monolithic AI diagnostic systems: adding a new disease domain requires training a new agent rather than retraining the entire system.

\needspace{24\baselineskip}
\begin{figure}[!htbp]
\centering
\resizebox{0.97\textwidth}{!}{%
\begin{tikzpicture}[
  layer/.style={rectangle, draw=ggunavy, rounded corners=4pt,
    minimum width=13.5cm, minimum height=1.6cm, align=center, font=\small, inner sep=6pt},
  l1/.style={layer, fill=ggunavy!30},
  l2/.style={layer, fill=ggunavy!20},
  l3/.style={layer, fill=ggunavy!12},
  l4/.style={layer, fill=ggunavy!5},
  agentbox/.style={rectangle, draw=ggunavy, fill=white, rounded corners=2pt,
    minimum height=1.0cm, minimum width=2.8cm, align=center, font=\small, inner sep=3pt},
  arr/.style={-{Stealth[length=4mm, width=3mm]}, very thick, ggunavy},
  lbl/.style={font=\small\bfseries, text=ggunavy, anchor=east},
  node distance=0.8cm
]
% Layer 1: Model Control Portal
\node[l1, fill=lightblue] (layer1) {\textbf{\color{ggunavy}Layer 1: Model Control Portal} --- REST API gateway, authentication, rate limiting, request routing};

% Layer 2: MCP Orchestrator
\node[l2, fill=lightorange, below=of layer1] (layer2) {\textbf{\color{ggunavy}Layer 2: MCP Orchestrator} --- JSON-RPC 2.0, tool registry, agent lifecycle management, audit logging};

% Layer 3: Disease-Specific Agents
\node[l3, fill=lightteal, below=of layer2] (layer3) {};
\node[font=\small\bfseries, text=ggunavy] at ([yshift=0.3cm]layer3.center) {Layer 3: Disease-Specific Agents};
\node[agentbox, fill=lightgreen] (asd) at ([xshift=-4cm, yshift=-0.2cm]layer3.center) {ASD Agent};
\node[agentbox, fill=lightyellow] (pd) at ([yshift=-0.2cm]layer3.center) {PD Agent};
\node[agentbox, fill=lightpink] (stress) at ([xshift=4cm, yshift=-0.2cm]layer3.center) {Stress Agent};

% Layer 4: Data Processing Pipeline
\node[l4, fill=lightpurple, below=of layer3] (layer4) {\textbf{\color{ggunavy}Layer 4: Data Processing Pipeline} --- EEG preprocessing, feature extraction, model inference, RAG generation};

% Arrows between layers
\draw[arr] (layer1) -- (layer2);
\draw[arr] (layer2) -- (layer3);
\draw[arr] (layer3) -- (layer4);

% Layer labels
\node[lbl] at ([xshift=-0.3cm]layer1.west) {};
\end{tikzpicture}%
}
\caption{NeuroMCP Four-Layer Multi-Agent Architecture}
\label{fig:neuromcp_layers}
\vspace{0.5em}\noindent\textit{Note.} MCP = Model Control Protocol; REST = representational state transfer; JSON-RPC = JavaScript Object Notation Remote Procedure Call. Each disease-specific agent encapsulates its own preprocessing, model, and explanation pipeline. The orchestrator layer manages agent lifecycle, tool access, and audit logging via JSON-RPC 2.0. \textit{Source:}~\cite{asthana2025eegmaster}.
\end{figure}
\fi

Table~\ref{tab:mcp_tools} documents the 12 MCP tools available to agents, organised by functional category. Each tool is exposed via the JSON-RPC 2.0 interface and subject to the access control and audit logging described in Table~\ref{tab:mcp_capabilities}.

\paragraph{NeuroMCP Tool Registry.}


{\tablestripes\tablefontsize

\needspace{20\baselineskip}
\begin{xltabular}{\textwidth}{@{} >{\raggedright\arraybackslash}p{1.0cm} >{\raggedright\arraybackslash}p{1.0cm} p{8.0cm} >{\raggedright\arraybackslash}p{4.2cm} @{}}
\caption[NeuroMCP Tool Registry]{NeuroMCP Tool Registry: 12 Tools Organised by Functional Category}\label{tab:mcp_tools} \\
\toprule
\thead{Category} & \thead{Tool Name} & \thead{Function} & \thead{Access} \\
\midrule
\endfirsthead
\endhead
\bottomrule
\endlastfoot
Data ingest & {\texttt{load\_\allowbreak eeg}} & Load and validate raw EEG files (EDF/BDF/GDF) & All agents \\
\addlinespace
Pre-process & {\texttt{filter\_\allowbreak signal}} & Apply band-pass and notch filters & All EEG agents \\
\addlinespace
 & {\texttt{extract\_\allowbreak features}} & Compute 127-feature vector per channel & All EEG agents \\
\addlinespace
Inference & {\texttt{classify}} & Run VotingClassifier or ResNet-18 inference & All agents \\
\addlinespace
 & {\texttt{explain\_\allowbreak shap}} & Generate SHAP feature importance values & All agents \\
\addlinespace
RAG & {\texttt{retrieve\_\allowbreak context}} & Query FAISS vector store for clinical literature & All agents \\
\addlinespace
 & {\texttt{generate\_\allowbreak narrative}} & Produce evidence-grounded clinical explanation via SLM & All agents \\
\addlinespace
Governance & {\texttt{check\_\allowbreak guardrails}} & Validate input/output against safety rules (Sec.~\ref{sec:guardrail_rules}) & All agents \\
\addlinespace
 & {\texttt{log\_\allowbreak audit}} & Write immutable audit record of prediction and explanation & System (auto) \\
\addlinespace
 & {\texttt{redact\_\allowbreak pii}} & Detect and mask PII & System (auto) \\
\end{xltabular}}
\vspace{0.5em}\noindent\textit{Note.} MCP = Model Control Protocol; EDF = European Data Format; BDF = BioSemi Data Format; GDF = General Data Format; PBS = peripheral blood smear; GAN = generative adversarial network; SHAP = SHapley Additive exPlanations; FAISS = Facebook AI Similarity Search; SLM = small language model; PII = personally identifiable information. Tool registry follows~\cite{asthana2025eegmaster}.

%%----------------------------------------------------------------------------
% S9: Guardrail AI input/output safety rules tables (tab:ch3_input_guardrails, tab:ch3_output_guardrails)
%%----------------------------------------------------------------------------
% MOVED TO CHAPTER 3 MAIN TEXT (Mar 2026)
\iffalse
\clearpage


\subsection{Guardrail AI}
\label{sec:app_guardrail_rules}

The Guardrail AI flow (Figure~\ref{fig:guardrail_flow}) implements six safety gates. This subsection documents the specific input and output rules that govern clinical content generation. These rules argue that patient-facing AI systems require \textit{both} input and output guardrails: input guardrails prevent dangerous queries from reaching the model, while output guardrails prevent unsafe content from reaching the patient.

\paragraph{Input Guardrail Rules.}

Five input rules are applied before any query reaches the RAG pipeline. Detailed rule specifications are provided in the Supplementary Volume.

\begin{enumerate}[leftmargin=*, itemsep=3pt]
    \item \textbf{Block diagnosis intent:} Reject queries that request a definitive medical diagnosis.
    \item \textbf{Block prescription requests:} Reject queries that request medication or dosage recommendations.
    \item \textbf{Redirect high-risk queries:} Route queries mentioning self-harm, suicide, or emergency symptoms to crisis resources and human clinicians.
    \item \textbf{Refuse out-of-scope queries:} Decline queries outside the five clinical domains (ASD, PD, Cancer, Stress, BCI) with an explanatory message.
    \item \textbf{Redact PII:} Detect and mask personally identifiable information before the query enters processing.
\end{enumerate}

\needspace{24\baselineskip}
\noindent Table~\ref{tab:app_ch3_input_guardrails} presents input Guardrail Rules: Five Pre-Processing Safety Checks.

{\tablestripes\tablefontsize

\noindent Each guardrail is paired with the prompt class it rejects, the action taken, and the audit row written when it fires.
\begin{xltabular}{\textwidth}{@{} >{\raggedright\arraybackslash}p{2cm} L L >{\raggedright\arraybackslash}p{3.5cm} @{}}
\caption[Input Guardrail Rules]{Input Guardrail Rules: Five Pre-Processing Safety Checks}\label{tab:app_ch3_input_guardrails} \\
\toprule
\thead{Rule} & \thead{Trigger Condition} & \thead{Action} & \thead{Example} \\
\midrule
\endfirsthead
\endhead
\bottomrule
\endlastfoot
Block diagnosis intent & Query requests definitive diagnosis (``do I have cancer?'') & Block; redirect to clinician & ``Please consult your healthcare provider for diagnosis'' \\
\addlinespace
Block prescription & Query requests medication or dosage (``what drug should I take?'') & Block; display disclaimer & ``RGAIG does not provide prescription advice'' \\
\addlinespace
Redirect high-risk & Query mentions self-harm, suicide, or emergency symptoms & Redirect to crisis hotline + human clinician & Display crisis resources (988, emergency contacts) \\
\addlinespace
Refuse out-of-scope & Query outside the five clinical domains (ASD, PD, Cancer, Stress, BCI) & Refuse with explanation & ``This system is designed for neurodegenerative conditions'' \\
\addlinespace
Redact PII & Query contains patient name, DOB, SSN, or medical record number & Redact before processing; log detection event & ``John Smith'' $\to$ ``[PATIENT]'' \\
\end{xltabular}}
\vspace{0.5em}\noindent\textit{Note.} PII = personally identifiable information; DOB = date of birth; SSN = Social Security number. Rules are evaluated in priority order; the first matching rule terminates processing.

\paragraph{Output Guardrail Rules.}

Five output rules are applied after LLM generation and before display to the user. Detailed rule specifications are provided in the Supplementary Volume.

\begin{enumerate}[leftmargin=*, itemsep=3pt]
    \item \textbf{Downgrade medical authority:} Inject hedging language into responses that contain authoritative medical claims.
    \item \textbf{Inject uncertainty statements:} Add explicit uncertainty disclaimers when model confidence falls below 85\%.
    \item \textbf{Block unsupported claims:} Remove claims that lack grounding evidence from the retrieved context.
    \item \textbf{Replace hallucinated outputs:} Substitute outputs failing the hallucination gate (faithfulness $<$0.70) with a safe fallback response.
    \item \textbf{Escalate sensitive topics:} Flag outputs involving sensitive topics (end-of-life, paediatric, mental health) for human clinician review before patient display.
\end{enumerate}

\noindent Table~\ref{tab:app_ch3_output_guardrails} presents output Guardrail Rules: Five Post-Generation Safety Checks.

\noindent Each guardrail is paired with the response class it inspects, the action taken, and the audit row written when it fires.
{\tablestripes\tablefontsize
\needspace{20\baselineskip}
\begin{xltabular}{\textwidth}{@{} >{\raggedright\arraybackslash}p{2.2cm} L L >{\raggedright\arraybackslash}p{2.8cm} @{}}
\caption[Output Guardrail Rules]{Output Guardrail Rules: Five Post-Generation Safety Checks}\label{tab:app_ch3_output_guardrails} \\
\toprule
\thead{Rule} & \thead{Detection Method} & \thead{Action} & \thead{Threshold} \\
\midrule
\endfirsthead
\endhead
\bottomrule
\endlastfoot
Downgrade medical authority & Tone classifier detects authoritative medical claims & Inject hedging language (``this may indicate,'' ``consider consulting'') & Authority score $>$0.7 \\
\addlinespace
Inject uncertainty & Confidence below threshold for classification & Add explicit uncertainty statement to output & Model confidence $<$85\% \\
\addlinespace
Block unsupported claims & Claim-to-evidence mapping finds ungrounded assertions & Remove claim; replace with ``insufficient evidence'' & Zero matching passages \\
\addlinespace
Hallucination gate & Faithfulness score below threshold & Replace output with ``I don't have enough information'' & Faithfulness $<$0.70 \\
\addlinespace
Human review escalation & Sensitive topic detected (end-of-life, paediatric, mental health) & Flag for clinician review before patient display & Topic classifier confidence $>$0.8 \\
\end{xltabular}}
\vspace{0.5em}\noindent\textit{Note.} Output rules are evaluated sequentially; multiple rules may apply to the same output. Human review escalation is the final safety net before content reaches the patient.
\fi

%%----------------------------------------------------------------------------
% S10: PII Protection Pipeline (tab:pii_pipeline) and Security Alignment (tab:security_alignment) tables
%%----------------------------------------------------------------------------
% [SPACE-FIX] clearpage removed to eliminate empty page


\subsection{PII Protection and Security Architecture}
\label{sec:pii_security}

The PII Protection flow (Figure~\ref{fig:pii_flow}) implements a seven-stage privacy pipeline. This subsection documents the specific detection methods, encryption standards, and regulatory alignment that ensure patient data protection throughout the pipeline.


\noindent This table lists each \textit{\#} across process, technology, and verification, so examiners can trace what was assessed and what evidence supports each row.

{\tablestripes\tablefontsize
\needspace{20\baselineskip}
\begin{xltabular}{\textwidth}{@{} >{\raggedright\arraybackslash}p{1.2cm} L L >{\raggedright\arraybackslash}p{3.8cm} @{}}
\caption[PII Protection Pipeline]{PII Protection Pipeline: Six-Stage Implementation}\label{tab:pii_pipeline} \\
\toprule
\thead{\#} & \thead{Process} & \thead{Technology} & \thead{Verification} \\
\midrule
\endfirsthead
\endhead
\bottomrule
\endlastfoot
1 & Detect: identify PII entities in input text & Regex patterns (SSN, DOB, email, phone) + spaCy NER & Recall $>$99\% on synthetic PII test set \\
\addlinespace
2 & Classify: categorise entities by sensitivity & PHI (HIPAA-defined), PII (general), sensitive (clinical notes) & Three-tier classification \\
\addlinespace
3 & Redact: replace entities with type tokens & ``John Smith'' $\to$ ``[PERSON]''; ``123-45-6789'' $\to$ ``[SSN]'' & Redacted text reviewed by audit layer \\
\addlinespace
4 & Encrypt: protect data at rest and in transit & AES-256-GCM at rest; TLS 1.3 in transit & FIPS 140-2 compliant; 90-day key rotation \\
\addlinespace
5 & Access control: restrict by role and purpose & RBAC with three roles; GDPR Art.~5(1)(b) purpose limitation & Access denied without valid role + purpose \\
\addlinespace
6 & Audit: immutable access and processing log & Append-only SQLite log; tamper-evident hash chain & HIPAA \S164.312(b) compliance \\
\end{xltabular}}
\vspace{0.5em}\noindent\textit{Note.} PHI = protected health information; PII = personally identifiable information; NER = named entity recognition; RBAC = role-based access control; GCM = Galois/Counter Mode; FIPS = Federal Information Processing Standards.

% MOVED TO CHAPTER 3 MAIN TEXT (Mar 2026) — tab:security_alignment only
\iffalse
\paragraph{Regulatory Framework Alignment.}

RGAIG security controls are mapped to five regulatory frameworks---HIPAA \S164.312, GDPR Articles~25/32, NIST AI RMF, ISO~27001, and EU AI Act Article~14---demonstrating compliance readiness across encryption, access control, audit, pseudonymisation, and human oversight requirements. The detailed alignment matrix is provided in the Supplementary Volume.

\noindent Table~\ref{tab:app_security_alignment} outlines security Framework Alignment: Five Regulatory Standards.

{\tablestripes\tablefontsize
\needspace{20\baselineskip}
\begin{xltabular}{\textwidth}{@{} >{\raggedright\arraybackslash}p{2.5cm} L L @{}}
\caption[Security Framework Alignment]{Security Framework Alignment: Five Regulatory Standards}\label{tab:app_security_alignment} \\
\toprule
\thead{Regulation} & \thead{Key Requirement} & \thead{RGAIG Control} \\
\midrule
\endfirsthead
\endhead
\bottomrule
\endlastfoot
HIPAA \S164.312 & Technical safeguards: access control, audit, encryption, integrity & RBAC, AES-256, audit trail, hash verification \\
\addlinespace
GDPR Art.~25/32 & Data protection by design; encryption; pseudonymisation & Privacy-by-design architecture; $k$-anonymity $\geq$5 \\
\addlinespace
NIST AI RMF & Risk management throughout AI lifecycle; trustworthy AI & 15 governance layers (A1--A15); continuous monitoring \\
\addlinespace
ISO 27001 & Information security management system; risk assessment & Documented risk register; annual review cycle \\
\addlinespace
EU AI Act Art.~14 & Human oversight for high-risk AI systems & Human-in-the-loop escalation; confidence thresholds \\
\end{xltabular}}
\vspace{0.5em}\noindent\textit{Note.} HIPAA = Health Insurance Portability and Accountability Act; GDPR = General Data Protection Regulation; NIST AI RMF = National Institute of Standards and Technology AI Risk Management Framework; EU AI Act = European Union Artificial Intelligence Act.
\fi

%%----------------------------------------------------------------------------
% S11: 13-Phase GenAI QA Framework overview table (tab:13phase_overview) — audit checklist detail
%%----------------------------------------------------------------------------
% [SPACE-FIX] clearpage removed to eliminate empty page

\section{13-Phase ASD GenAI QA Framework overview table}
\label{sec:appendix_ch3_s11}

Table~\ref{tab:13phase_overview} presents the thirteen-phase GenAI quality-assurance framework used to assess the ASD build, detailing each phase's focus area, module count, key metrics, and sign-off gates.

\tablefontsize

\needspace{20\baselineskip}
\begin{xltabular}
{\textwidth}{@{} >{\raggedright\arraybackslash}p{0.8cm} >{\raggedright\arraybackslash}p{2.2cm} >{\raggedright\arraybackslash}p{0.6cm} L >{\raggedright\arraybackslash}p{3.2cm} @{}}
\caption[Thirteen-Phase ASD GenAI Quality Assurance Framework]{Thirteen-Phase ASD GenAI Quality Assurance Framework}\label{tab:13phase_overview} \\
\toprule
\thead{Phase} & \thead{Focus Area} & \thead{Mod.} & \thead{Key Metrics} & \thead{Sign-Off Gate} \\
\midrule
\endfirsthead
\endhead
\bottomrule
\endfoot
1 & Knowledge \& Data & 16 & Source inventory, authority scoring, coverage gaps, PHI/PII scan, ingestion completeness & All sources registered; 0 unhandled PHI \\
\addlinespace
2 & Representation \& Retrieval & 17 & Chunk integrity, embedding quality, Recall@K, Precision@K, MMR, Table-RAG routing & Recall@K $\geq$ baseline; no index corruption \\
\addlinespace
3 & Generation \& Reasoning & 17 & Citation coverage, claim verification, hallucination rate, numeric integrity, answer relevancy & Unsupported claims $\leq$3\%; faithfulness $\geq$0.85 \\
\addlinespace
4 & Decision Policy & 17 & Intent classification, risk-tier assignment, confidence calibration, escalation accuracy & False answers $\leq$ limit; 0 unsafe passes \\
\addlinespace
5 & Agent Behaviour & 17 & Task completion, tool selection, loop detection, error recovery, autonomy bounds & Goal completion $\geq$ target; 0 unsafe actions \\
\addlinespace
6 & Agent-to-Agent & 17 & Role adherence, message protocol, conflict resolution, coordination, failure isolation & $\geq$95\% role adherence; 0 cascade failures \\
\addlinespace
7 & Model Control Protocol & 17 & Policy coverage, tool access enforcement, bypass detection, prompt override resistance & 0 policy bypasses; 0 unauthorised tool calls \\
\addlinespace
8 & Explainability \& Trust & 17 & Evidence visibility, explanation consistency, confidence communication, comprehension & $\geq$97\% answers with visible sources \\
\addlinespace
9 & Robustness \& Sensitivity & 17 & Prompt perturbation, context removal, adversarial inputs, sparsity handling & Graceful degradation; 0 adversarial successes \\
\addlinespace
10 & Statistical Validation & 17 & Hypothesis testing, sample adequacy, CI estimation, multiple correction, reproducibility & Power $\geq$0.8; $p < \alpha$; low variance \\
\addlinespace
11 & Benchmarking & 17 & RAG vs non-RAG, hybrid vs dense retrieval, chunking ablation, SOTA comparison, cost & RAG $\geq$ baseline; cost within budget \\
\addlinespace
12 & Scalability \& Deploy & 17 & Token budget, retrieval/generation latency, throughput, cache hit rate, rollback readiness & SLA met; 100\% reproducible builds \\
\addlinespace
13 & Governance \& Compliance & 17 & Lineage completeness, PII exposure, consent validation, incident response, regulatory mapping & $\geq$99\% lineage; 0 PII exposure \\
\end{xltabular}
\vspace{0.5em}\noindent\textit{Note.} Each phase contains 16--17 monitoring modules, each with defined purpose, measurement method, outputs, pass/fail criteria, and edge-case fixes. In the ASD build, the 13-phase framework extends classical validation (Phases 10--11) with six GenAI-specific phases (1--3: knowledge and generation quality), four agentic phases (4--7: agent and policy governance), and two operational phases (8--9: trust and robustness). All 13 gates must pass before evaluation in the bounded ASD screening context. PHI = protected health information; PII = personally identifiable information; MMR = maximal marginal relevance; MCP = Model Control Protocol; A2A = agent-to-agent; CI = confidence interval; SLA = service-level agreement; SOTA = state of the art.

%%----------------------------------------------------------------------------
% S12: 8-Phase DL Analysis Framework overview table (tab:8phase_overview) — audit checklist detail
%%----------------------------------------------------------------------------
% [SPACE-FIX] clearpage removed to eliminate empty page

\section{8-Phase DL Analysis Framework overview table}
\label{sec:appendix_ch3_s12}

Table~\ref{tab:8phase_overview} summarises the eight-phase deep learning analysis framework comprising 111 modules across data, model, performance, subject, sensitivity, statistical, benchmarking, and conceptual monitoring reference (not deployed) phases.

\tablefontsize

\needspace{20\baselineskip}
\begin{xltabular}
{\textwidth}{@{} >{\raggedright\arraybackslash}p{0.5cm} >{\raggedright\arraybackslash}p{2.2cm} >{\raggedright\arraybackslash}p{0.6cm} X >{\raggedright\arraybackslash}p{3.2cm} @{}}
\caption[Eight-Phase DL Analysis Framework Overview (111 Modules)]{Eight-Phase DL Analysis Framework Overview (111 Modules)}\label{tab:8phase_overview} \\
\toprule
\thead{\#} & \thead{Phase} & \thead{Mod.} & \thead{Focus and Key Analyses} & \thead{RGAIG Evidence} \\
\midrule
\endfirsthead
\endhead
\midrule
\multicolumn{5}{r}{\normalsize\itshape Continued on next page} \\
\endfoot
\bottomrule
\endlastfoot
1 & Data Analysis & 9 & Data inventory, class/label distribution, missingness, SQI, cross-subject variability, session drift, harmonisation, leakage audit, documentation & 305\,GB; 131 subjects + 20{,}000 images; SQI per channel; grouped split \\
\addlinespace
2 & Model Analysis & 14 & Model--data alignment, baselines (SVM/RF/LDA), complexity vs generalisation, training stability, regularisation, class-imbalance, calibration, interpretability, computational cost, failure-mode anticipation, readiness gate & 7 baseline + 4 DL architectures; learning curves; SHAP; cost profiling \\
\addlinespace
3 & Performance Analysis & 14 & Metric correctness, fold-level performance, seed stability, CI, operating-point analysis, calibration, robustness, subgroup performance, error breakdown, baseline deltas, SOTA benchmark, sign-off gate & 10-fold CV; bootstrap CI; ablation $\Delta$5.3--8.1\%; SOTA comparison \\
\addlinespace
4 & Subject Analysis & 14 & Subject coverage, per-subject F1/AUC, worst-case, session drift, similarity clusters, device confounding, artifact sensitivity, label reliability, personalisation vs generalisation, demographic fairness, sign-off gate & Per-subject metrics; age/gender stratification; worst-case floor \\
\addlinespace
5 & Sensitivity Analysis & 16 & Filter cutoffs, notch settings, referencing (CAR/Cz), artifact strictness, ICA/ASR, window length (1--8\,s), overlap, normalisation, class-imbalance strategy, HP neighbourhood, seed sensitivity, cross-dataset, sign-off & 3-scenario + tornado; HP stability; 10-seed $\pm$0.8\% \\
\addlinespace
6 & Statistical Analysis & 14 & Hypothesis definition, effect sizes (Cohen's $d$/Hedges' $g$), CI, paired comparison (McNemar/DeLong), Holm/BH-FDR correction, Wilcoxon, subgroup disparity, model ranking, power adequacy, reporting pack & $d$=0.82--1.48; bootstrap CI; Holm correction; subject-level bootstrap \\
\addlinespace
7 & Benchmarking & 14 & Results registry, dataset/split table, preprocessing config, main results with CI, baseline delta, significance, robustness, sensitivity matrix, error taxonomy, explainability, literature matrix, claims-to-evidence map & Per-ASD pack; SOTA matrix; claims traceability \\
\addlinespace
8 & Production Monitoring & 16 & Data health, distribution drift (PSD), representation drift, prediction drift, confidence monitoring, latency, silent failure detection, drift-to-action policy, retraining readiness, A/B shadow testing, governance logging, illustrative financial scenario monitoring & Drift monitor; batch-wise accuracy; retraining triggers; model cards \\
\end{xltabular}
\vspace{0.5em}\noindent\textit{Note.} The eight-phase framework draws on implementations totalling 24,287 lines of analysis code across \texttt{data\_\allowbreak analysis.py} (1,968 lines), \texttt{statistical\_\allowbreak analysis.py} (1,405 lines), \texttt{cross\_\allowbreak validation.py} (692 lines), \texttt{eda\_\allowbreak analysis.py} (995 lines), \texttt{reliability\_\allowbreak analysis.py} (995 lines), \texttt{filter\_\allowbreak analysis.py} (727 lines), \texttt{drift\_\allowbreak monitor.py} (664 lines), and \texttt{generate\_\allowbreak ml\_\allowbreak dl\_\allowbreak analysis.py} (1,361 lines). SQI = signal quality index; PSD = power spectral density; HP = hyperparameter; BH-FDR = Benjamini--Hochberg false discovery rate.

%%----------------------------------------------------------------------------
% S13: 10 Pillars of Cross-Cutting AI Governance overview table (tab:10pillars) — audit checklist detail
%%----------------------------------------------------------------------------
% [SPACE-FIX] clearpage removed to eliminate empty page

\section{10 Pillars of Cross-Cutting AI Governance overview table}
\label{sec:appendix_ch3_s13}

\tablefontsize

\needspace{20\baselineskip}

\noindent This table lists each \textit{\#} across pillar, mod., core question and key analyses, and 1 more dimensions, so examiners can trace what was assessed and what evidence supports each row.

\begin{xltabular}
{\textwidth}{@{} >{\raggedright\arraybackslash}p{0.5cm} >{\raggedright\arraybackslash}p{2cm} >{\raggedright\arraybackslash}p{0.6cm} X >{\raggedright\arraybackslash}p{3.2cm} @{}}
\caption[Ten Pillars of Cross-Cutting AI Governance (194 Modules)]{Ten Pillars of Cross-Cutting AI Governance (194 Modules)}\label{tab:10pillars} \\
\toprule
\thead{\#} & \thead{Pillar} & \thead{Mod.} & \thead{Core Question and Key Analyses} & \thead{RGAIG Implementation} \\
\midrule
\endfirsthead
\endhead
\midrule
\multicolumn{5}{r}{\normalsize\itshape Continued on next page} \\
\endfoot
\bottomrule
\endlastfoot
1 & Energy Efficient AI & 18 & Is energy justified? Workload characterisation, architecture efficiency, right-sizing, PEFT/LoRA, quantisation (FP16$\to$INT8), batching/caching, latency--energy trade-off, scaling, ESG alignment & Local SLM (3B params); INT8 quantisation; L1/L2 cache; 512-token limit \\
\addlinespace
2 & Hallucination Prevention & 20 & How are fabricated outputs prevented? Taxonomy, knowledge boundaries, prompt robustness, RAG grounding, source attribution, faithfulness, reasoning chain, uncertainty/abstention, adversarial tests, HITL, monitoring & Faithfulness 0.89; hallucination 2.1\%; citation 94.2\%; 6-gate guardrail \\
\addlinespace
3 & Hypothesis in AI & 20 & Are assumptions stated and tested? Problem framing, data sufficiency, label validity, feature relevance, leakage, model capacity, convergence, generalisation, imbalance, error patterns, robustness, explainability, causal mechanisms, drift, safety & H1--H5 pre-registered; 5 falsifiable criteria; power $\geq$0.80 \\
\addlinespace
4 & Threat AI & 20 & What adversarial threats exist? Asset identification, attack surface, data poisoning, prompt injection, model extraction, membership inference, adversarial examples, RAG poisoning, DoS, supply chain, incident response, residual risk & PII scrubbing; RBAC; prompt override testing; red-team validation \\
\addlinespace
5 & SWOT for AI & 16 & Strategic assessment: capability, data advantage, efficiency (S); data dependency, XAI gaps, generalisation limits (W); new products, scale, differentiation (O); regulatory pressure, ethical risk, security (T) & SWOT by lifecycle stage (data, model, deploy, monitor, govern) \\
\addlinespace
6 & Governance AI & 20 & Who owns AI decisions? Scope/authority, RACI, policy alignment, risk management, ethics review, data governance, model lifecycle, monitoring, incident escalation, HITL, vendor governance, documentation, enforcement & 4-level escalation ladder; RACI matrix; MCP policy engine \\
\addlinespace
7 & Compliance AI & 20 & Which laws apply? Jurisdiction mapping (5), risk classification, GDPR/HIPAA, fairness, safety controls, human oversight, right-to-explanation, post-market surveillance, vendor compliance, audit readiness, accountability & 35+ standards; 5-jurisdiction compliance; audit-ready docs \\
\addlinespace
8 & Responsible AI & 20 & Are outcomes equitable and harm-minimised? Stakeholder impact, misuse analysis, fairness, transparency, oversight, accountability, privacy-by-design, safety, monitoring, contestability, environmental/social impact & 9 responsibility dimensions; Fairlearn bias audit; SHAP \\
\addlinespace
9 & Explainable AI & 20 & Can decisions be understood? SHAP/LIME local, global feature importance, PDP/ICE, counterfactuals, faithfulness, stability, bias audit, temporal XAI, Grad-CAM, token attribution, interpretability, reproducibility & SHAP top-5; 3-tier RAG explanation; expert $M$=4.6/5.0; $\alpha$=0.84 \\
\addlinespace
10 & Secure AI & 20 & Are AI assets protected? Asset inventory, threat model, data integrity, training security, model IP, adversarial robustness, prompt injection defence, RBAC, privacy/leakage, supply chain, logging/audit, incident response, penetration testing & AES-256; RBAC; PHI/PII detection; immutable audit logs \\
\end{xltabular}
\vspace{0.5em}\noindent\textit{Note.} The ten pillars address cross-cutting concerns that span all pipeline stages. Module counts: Energy-Efficient (18), Hallucination Prevention (20), Hypothesis (20), Threat (20), SWOT (16), Governance (20), Compliance (20), Responsible AI (20), Explainable AI (20), Secure AI (20). HITL = human-in-the-loop; PEFT = parameter-efficient fine-tuning; LoRA = low-rank adaptation; SWOT = strengths, weaknesses, opportunities, threats; DoS = denial of service; ESG = environmental, social, and governance.

%%----------------------------------------------------------------------------
% S14: Responsible AI Governance Summary table (tab:responsible_ai_summary) — audit-level specification
%%----------------------------------------------------------------------------
% [SPACE-FIX] clearpage removed to eliminate empty page

\section{Responsible AI Governance Summary table}
\label{sec:appendix_ch3_s14}

\tablefontsize

\noindent Each principle is paired with its operational mechanism, primary owner, and verification cadence.
\needspace{20\baselineskip}

\noindent This table lists each \textit{governance pillar} across requirement, implementation in this study, and evidence / verification, so examiners can trace what was assessed and what evidence supports each row.

\begin{xltabular}
{\textwidth}{@{} >{\raggedright\arraybackslash}p{2cm} >{\raggedright\arraybackslash}p{2.5cm} L L @{}}
\caption[Executive Summary]{Executive Summary: Responsible AI Governance Framework}\label{tab:responsible_ai_summary} \\
\toprule
\thead{Governance Pillar} & \thead{Requirement} & \thead{Implementation in This Study} & \thead{Evidence / Verification} \\
\midrule
\endfirsthead
\endhead
\bottomrule
\endfoot
\multicolumn{4}{l}{Data Compliance} \\
\addlinespace
& HIPAA de-identification & All datasets pre-anonymized by original collectors & Data use certificates (Appendix~\ref{app:data_certificates}) \\
& GDPR data minimization & Only task-relevant features extracted; no surplus data retained & Preprocessing pipeline documentation (Section~\ref{sec:data_prep}) \\
& IRB exemption & Publicly available de-identified data; 45 CFR 46 exemption & Institutional compliance letter \\
& Encrypted storage & Raw data on AES-256 encrypted local workstation; no cloud & Single-user access; 5-year retention policy \\
\addlinespace
\multicolumn{4}{l}{Algorithmic Fairness} \\
\addlinespace
& Bias detection & Demographic subgroup accuracy analysis & $\leq$2\% variance across available subgroups (Ch.~4) \\
& Transparency & SHAP feature attribution + RAG literature-grounded narratives & Expert panel validation (\textit{M} = 4.6/5) \\
& Human oversight & AI recommends, clinician decides; no autonomous decisions & Human-in-the-loop design; escalation protocols \\
\addlinespace
\multicolumn{4}{l}{Data Provenance} \\
\addlinespace
& Source verification & NIH/NITRC, PhysioNet, university repositories & KAU ASD dataset with documented provenance (Table~\ref{tab:dataset_inventory_full}) \\
& Access protocol & Data use agreements signed per dataset & DUA records maintained; access logging enabled \\
& Audit trail & All preprocessing steps logged with parameters & Reproducible via Docker + fixed seeds + Git \\
\addlinespace
\multicolumn{4}{l}{Reproducibility} \\
\addlinespace
& Environment & Docker containers with pinned dependency versions & Dockerfile versioned in repository \\
& Determinism & Fixed random seed (42) across all frameworks & PyTorch, TensorFlow, NumPy, Scikit-learn \\
& Experiment tracking & MLflow for hyperparameters, metrics, artifacts & Tagged experiment commits in Git \\
\addlinespace
\multicolumn{4}{l}{Regulatory Alignment (Future Deployment)} \\
\addlinespace
& FDA SaMD & Framework designed for IEC 62304 software lifecycle & Governance architecture supports SaMD pathway \\
& EU AI Act & High-risk classification; human oversight mandated & Audit trails, bias docs, decision logic embedded \\
& WHO/ITU GMLP & Good Machine Learning Practice principles & Validation, monitoring, documentation controls \\
\end{xltabular}

\vspace{0.5em}\noindent\textit{Note.} HIPAA = Health Insurance Portability and Accountability Act; GDPR = General Data Protection Regulation; IRB = Institutional Review Board; SHAP = SHapley Additive exPlanations; RAG = retrieval-augmented generation; FDA = Food and Drug Administration; SaMD = Software as a Medical Device; EU = European Union; WHO = World Health Organization; ITU = International Telecommunication Union; GMLP = Good Machine Learning Practice; DUA = data use agreement. Regulatory alignment items represent design intent for future future clinical-validation pathway, not current certification status.


%%============================================================================
%% S15: Engineering Pipeline Flowcharts
%%============================================================================
% [SPACE-FIX] clearpage removed to eliminate empty page

%%%%%%%%%%%%%%%%%%%%%%%%%%%%%%%%%%%%%%%%%%%%%%%%%%%%%%%%%%%%%%%%%%%%%%%%%%%%%%%%
%% EEG PROCESSING DEFENSIBILITY PACK
%% Reframes the technical EEG pipeline as governance-embedded supporting
%% context for the DBA contribution. Pre-empts the supervisor critique
%% that the thesis could revert into a biomedical-engineering frame.
%%%%%%%%%%%%%%%%%%%%%%%%%%%%%%%%%%%%%%%%%%%%%%%%%%%%%%%%%%%%%%%%%%%%%%%%%%%%%%%%

\addcontentsline{toc}{section}{EEG Processing Defensibility Pack} % [TOC-appendix-section-insertion]
\section*{EEG Processing Defensibility Pack}
\addcontentsline{toc}{section}{EEG Processing Defensibility Pack}
\label{sec:eeg_defensibility_pack}

\noindent\textbf{Critical disclaimer (read first).} The EEG processing workflow documented in this section represents a \emph{governance-enabled AI-assisted decision-support environment} and is \textbf{not} intended as an autonomous clinical diagnostic system. The pipeline supports clinician decision-making with explainable, governance-monitored predictions; the clinician remains the decision-maker. All technical details below exist to operationalize the dissertation's Governance $\to$ Explainability $\to$ Trust $\to$ Adoption pathway, not to claim biomedical-engineering novelty.

\noindent\textbf{Reading order.} This pack is organised in two layers, matching the Architecture Appendix convention:
\begin{itemize}[leftmargin=*, itemsep=2pt]
    \item \textbf{Layer 1 (this pack):} Governance-embedded EEG workflow + human-centered pipeline + research-to-technical mapping. \emph{Doctoral-relevant.}
    \item \textbf{Layer 2 (subsequent engineering flowcharts):} Five-phase processing detail + three end-to-end pipeline diagrams. \emph{Engineering-relevant only; provided for completeness, not as dissertation contribution.}
\end{itemize}

\noindent Table~\ref{tab:eeg_why_domain} gives the short decision logic for selecting EEG-based ASD screening as the empirical domain.

%% =========================================================================
%% TABLE 1: WHY EEG/ASD AS THE EMPIRICAL DOMAIN
%% =========================================================================
\needspace{24\baselineskip}
\noindent Table~\ref{tab:eeg_why_domain} explains why each clinical domain in the empirical evaluation was selected and what EEG-specific evidence base supports its inclusion.

\paragraph{Why EEG/ASD as the.}




{\tablestripes\tablefontsize
\needspace{20\baselineskip}
\begin{xltabular}{\textwidth}{@{} >{\raggedright\arraybackslash}p{3.0cm} L @{}}
\caption[Why EEG/ASD as the Empirical Domain]{Why EEG-Based ASD Screening Is the Empirical Domain}\label{tab:eeg_why_domain} \\
\toprule
\thead{Reason} & \thead{Organisational value} \\
\midrule
\endfirsthead
\endhead
\bottomrule
\endlastfoot
\textbf{Early intervention}      & Faster screening can reduce care delay and downstream cost. \\
\textbf{Subjective diagnosis}    & Multi-source assessment creates a clear need for calibrated AI support. \\
\textbf{Clinician uncertainty}   & Explainable second opinion has practical value when cases are borderline. \\
\textbf{Decision workflow}       & Referral, EEG, behaviour, and diagnosis fit a staged support model. \\
\textbf{High-risk context}       & Paediatric screening makes explainability and governance mandatory. \\
\textbf{Ethical sensitivity}     & Consent, minimisation, and equity controls are materially testable. \\
\end{xltabular}}
\vspace{0.4em}
\noindent\textit{Note.} EEG-based ASD screening is not chosen for technical convenience but for being a domain where every variable in the conceptual framework (governance, explainability, trust, adoption) has material clinical significance.

%% =========================================================================
%% FIGURE 1: HUMAN-CENTERED EEG WORKFLOW (governance EMBEDDED, not separate)
%% =========================================================================
\needspace{24\baselineskip}
\begin{figure}[!htbp]
\centering
\begin{tikzpicture}[every node/.style={font=\fignodefont},
    eegbx/.style={draw=ggunavy, fill=ggunavy!10, rounded corners=4pt, thick, minimum width=4.0cm, minimum height=0.9cm, align=center, font=\fignodefont},
    govbx/.style={draw=ggunavy, fill=orange!12, rounded corners=4pt, thick, minimum width=4.0cm, minimum height=0.9cm, align=center, font=\fignodefont\itshape},
    humanbx/.style={draw=ggunavy, fill=ggunavy!18, rounded corners=4pt, very thick, minimum width=4.0cm, minimum height=0.9cm, align=center, font=\fignodefont\bfseries},
    arr/.style={-{Stealth[length=3mm, width=2mm]}, thick, ggunavy}]
\node[eegbx] (s1)  at (0,0)     {\textbf{1.\ EEG Data}\\\scriptsize De-identified, lineage-tagged};
\node[govbx] (g1)  at (0,-1.5)  {\textbf{Gov Step:}\ Lineage + Access ctrl};
\node[eegbx] (s2)  at (0,-3.0)  {\textbf{2.\ AI-Assisted Analysis}\\\scriptsize ASD classifier (decision support, not diagnosis)};
\node[govbx] (g2)  at (0,-4.5)  {\textbf{Gov Step:}\ Bias monitoring + drift check};
\node[eegbx] (s3)  at (0,-6.0)  {\textbf{3.\ Explainability Layer}\\\scriptsize SHAP attribution + RAG citation grounding};
\node[govbx] (g3)  at (0,-7.5)  {\textbf{Gov Step:}\ Audit log + traceability};
\node[humanbx](s4) at (0,-9.0)  {\textbf{4.\ Clinician Review (HITL)}\\\scriptsize Inspect explanation; override available};
\node[govbx] (g4)  at (0,-10.5) {\textbf{Gov Step:}\ Override rationale capture};
\node[eegbx] (s5)  at (0,-12.0) {\textbf{5.\ Decision Support Output}\\\scriptsize EHR-integrated; clinician decision recorded};
\foreach \a/\b in {s1/g1, g1/s2, s2/g2, g2/s3, s3/g3, g3/s4, s4/g4, g4/s5} { \draw[arr] (\a) -- (\b); }
\draw[arr, dashed, bend left=70] (s5.east) to node[right, font=\scriptsize] {feedback to governance} (s2.east);
\end{tikzpicture}
\caption[Human-Centered EEG Workflow]{Human-Centered EEG Workflow with Embedded Governance. The clinician remains the decision-maker at Stage~4 (HITL); governance steps (orange italic boxes) are not separate from the workflow but are embedded between every EEG processing stage. This figure replaces the engineering-centric pipeline view in this appendix's deeper sections; the technical phases appear here only as supporting context.}
\label{fig:eeg_human_workflow}
\end{figure}

%% =========================================================================
%% TABLE 2: OPERATIONAL WORKFLOW → GOVERNANCE ROLE MAPPING
%% =========================================================================
\needspace{24\baselineskip}
\noindent Table~\ref{tab:eeg_workflow_governance} summarises operational eeg workflow mapped to governance role for appendix review.

\paragraph{Operational Workflow vs. Governance.}




{\tablestripes\tablefontsize
\needspace{20\baselineskip}
\begin{xltabular}{\textwidth}{@{} >{\raggedright\arraybackslash}p{3.9cm} p{8.9cm} >{\raggedright\arraybackslash}p{1.5cm} @{}}
\caption[Operational Workflow vs.\ Governance Role]{Operational EEG Workflow Mapped to Governance Role. Each technical workflow step is paired with the governance control family it operationalizes, ensuring no step exists without governance accountability.}\label{tab:eeg_workflow_governance} \\
\toprule
\thead{Workflow step} & \thead{Governance role} & \thead{Supports H\#} \\
\midrule
\endfirsthead
\endhead
\bottomrule
\endlastfoot
EEG ingestion          & Data lineage (source, time, device, patient ID-hash); access control on raw data; immutable arrival log & H1 \\
Preprocessing          & Audit log of every preprocessing transformation (filter parameters, channel rejection, artefact removal); reproducibility record & H1, H2 \\
Feature extraction     & Lineage from raw signal to feature vector; feature catalog with derivation source                                                  & H1 \\
AI classification      & Model version recorded; SHAP attribution generated alongside prediction; subgroup fairness logged                              & H1, H2 \\
Explainability render  & SHAP + RAG citation grounding rendered in clinical UI; explanation traceability recorded                                       & H2 \\
Clinician review (HITL) & Clinician inspection logged; override capability; rationale capture if override exercised                                     & H3, H4 \\
Output validation      & Final clinician decision recorded; clinician accountability; outcome feedback to governance monitoring                          & H3, H5 \\
\end{xltabular}}
\vspace{0.4em}
\noindent\textit{Note.} The technical workflow is not separable from governance; governance is embedded at every step. This is the operational instantiation of the H1--H5 pathway tested in Chapter~4.

%% =========================================================================
%% TABLE 3: WHY EXISTING EEG AI SYSTEMS FAIL ORGANIZATIONALLY
%% =========================================================================
\needspace{24\baselineskip}
\noindent Table~\ref{tab:eeg_existing_failures} summarises why existing eeg ai systems fail to achieve clinical adoption for appendix review.

\paragraph{Why Existing EEG AI.}




{\tablestripes\tablefontsize
\needspace{20\baselineskip}
\begin{xltabular}{\textwidth}{@{} >{\raggedright\arraybackslash}p{3.8cm} p{6.8cm} >{\raggedright\arraybackslash}p{3.8cm} @{}}
\caption[Why Existing EEG AI Systems Fail Organizationally]{Why Existing EEG AI Systems Fail to Achieve Clinical Adoption. The technical gap is not accuracy (already achieved at $>$95\%); it is the governance/explainability/trust pathway this dissertation addresses.}\label{tab:eeg_existing_failures} \\
\toprule
\thead{Existing limitation} & \thead{Organizational impact} & \thead{Addressed by} \\
\midrule
\endfirsthead
\endhead
\bottomrule
\endlastfoot
Poor explainability             & Clinician distrust; refusal to override own judgement in favour of AI; adoption stalls at pilot                  & H1, H2 (SHAP+RAG layer) \\
Black-box models                & Adoption resistance; legal exposure under EU AI Act high-risk classification                                  & H1 (governance maturity) \\
Weak governance                 & Accountability concerns; CIO / compliance hesitation to conceptually evaluate                                                  & H1 (RAI controls) \\
Lack of clinician oversight     & Operational hesitation; over-reliance risk (per Chapter~4 negative findings)                                    & H3 (HITL workflow) \\
No workflow integration         & Even when accurate, AI cannot be used at point of care                                                          & H5 (workflow readiness) \\
No fairness monitoring          & Bias-amplification risk; subgroup harm; reputational exposure                                                   & H1 (bias control) \\
No rollback path                & Single-failure mode disables AI across the institution                                                          & H1 (rollback registry) \\
\end{xltabular}}
\vspace{0.4em}
\noindent\textit{Note.} Each limitation is mapped to the dissertation hypothesis that addresses it. The dissertation's claim is not that existing EEG AI is technically inadequate but that it is \emph{governance-inadequate} for clinical adoption.

%% =========================================================================
%% TABLE 4: RESEARCH-CONSTRUCT → TECHNICAL-SUPPORT MAPPING
%% =========================================================================
\needspace{24\baselineskip}
\noindent Table~\ref{tab:eeg_research_technical_map} summarises research construct mapped to technical support for appendix review.

\paragraph{Research Construct vs. Technical.}




{\tablestripes\tablefontsize
\needspace{20\baselineskip}
\begin{xltabular}{\textwidth}{@{} >{\raggedright\arraybackslash}p{3.0cm} L @{}}
\caption[Research Construct vs.\ Technical Support]{Research Construct Mapped to Technical Support. Every technical component in this appendix is justified by its role in supporting one research construct; engineering details outside this mapping are demoted to deep-appendix material.}\label{tab:eeg_research_technical_map} \\
\toprule
\thead{Research construct} & \thead{Technical support layer} \\
\midrule
\endfirsthead
\endhead
\bottomrule
\endlastfoot
\textbf{RAI Governance (IV)}            & Lineage tracking; immutable audit log; bias-monitoring jobs; HITL escalation queue; model rollback registry; access control gateway; subgroup fairness dashboard \\
\textbf{Perceived Explainability (M1)}  & SHAP attribution generator; RAG citation grounding from peer-reviewed sources; clinical-UI reasoning panel; per-feature contribution visualization \\
\textbf{Clinician Trust (M2)}           & Inspection-capable explanation rendering; override capability with rationale capture; trust-calibration feedback loop; per-clinician override audit \\
\textbf{Adoption Intention / Workflow (DV)} & EHR integration via HL7-FHIR; in-workflow decision-support display; turnaround-time monitoring; training and onboarding support tooling \\
\end{xltabular}}
\vspace{0.4em}
\noindent\textit{Note.} Pure ML engineering capabilities (CNN architecture depth, optimizer hyperparameters, augmentation internals, GPU pipeline tuning) exist in the engineering flowcharts that follow but are not part of the dissertation contribution. They are documented for engineering completeness only.

%% =========================================================================
%% PRIMARY vs SECONDARY METRICS
%% =========================================================================
\subsection*{Primary vs.\ Secondary Metrics}
\label{sec:eeg_metrics_priority}

\noindent\textbf{Primary metrics (DBA contribution):} clinician trust scores (Madsen \& Gregor TUI); adoption intention (TAM-BI); workflow-integration readiness; perceived explainability; governance maturity score (12-item RGAIG instrument).

\noindent\textbf{Secondary metrics (supporting evidence only):} classification accuracy ($\geq$~85\% pre-registered threshold); F1; ROC-AUC; SHAP attribution stability; RAG faithfulness and citation accuracy.

\noindent The dissertation's claim rests on primary metrics; technical performance metrics serve only to demonstrate that the AI system is competent enough to be a serious candidate for clinician adoption. A model that scored 80\% accuracy with strong governance-driven explainability would still be a valid test of the framework.

\noindent\textbf{End of EEG Processing Defensibility Pack.} The technical phase-by-phase flowcharts that follow are engineering supporting context, provided for completeness; examiners assessing the doctoral contribution should focus on Tables~\ref{tab:eeg_why_domain}--\ref{tab:eeg_research_technical_map} and Figure~\ref{fig:eeg_human_workflow} above.

\section{Engineering Pipeline Flowcharts}
\label{sec:appendix_ch3_engineering_flowcharts}

This section presents eight engineering-detail flowcharts that document the RGAIG pipeline at implementation level. Figures~\ref{fig:phase1_data_acquisition}--\ref{fig:phase5_model_training} trace the five core processing phases from data acquisition through model training. Figures~\ref{fig:signal_processing_flow}--\ref{fig:data_flow_e2e} provide three end-to-end pipeline overviews at increasing levels of scope.

%%--- Phase 1--5 Workflow Diagrams ---

\noindent\textbf{Phase 1: Data Acquisition and Ingestion.} Figure~\ref{fig:phase1_data_acquisition} documents the data acquisition and ingestion framework, covering multi-source data loading, format validation, quality checks, and the initial data registry that feeds all downstream phases.

\input{figures/fig_phase1_data_acquisition}

% [SPACE-FIX] clearpage removed to eliminate empty page

\noindent\textbf{Phase 2: Signal Preprocessing.} Figure~\ref{fig:phase2_signal_preprocessing} details the signal preprocessing and cleaning pipeline, including bandpass filtering, artifact rejection, re-referencing, and segment extraction that transforms raw EEG recordings into analysis-ready segments.

\input{figures/fig_phase2_signal_preprocessing}

% [SPACE-FIX] clearpage removed to eliminate empty page

\noindent\textbf{Phase 3: ASD Feature Extraction.} Figure~\ref{fig:phase3_feature_extraction} illustrates the ASD EEG feature extraction pipeline, where clean EEG segments are processed through four parallel branches---statistical, spectral, temporal, and nonlinear---yielding a 47-feature vector per sample.

\input{figures/fig_phase3_feature_extraction}

% [SPACE-FIX] clearpage removed to eliminate empty page

Figure~\ref{fig:phase4_feature_selection} presents the feature selection and dimensionality reduction funnel, which reduces the 47 extracted features to an optimised set of 25 through importance ranking, redundancy removal, and variance filtering.

\input{figures/fig_phase4_feature_selection}

% [SPACE-FIX] clearpage removed to eliminate empty page

Figure~\ref{fig:phase5_model_training} depicts the hybrid AI model training architecture, showing the dual-branch design (three baseline classifiers and three deep learning models) that converges into a majority-vote ensemble with calibrated probability outputs.

\input{figures/fig_phase5_model_training}

% [SPACE-FIX] clearpage removed to eliminate empty page

%%--- Pipeline Overview Diagrams ---

Figure~\ref{fig:signal_processing_flow} presents the 13-stage signal processing pipeline from raw EEG capture through to model-ready feature vectors, organised into four functional groups: signal acquisition, signal cleaning, signal transformation, and feature engineering.

\input{figures/fig_signal_processing_flow}

% [SPACE-FIX] clearpage removed to eliminate empty page

Figure~\ref{fig:model_training_pipeline} documents the 13-stage model training and computer vision pipeline, covering data preparation, model architecture selection, cross-validated training, and evaluation through to conceptually evaluatement-ready model output.

\input{figures/fig_model_training_pipeline}

% [SPACE-FIX] clearpage removed to eliminate empty page

Figure~\ref{fig:data_flow_e2e} provides the complete 16-stage end-to-end data flow from initial EEG capture through preprocessing, feature engineering, model inference, RAG-based explanation generation, and clinical decision support output.

\input{figures/fig_data_flow_e2e}

% [SPACE-FIX] clearpage removed to eliminate empty page

%========================================
% APPENDIX BLOCK: CHAPTER PLANNING TABLES
%========================================

\addcontentsline{toc}{section}{Chapter Planning and Alignment Tables} % [TOC-appendix-section-insertion]
\section*{Chapter Planning and Alignment Tables}

\noindent The following tables provide detailed chapter-wise planning, alignment, and readiness assessment material supporting the methodology design in Chapter~3.

\tablefontsize
\noindent Table~\ref{tab:chapter_structure_apa_appendix} presents chapter-Wise Objective, Process, and Decision Structure (Appendix Version).

\begin{longtable}
{@{} >{\raggedright\arraybackslash}p{0.7cm} >{\raggedright\arraybackslash}p{1.9cm} >{\raggedright\arraybackslash}p{1.8cm} >{\raggedright\arraybackslash}p{2.1cm} >{\raggedright\arraybackslash}p{1.8cm} >{\raggedright\arraybackslash}p{2.1cm} >{\raggedright\arraybackslash}p{1.9cm} @{}}
\caption[Chapter-Wise Objective, Process, and Decision Structure]{Chapter-Wise Objective, Process, and Decision Structure (Appendix Version)}
\label{tab:chapter_structure_apa_appendix} \\

\toprule
\thead{Chapter} & \thead{Objective} & \thead{Goal} & \thead{Task} & \thead{Input} & \thead{Output} & \thead{Decision Role} \\
\midrule
\endfirsthead
\toprule
\thead{Chapter} & \thead{Objective} & \thead{Goal} & \thead{Task} & \thead{Input} & \thead{Output} & \thead{Decision Role} \\
\midrule
\endhead
\bottomrule
\endfoot

1 & Introduce the study & Establish need and scope & Frame problem, gap, objectives, hypotheses, scope & Research problem, context, preliminary literature & Problem statement and study direction & Justifies what to address \\
2 & Review literature & Justify study intellectually & Compare studies, identify gaps, build conceptual logic & Published studies, theories, benchmarks & Gap-driven literature synthesis & Whether study is justified \\
3 & Explain methods & Show how claims are tested & Define data, variables, preprocessing, modelling, validation & Primary instruments, EEG data, method choices & Reproducible study design & Whether claims can be tested \\
4 & Present findings & Show what evidence demonstrates & Report results, test hypotheses, integrate findings & Primary results, EEG outputs, robustness evidence & Empirical findings and verdicts & Whether framework is supported \\
5 & Interpret and conclude & Translate evidence into action & Assess readiness, implications, final recommendation & Ch.~4 findings, governance logic & Adopt/pilot/defer conclusion & What should be done next \\

\end{longtable}


\tablefontsize
\noindent Table~\ref{tab:detailed_chapter_smart_matrix_appendix} documents detailed Chapter-Wise SMART, Evidence, Reliability, and Decision Matrix (Appendix Version).

\begin{longtable}
{@{} >{\raggedright\arraybackslash}p{0.4cm} >{\raggedright\arraybackslash}p{1.9cm} >{\raggedright\arraybackslash}p{0.7cm} >{\raggedright\arraybackslash}p{1.4cm} >{\raggedright\arraybackslash}p{1.4cm} >{\raggedright\arraybackslash}p{0.7cm} >{\raggedright\arraybackslash}p{0.7cm} >{\raggedright\arraybackslash}p{1.8cm} >{\raggedright\arraybackslash}p{1.8cm} >{\raggedright\arraybackslash}p{1.8cm} @{}}
\caption[Detailed Chapter-Wise SMART, Evidence, Reliability, and]{Detailed Chapter-Wise SMART, Evidence, Reliability, and Decision Matrix (Appendix Version)}
\label{tab:detailed_chapter_smart_matrix_appendix} \\

\toprule
\thead{Ch.} & \thead{Objective} & \thead{Spec.} & \thead{Measurable} & \thead{Achievable} & \thead{Rel.} & \thead{Time} & \thead{Evidence} & \thead{Reliability} & \thead{Decision Use} \\
\midrule
\endfirsthead
\toprule
\thead{Ch.} & \thead{Objective} & \thead{Spec.} & \thead{Measurable} & \thead{Achievable} & \thead{Rel.} & \thead{Time} & \thead{Evidence} & \thead{Reliability} & \thead{Decision Use} \\
\midrule
\endhead
\bottomrule
\endfoot

1 & Define problem, gap, objectives, hypotheses, scope & Yes & Problem clarity, objective alignment, scope clarity & Yes, if framing stays bounded & Yes & Yes & Problem statement, gap, stakeholder relevance, hypothesis map & Moderate--strong if duplication removed & What the thesis addresses and why \\
2 & Synthesize literature, identify gaps & Yes & Theme coverage, contradiction count, gap clarity & Yes, if focused on synthesis & Yes & Yes & Comparative review, gap matrix, theory linkage & Strong if selective and critical & Whether the study is justified \\
3 & Design methodological framework & Yes & Variables, datasets, metrics, validation, KPI thresholds & Yes, if protocol detail controlled & Yes & Yes & Method tables, EEG pipeline, variable map, governance logic & Strong if explicit and reproducible & Whether claims can be tested credibly \\
4 & Present findings, evaluate hypotheses & Yes & Accuracy, F1, AUC, trust, usability, robustness & Yes, if selective evidence displays & Yes & Yes & Performance tables, ROC, SHAP, trust findings & Very strong if clean reporting & Whether framework is supported \\
5 & Interpret findings, evaluation decisions & Yes & Readiness level, contribution quality, adoption status & Yes, if bounded and evidence-led & Yes & Yes & Readiness matrix, contribution summary, decision framework & Strong if domain-sensitive & Whether to adopt, pilot, or defer \\

\end{longtable}

\needspace{24\baselineskip}
\noindent Table~\ref{tab:chapter_decision_role_summary_appendix} summarises chapter-Wise Decision Role Summary (Appendix Version).

\paragraph{Chapter-Wise Decision Role Summary.}




{\tablestripes\tablefontsize
\needspace{20\baselineskip}
\begin{xltabular}{\textwidth}{@{} >{\raggedright\arraybackslash}p{1.5cm} L L @{}}
\caption[Chapter-Wise Decision Role Summary (Appendix Version)]{Chapter-Wise Decision Role Summary (Appendix Version)}\label{tab:chapter_decision_role_summary_appendix} \\
\toprule
\thead{Chapter} & \thead{Core Function} & \thead{Decision Role} \\
\midrule
\endfirsthead
\endhead
\bottomrule
\endlastfoot
1 & Frame the problem and define scope & Decides what the thesis must address \\
2 & Justify the study through literature and gaps & Decides why the study is needed \\
3 & Design the method and evaluation logic & Decides how the claims will be tested \\
4 & Report evidence and hypothesis results & Decides whether the claims are supported \\
5 & Interpret results and set evaluation stance & Decides what should be done in practice \\
\end{xltabular}}
\vspace{0.3em}\noindent\textit{Note.} Source: dissertation analysis; abbreviations defined in chapter prose.
\needspace{24\baselineskip}
\begin{figure}[!htbp]
\centering
\resizebox{0.97\textwidth}{!}{%
\begin{tikzpicture}[
  thesisstep/.style={
    rectangle,
    rounded corners=4pt,
    draw=ggunavy!70,
    fill=ggunavy!6,
    minimum width=3.0cm,
    minimum height=1.0cm,
    text width=2.8cm,
    align=center,
    font=\small
  },
  thesisarrow/.style={-{Stealth[length=4mm,width=3mm]}, thick}
]
\node[thesisstep] (a) at (0,0) {Chapter 1\\Problem and Scope};
\node[thesisstep] (b) at (4,0) {Chapter 2\\Literature and Gaps};
\node[thesisstep] (c) at (8,0) {Chapter 3\\Method and Evaluation};
\node[thesisstep] (d) at (12,0) {Chapter 4\\Evidence and Verdicts};
\node[thesisstep] (e) at (16,0) {Chapter 5\\Interpretation and Decision};

\draw[thesisarrow] (a) -- (b);
\draw[thesisarrow] (b) -- (c);
\draw[thesisarrow] (c) -- (d);
\draw[thesisarrow] (d) -- (e);
\end{tikzpicture}%
}
\caption{Chapter-Wise Thesis Decision Flow (Appendix Version)}
\label{fig:chapter_decision_flow_appendix}
\vspace{0.5em}
\noindent\textit{Note.} The thesis progresses from problem framing to literature justification, methodological design, empirical evidence, and final decision logic.
\end{figure}

%========================================
% APPENDIX BLOCK: ALIGNMENT, DIAGNOSTIC,
% AND READINESS TABLES
%========================================

\addcontentsline{toc}{section}{Thesis Alignment and Readiness Assessment} % [TOC-appendix-section-insertion]
\section*{Thesis Alignment and Readiness Assessment}

\tablefontsize
\noindent Table~\ref{tab:thesis_objective_alignment_appendix} presents thesis Objective Alignment Across Chapters (Appendix Version).


\begin{longtable}
{@{} >{\raggedright\arraybackslash}p{2.7cm} >{\raggedright\arraybackslash}p{1.7cm} >{\raggedright\arraybackslash}p{1.8cm} >{\raggedright\arraybackslash}p{1.9cm} >{\raggedright\arraybackslash}p{1.7cm} >{\raggedright\arraybackslash}p{2.0cm} >{\raggedright\arraybackslash}p{1.8cm} @{}}
\caption{Thesis Objective Alignment Across Chapters (Appendix Version)}
\label{tab:thesis_objective_alignment_appendix} \\

\toprule
\thead{Objective Component} & \thead{Ch.\ 1} & \thead{Ch.\ 2} & \thead{Ch.\ 3} & \thead{Ch.\ 4} & \thead{Ch.\ 5} & \thead{Overall} \\
\midrule
\endfirsthead
\toprule
\thead{Objective Component} & \thead{Ch.\ 1} & \thead{Ch.\ 2} & \thead{Ch.\ 3} & \thead{Ch.\ 4} & \thead{Ch.\ 5} & \thead{Overall} \\
\midrule
\endhead
\bottomrule
\endfoot

Define the biomedical AI problem & Strong & Supportive & Indirect & Indirect & Indirect & Strong \\
Justify why literature is insufficient & Moderate & Strong & Supportive & Indirect & Indirect & Strong \\
Develop unified ASD-focused framework & Framing only & Justification only & Strong & Tested empirically & Interpreted strategically & Strong \\
Integrate primary and secondary data & Previewed & Justified as gap & Strongly operationalized & Partially evidenced & Interpreted in deploy terms & Strong \\
Evaluate technical EEG performance & Previewed early & Justified & Strongly designed & Strongly evidenced & Interpreted & Strong \\
Evaluate explainability and trust & Framed strongly & Justified strongly & Methodologically supported & Evidenced & Interpreted for adoption & Strong \\
Evaluate governance and safe evaluation & Framed strongly & Justified conceptually & Strongly designed & Partly evidenced & Heavily interpreted & Moderate--strong \\
Evaluate workflow and operational value & Framed strongly & Justified moderately & Methodologically included & Partly evidenced & Strongly interpreted & Moderate \\
Support B2B hospital decision making & Framed strongly & Justified moderately & Operationalized & Partly evidenced & Strongly interpreted & Strong \\
Bound B2C as extension or future path & Over-expanded & Moderately justified & Included but broad & Limited evidence & Interpreted but broad & Moderate \\
Produce evidence-maturity interpretation decision logic & Previewed early & Not central & Designed & Partially evidenced & Main destination & Moderate--strong \\

\end{longtable}
\normalsize

\noindent Table~\ref{tab:chapter_vs_thesis_matrix_appendix} documents chapter Objective versus Thesis Objective Matrix (Appendix Version).

\begin{longtable}
{@{} >{\raggedright\arraybackslash}p{1.2cm} >{\raggedright\arraybackslash}p{3.4cm} >{\raggedright\arraybackslash}p{3.4cm} >{\raggedright\arraybackslash}p{1.8cm} >{\raggedright\arraybackslash}p{3.0cm} @{}}
\caption{Chapter Objective versus Thesis Objective Matrix (Appendix Version)}
\label{tab:chapter_vs_thesis_matrix_appendix} \\

\toprule
\thead{Chapter} & \thead{Chapter Objective} & \thead{Link to Thesis Objective} & \thead{Strength} & \thead{Main Problem} \\
\midrule
\endfirsthead
\toprule
\thead{Chapter} & \thead{Chapter Objective} & \thead{Link to Thesis Objective} & \thead{Strength} & \thead{Main Problem} \\
\midrule
\endhead
\bottomrule
\endfoot

Ch.\ 1 & Introduce problem, gap, objectives, hypotheses, scope & Defines why the thesis exists & Moderate--strong & Method and architecture creep in introduction \\
Ch.\ 2 & Synthesize literature and identify gaps & Justifies why the thesis is needed & Strong & Synthesis weakened by duplication \\
Ch.\ 3 & Explain framework and study design & Operationalizes how thesis objective will be tested & Strong & Architecture and protocol overload \\
Ch.\ 4 & Present primary, secondary, and integrated evidence & Shows whether thesis objective is empirically supported & Very strong & Too many tables and support displays \\
Ch.\ 5 & Interpret findings and make final decisions & Explains what thesis objective means in practice & Moderate--strong & Decision diluted by adjacent frameworks \\

\end{longtable}

\noindent Table~\ref{tab:chapter_strengths_weaknesses_action_appendix} presents chapter-Wise Strengths, Weaknesses, and Action Plan (Appendix Version).

\begin{longtable}
{@{} >{\raggedright\arraybackslash}p{1.5cm} >{\raggedright\arraybackslash}p{4cm} >{\raggedright\arraybackslash}p{4cm} >{\raggedright\arraybackslash}p{4cm} @{}}
\caption{Chapter-Wise Strengths, Weaknesses, and Action Plan (Appendix Version)}
\label{tab:chapter_strengths_weaknesses_action_appendix} \\

\toprule
\thead{Chapter} & \thead{Strengths} & \thead{Weaknesses} & \thead{Action Plan} \\
\midrule
\endfirsthead
\multicolumn{4}{l}{\textit{Table~\ref{tab:chapter_strengths_weaknesses_action_appendix} continued}} \\
\midrule
\endhead
\bottomrule
\endfoot

Ch.\ 1 & Strong problem framing, stakeholder relevance, broad scope & Overloaded introduction, methods creep, workflow clutter & Cut operational detail; move method content to Ch.~3; retain framing only \\
Ch.\ 2 & Broad evidence base, meaningful gap identification & Duplication, inventory-style summaries, too many side frameworks & Tighten synthesis; remove duplicates; focus on gap-to-study logic \\
Ch.\ 3 & Strong methodological ambition, integrated design, good coverage & Bloated chapter, too many master tables, protocol overload & One method flow per strand; move protocols to appendix \\
Ch.\ 4 & Strongest empirical chapter, explicit hypothesis testing & Too many support tables, repeated displays, discussion leakage & One table + one figure per claim; move helpers to appendix \\
Ch.\ 5 & Rich interpretation, strong DBA orientation, practical relevance & Too many matrices, policy sprawl, weak prioritization & Compress around readiness, ROI, evidence-maturity interpretation conclusion \\

\end{longtable}
\needspace{24\baselineskip}
\noindent Table~\ref{tab:adopt_pilot_defer_readiness} documents adopt, Pilot, or Defer Readiness Matrix.

\paragraph{Adopt, Pilot, or Defer.}




{\tablestripes\tablefontsize
\needspace{20\baselineskip}
\begin{xltabular}{\textwidth}{@{} >{\raggedright\arraybackslash}p{3.0cm} >{\centering\arraybackslash}p{2.1cm} >{\centering\arraybackslash}p{2.1cm} >{\centering\arraybackslash}p{2.1cm} L @{}}
\caption[Adopt, Pilot, or Defer Readiness Matrix]{Adopt, Pilot, or Defer Readiness Matrix}\label{tab:adopt_pilot_defer_readiness} \\
\toprule
\thead{Evaluation Area} & \thead{Adopt} & \thead{Pilot} & \thead{Defer} & \thead{Interpretation} \\
\midrule
\endfirsthead
\endhead
\bottomrule
\endlastfoot
Technical performance & High & Moderate to high & Low & Accuracy, recall, F1, AUC, and robustness status \\
Explainability and trust & High & Moderate & Low & SHAP clarity, RAG usefulness, expert trust, and user understanding \\
Governance readiness & High & Moderate & Low & Privacy, auditability, access control, and human-in-the-loop status \\
Workflow fit & High & Moderate & Low & Time reduction, usability, process fit, and burden reduction \\
Domain readiness & Strong & Conditional & Weak & Domain-specific maturity and validation sufficiency \\
Overall decision & Deploy selectively & Controlled trial use & Hold for future work & Final evaluation stance based on combined evidence \\
\end{xltabular}}
\vspace{0.3em}\noindent\textit{Note.} SHAP = SHapley Additive exPlanations; RAG = Retrieval-Augmented Generation; AUC = Area Under (ROC) Curve. See chapter prose for full context.
\needspace{24\baselineskip}
\noindent Table~\ref{tab:domain_sensitive_decision_matrix} documents domain-Sensitive Evaluation Decision Matrix.

\paragraph{Domain-Sensitive Evaluation Decision Matrix.}




{\tablestripes\tablefontsize
\needspace{20\baselineskip}
\begin{xltabular}{\textwidth}{@{} >{\raggedright\arraybackslash}p{2.7cm} >{\centering\arraybackslash}p{2.0cm} >{\centering\arraybackslash}p{2.0cm} >{\centering\arraybackslash}p{2.0cm} L @{}}
\caption[Domain-Sensitive Evaluation Decision Matrix]{Domain-Sensitive Evaluation Decision Matrix}\label{tab:domain_sensitive_decision_matrix} \\
\toprule
\thead{Domain} & \thead{Technical} & \thead{Trust} & \thead{Governance} & \thead{Recommended Decision} \\
\midrule
\endfirsthead
\endhead
\bottomrule
\endlastfoot
ASD & Strong & Strong & Moderate to strong & Pilot or selective adopt \\
\end{xltabular}}
\vspace{0.3em}\noindent\textit{Note.} ASD = Autism Spectrum Disorder. See chapter prose for full context.

% [SPACE-FIX] clearpage removed to eliminate empty page

%========================================
% CHAPTER-WISE READINESS Q&A TABLE
%========================================
%% MOVED TO STANDALONE: chapter_readiness_qa.pdf
%% Per user directive 2026-05-11 — readiness Q&A is defense-prep, not thesis evidence.
%% Phantom label kept so cross-references resolve to "??" rather than error.
\phantomsection\label{tab:detailed_chapter_qa}

\iffalse %% SUPPRESSED — moved to chapter_readiness_qa.tex (standalone)

\addcontentsline{toc}{section}{Chapter-Wise Readiness Questions and Answers} % [TOC-appendix-section-insertion]
\section*{Chapter-Wise Readiness Questions and Answers}

\noindent Table~\ref{tab:detailed_chapter_qa} provides a comprehensive set of likely examiner questions for each chapter, with structured answers prepared from the thesis content.

\renewcommand{\arraystretch}{1.15}

\begin{longtable}
{@{} >{\raggedright\arraybackslash}p{1.6cm} >{\raggedright\arraybackslash}p{4.6cm} >{\raggedright\arraybackslash}p{8.0cm} @{}}
\caption{Detailed Chapter-Wise Readiness Questions and Answers}
\\

\toprule
\thead{Chapter} & \thead{Likely Question} & \thead{Prepared Answer} \\
\midrule
\endfirsthead

\toprule
\thead{Chapter} & \thead{Likely Question} & \thead{Prepared Answer} \\
\midrule
\endhead

\bottomrule
\endfoot

Chapter 1 & What is the exact research problem? &
Biomedical AI systems are often accurate only in isolated settings. Existing systems remain fragmented across diseases, tools, and workflows. Many models are weakly explainable and difficult to trust. Governance, accountability, and workflow integration are often insufficient. The real problem is not only prediction quality, but evidence-maturity interpretation. \\
\addlinespace

Chapter 1 & Why is this problem important? &
Healthcare organisations need systems that support real decisions. Accuracy alone is not enough in clinical environments. Poor explainability reduces clinician trust and adoption. Weak workflow fit increases burden. Governance failures create safety, privacy, and accountability risks. \\
\addlinespace

Chapter 1 & What gap does the thesis address? &
Prior studies often focus on one disease or one model only. Few studies integrate technical, human, and governance dimensions together. ASD-focused readiness remains weakly tested. Explainability is often technical rather than deployment-oriented. This thesis addresses the lack of a unified, explainable, governance-aware framework. \\
\addlinespace

Chapter 1 & What is the main objective? &
To develop a ASD AI framework, evaluate technical EEG-based performance across domains, integrate primary and secondary evidence, examine explainability, trust, workflow, and governance readiness, and translate findings into evidence-maturity interpretation logic. \\
\addlinespace

Chapter 1 & What is the novelty? &
It does not stop at single-domain benchmarking. It combines technical, human-centred, and governance logic. It uses a ASD-focused evaluation perspective. It includes decision-oriented evaluation logic. It treats readiness as multidimensional rather than purely technical. \\
\addlinespace

Chapter 1 & What is the scope? &
The primary scope is supervised B2B clinical decision support. B2C is a bounded extension, not the main validated use case. The thesis focuses on readiness, not full commercial rollout. Domain-sensitive interpretation is central to scope control. \\
\addlinespace

Chapter 2 & How was the literature selected? &
The review was structured around core thesis themes. Sources were selected for relevance to EEG diagnostics and AI readiness. The review included both technical and managerial literature. Priority was given to studies that inform the framework logic. \\
\addlinespace

Chapter 2 & What does the literature show? &
EEG-based systems can perform strongly in some domains. Explainability is increasingly recognised as important. Governance concerns are widely discussed. Workflow and adoption issues are acknowledged but often underdeveloped. \\
\addlinespace

Chapter 2 & What are the major weaknesses in prior work? &
Many studies are single-domain only. Explainability is often added after modelling rather than embedded in workflow. Governance is often discussed conceptually, not operationally. Human-centred evidence is often weak. Real evidence-maturity interpretation is rarely tested systematically. \\
\addlinespace

Chapter 2 & What are the main research gaps? &
Lack of ASD-focused unification. Lack of integrated technical and human-centred evaluation. Weak governance-ready biomedical AI design. Weak workflow-oriented explainability. Limited adoption-oriented readiness assessment. \\
\addlinespace

Chapter 2 & How does Chapter~2 justify Chapter~3? &
It identifies what existing studies do not solve. It shows why mixed methods are needed. It justifies why both technical and human-centred data are included. It supports the inclusion of governance and explainability layers. It provides the logic for the methodological response. \\
\addlinespace

Chapter 3 & Why was this research design chosen? &
A single-method design would not capture full readiness. Technical evidence alone does not show adoption potential. Human-centred evidence alone does not prove technical viability. Mixed methods support integrated judgement. The design fits the thesis objective directly. \\
\addlinespace

Chapter 3 & What are the data sources? &
Primary data include human-centred instruments and structured feedback. Secondary data include EEG-related biomedical datasets. The study combines both strands deliberately to support both technical and practical analysis. \\
\addlinespace

Chapter 3 & How was EEG data handled? &
Raw EEG underwent signal-quality assessment, preprocessing (filtering, artefact handling), segmentation, normalisation, and transformation (FFT, PSD, SWT). The output became model-ready input for comparative benchmarking. \\
\addlinespace

Chapter 3 & How was performance validated? &
Validation used subject-safe splitting to reduce leakage, cross-validation for credibility, LOSO for stronger generalisation, and robustness methods to reduce reliance on one favourable result. \\
\addlinespace

Chapter 3 & How were primary and secondary findings integrated? &
Both strands were analysed separately first. Integration used joint display and triangulation. Convergence and divergence were treated analytically. Meta-inference linked the strands to readiness decisions. \\
\addlinespace

Chapter 4 & What are the main findings? &
The framework performs strongly in most included domains. Performance is not equally strong for ASD. Results support the framework in a qualified rather than universal sense. Domain variation is a central empirical finding. \\
\addlinespace

Chapter 4 & Which hypotheses were supported? &
Some hypotheses were strongly supported. Some were partially supported depending on domain or evidence type. Each verdict links to measured evidence. The thesis avoids overclaiming full support where evidence is mixed. \\
\addlinespace

Chapter 4 & Which domain performed best? &
The best-performing domain is identified clearly with both performance and stability evidence. Strong performance is still interpreted cautiously if sample limits exist. Best performance does not automatically mean full evidence-maturity interpretation. \\
\addlinespace

Chapter 4 & Which domain performed worst? &
The weakest domain functions as a boundary condition. Its lower performance helps define framework limits. This strengthens honesty and credibility and informs the defer decision. \\
\addlinespace

Chapter 4 & Are the findings robust? &
Robustness depends on stability across repeated evaluations. Cross-validation, LOSO, Monte Carlo, and ablation strengthen trust in the results. Robust findings support stronger readiness conclusions. \\
\addlinespace

Chapter 4 & Did technical and human-centred findings align? &
Alignment occurred in some domains more than others. Convergence strengthens readiness claims. Divergence reveals adoption or trust barriers. This prevents equating accuracy with readiness. \\
\addlinespace

Chapter 5 & What do the findings mean overall? &
The framework is promising but not universally mature. Results support selective rather than blanket readiness. Technical, human-centred, and governance evidence must be read together. The thesis supports a bounded evaluation stance. \\
\addlinespace

Chapter 5 & What are the contributions? &
Theoretical: reframes readiness as multidimensional. Methodological: mixed-methods ASD-focused design. Practical: decision-oriented evaluation logic with evidence-maturity interpretation framework. \\
\addlinespace

Chapter 5 & Is the framework ready for deployment? &
It is selectively ready, not universally ready. Stronger domains support pilot or limited adoption. Weaker domains should be deferred. Governance and trust conditions must remain in place. Human review remains essential. \\
\addlinespace

Chapter 5 & What is the B2B conclusion? &
B2B is the primary validated pathway. Hospital-facing supervised evaluation is the strongest practical conclusion. Pilot-first logic is appropriate. \\
\addlinespace

Chapter 5 & What is the B2C conclusion? &
B2C should remain a bounded extension pathway. It may support future continuity-of-care applications. Full autonomous consumer evaluation is not yet justified. \\
\addlinespace

Chapter 5 & What are the main limitations? &
Domain maturity is uneven. Validation remains bounded by available data. Operational evaluation evidence is still limited. Some conclusions are pilot-oriented rather than adoption-level. Real-world testing remains necessary. \\
\addlinespace

Chapter 5 & What is the final thesis conclusion? &
The framework should be interpreted as selectively and conditionally ready. Stronger domains support pilot or limited supervised adoption. Weaker domains support deferral. B2B is the strongest evaluation pathway. Final claims remain evidence-led and bounded. \\

\end{longtable}

\renewcommand{\arraystretch}{1.0}

\fi %% END SUPPRESSED readiness Q&A — see chapter_readiness_qa.pdf

%% ================================================================
%% Migrated from Chapter 3 main body (engineering detail)
%% ================================================================
\section{Engineering Detail Migrated from Chapter~3 Main Body}\label{sec:appendix_ch3_migrated}
The following tables and figures present full engineering detail supporting the DBA research methods narrative in Chapter~3. They were relocated here so the main chapter reads as a complete research-methods story while preserving all technical evidence for examination.

\subsection{Master Logic Chain}\label{sec:appendix_ch3_master_logic_chain}
\needspace{24\baselineskip}
{%
\tablefontsize

\noindent Table~\ref{tab:master_logic_chain} presents master Research Logic: Gap to Decision Chain.


\needspace{20\baselineskip}
\begin{xltabular}
{\textwidth}{@{} L >{\raggedright\arraybackslash}p{1cm} L >{\raggedright\arraybackslash}p{2cm} L L @{}}
\caption[Master Research Logic]{Master Research Logic: Gap to Decision Chain} \label{tab:master_logic_chain} \\
\toprule
\thead{Research Gap} & \thead{Obj.} & \thead{Variable} & \thead{Analysis} & \thead{Hypothesis / Logic} & \thead{Decision Implication} \\
\midrule
\endfirsthead
\endhead
\bottomrule
\endlastfoot
No integration of primary and secondary evidence & O7, O11 & Primary--secondary alignment, contextual completeness & Correlation, concordance, joint display & If evidence aligns, integrated workflow justified & Retain integrated assessment design \\
\addlinespace
Behavioural/psychological assessment not captured in AI workflows & O7, O9 & Behavioural, psychological, functional scores & Descriptive, reliability, severity profiling & If stable scores derivable, primary data adds valid context & Include behavioural/psychological layer \\
\addlinespace
Clinical AI explainability weak and hard to trust & O8 & Trust, clarity, usefulness scores & Descriptive, Cronbach's $\alpha$, regression & H5: Better explanation quality increases trust & Keep or improve explanation module \\
\addlinespace
Hospital onboarding fragmented and inefficient & O5, O6 & Usability, completeness scores, time & Descriptive, group comparison & H4: Guided onboarding improves completeness & Adopt chatbot intake if thresholds met \\
\addlinespace
Existing systems ignore caregiver perspective (ASD) & O9 & Caregiver behavioural, developmental scores & Descriptive, reliability, correlation & Caregiver patterns provide interpretable context & Use caregiver data as formal primary input \\
\addlinespace
Patient symptom burden not linked to AI outputs & O9, O11 & Symptom severity, model confidence & Correlation, cross-tabulation & Higher burden may align with disease-specific output & Use patient report as contextual support \\
\addlinespace
AI workflows lack human acceptance focus & O8, O15 & Adoption, usability, trust scores & Descriptive, regression, triangulation & If acceptance low, technical performance insufficient & Pilot only, not broad evaluation \\
\addlinespace
Primary-data instruments not validated & O9 & Internal consistency of subscales & Cronbach's $\alpha$, item-total & Scales should meet acceptable reliability & Retain stable scales, revise weak ones \\
\addlinespace
No evidence context improves decision usefulness & O11, O15 & Contextual usefulness, expert rating & Comparison, regression, joint display & Integrated context improves decision usefulness & Support mixed-method justification \\
\addlinespace
B2C continuity proposed but not assessed & O14 & Monitoring readiness, privacy confidence & Descriptive, thematic interpretation & If readiness weak, B2C remains future scope & Keep B2C as secondary extension \\
\addlinespace
Governance requires user trust and privacy comfort & O12 & Privacy confidence, trust, explanation acceptance & Descriptive, threshold testing & Acceptable governance scores support evaluation & Move toward pilot if thresholds met \\
\addlinespace
Prior studies lack mixed-method validation & O11, O15 & Convergent/divergent evidence & Joint display, meta-inference & If quant and qual converge, evaluation case strengthens & Final evidence-maturity interpretation decision \\
\end{xltabular}
\vspace{0.3em}\noindent\textit{Note.} ASD = Autism Spectrum Disorder. See chapter prose for full context.
}%

\subsection{ASD Questionnaire Blueprint}\label{sec:appendix_ch3_asd_questionnaire_blueprint}
\clearpage
\noindent Table~\ref{tab:asd_questionnaire_blueprint} presents the ASD Parent Questionnaire Blueprint, specifying respondent roles, item counts, and output variables across the seven assessment sections.

{\tablestripes\tablefontsize
\begin{xltabular}{\textwidth}{@{} >{\raggedright\arraybackslash}p{3.7cm} >{\raggedright\arraybackslash}p{2.7cm} p{5.5cm} >{\centering\arraybackslash}p{1.0cm} >{\raggedright\arraybackslash}p{1.0cm} @{}}
\caption[ASD Parent Questionnaire Blueprint]{ASD Parent Questionnaire Blueprint: Sections, Items, and Output Variables}\label{tab:asd_questionnaire_blueprint} \\
\toprule
\thead{Section} & \thead{Respondent} & \thead{Focus} & \thead{Items} & \thead{Output Variable} \\
\midrule
\endfirsthead
\endhead
\bottomrule
\endlastfoot
Develop.\ hist. & Parent & Speech delay, milestones, onset of concern & 5--8 & \texttt{asd\_\allowbreak dev\_\allowbreak score} \\
Social comm. & Parent & Eye contact, name response, peer engagement & 6--8 & \texttt{asd\_\allowbreak social\_\allowbreak score} \\
Repet.\ behav. & Parent & Routines, repetitive actions, fixation & 4--6 & \texttt{asd\_\allowbreak repet\_\allowbreak score} \\
Sensory proc. & Par./Pat. & Sound, touch, taste, visual overload & 4--6 & \texttt{asd\_\allowbreak sensory\_\allowbreak score} \\
Emotional reg. & Par./Pat. & Meltdown frequency, frustration, recovery & 4--6 & \texttt{asd\_\allowbreak emotion\_\allowbreak score} \\
Function.\ impact & Parent & School, home, communication, social & 4--5 & \texttt{asd\_\allowbreak function\_\allowbreak score} \\
Self-report & Patient & Stress, anxiety, overload, comfort & 3--5 & \texttt{asd\_\allowbreak selfrep\_\allowbreak score} \\
\midrule
\multicolumn{3}{@{}l}{\textbf{Total estimated items}} & \textbf{30--44} & --- \\
\end{xltabular}}
\vspace{2pt}
\noindent\textit{Note.} Par.\ = Parent; Pat.\ = Patient. All structured items use 5-point Likert or frequency scales (0 = Never to 4 = Very Often). Composite scores are computed as section means. Reliability target: Cronbach's $\alpha \geq 0.70$.

\subsection{PD Questionnaire Blueprint (Comparator)}\label{sec:appendix_ch3_pd_questionnaire_blueprint}
\needspace{24\baselineskip}
\noindent Table~\ref{tab:pd_questionnaire_blueprint} presents the Parkinson's Disease (PD) Patient Questionnaire Blueprint, included here as a methodological comparator illustrating how the ASD instrument's section/respondent/item/output structure generalises to an adult-onset neurodegenerative condition. The PD instrument is not administered as part of this dissertation; the comparator establishes content-validity portability of the questionnaire design pattern.

\noindent Table~\ref{tab:pd_questionnaire_blueprint} pd questionnaire blueprint (comparator).

{\tablestripes\tablefontsize
\begin{xltabular}{\textwidth}{@{} >{\raggedright\arraybackslash}p{3.8cm} >{\raggedright\arraybackslash}p{2.7cm} p{5.5cm} >{\centering\arraybackslash}p{1.0cm} >{\raggedright\arraybackslash}p{1.0cm} @{}}
\caption[PD Questionnaire Blueprint (Comparator)]{PD Patient Questionnaire Blueprint (Comparator): Sections, Items, and Output Variables}\label{tab:pd_questionnaire_blueprint} \\
\toprule
\thead{Section} & \thead{Respondent} & \thead{Focus} & \thead{Items} & \thead{Output Variable} \\
\midrule
\endfirsthead
\endhead
\bottomrule
\endlastfoot
Motor symptoms     & Patient   & Tremor, rigidity, slowness, gait, balance       & 5--8 & \texttt{pd\_\allowbreak motor\_\allowbreak score} \\
Non-motor sympt.   & Patient   & Fatigue, sleep, constipation, smell loss        & 4--6 & \texttt{pd\_\allowbreak nonmotor\_\allowbreak score} \\
Cognitive sympt.   & Patient   & Memory, concentration, planning difficulty      & 4--5 & \texttt{pd\_\allowbreak cognitive\_\allowbreak score} \\
Emotional sympt.   & Patient   & Anxiety, depression, apathy, frustration        & 4--5 & \texttt{pd\_\allowbreak emotional\_\allowbreak score} \\
Daily living       & Patient   & Dressing, eating, movement, speaking            & 4--6 & \texttt{pd\_\allowbreak adl\_\allowbreak score} \\
Caregiver obs.     & Caregiver & Falls, freezing, mood shifts, variability       & 3--5 & \texttt{pd\_\allowbreak caregiver\_\allowbreak score} \\
\midrule
\multicolumn{3}{@{}l}{\textbf{Total estimated items}} & \textbf{24--35} & --- \\
\end{xltabular}}
\vspace{2pt}
\noindent\textit{Note.} Comparator instrument; not administered in this dissertation. Motor and non-motor sections use 5-point severity scales (1 = Very Low to 5 = Very High). Cognitive and emotional sections use 5-point agreement scales. Reliability target: Cronbach's $\alpha \geq 0.70$.

\subsection{EEG Analysis Pipeline}\label{sec:appendix_ch3_eeg_analysis_pipeline}

\noindent\textbf{Secondary EEG Data: Quantitative Analysis Pipeline.} Table~\ref{tab:eeg_analysis_pipeline} below presents the detailed analysis.

{\tablestripes\tablefontsize

\noindent Table~\ref{tab:eeg_analysis_pipeline} presents the secondary eeg data.

\needspace{20\baselineskip}
\begin{xltabular}{\textwidth}{@{} >{\raggedright\arraybackslash}p{1.2cm} >{\raggedright\arraybackslash}p{3.8cm} L @{}}
\caption[Secondary EEG Data]{Secondary EEG Data: Quantitative Analysis Pipeline}\label{tab:eeg_analysis_pipeline} \\
\toprule
\thead{Step} & \thead{Action} & \thead{Output} \\
\midrule
\endfirsthead
\endhead
\bottomrule
\endlastfoot
1 & Acquire or load EEG data & Analysis dataset \\
2 & Clean artefacts and noise & Usable EEG segments \\
3 & Segment / epoch data & Comparable analysis windows \\
4 & Extract features or generate deep representations & Model-ready inputs \\
5 & Split data for training/validation/testing & Valid evaluation design \\
6 & Train baseline and advanced models & Comparative results \\
7 & Compute performance metrics & Accuracy, AUC, F1, sensitivity, specificity \\
8 & Test robustness & Confidence intervals, ablation, LOSO/CV \\
9 & Run explainability analysis & Feature importance / signal relevance \\
10 & Translate results into decision usefulness & Adopt / pilot / defer logic \\
\end{xltabular}}
\vspace{0.3em}\noindent\textit{Note.} EEG = Electroencephalography; AUC = Area Under (ROC) Curve. See chapter prose for full context.

\subsection{EEG 40-Step Pipeline}\label{sec:appendix_ch3_eeg_40step_pipeline}
{%
\tablefontsize

\noindent Table~\ref{tab:eeg_40step_pipeline} presents complete End-to-End EEG Pipeline: 40-Step Methodology.


\needspace{20\baselineskip}
\begin{xltabular}
{\textwidth}{@{} r >{\raggedright\arraybackslash}p{2.1cm} L L L @{}}
\caption[Complete End-to-End EEG Pipeline]{Complete End-to-End EEG Pipeline: 40-Step Methodology} \label{tab:eeg_40step_pipeline} \\
\toprule
\thead{\#} & \thead{Stage} & \thead{What Happens} & \thead{Why It Matters} & \thead{Risk if Missing} \\
\midrule
\endfirsthead
\endhead
\bottomrule
\endlastfoot
1 & Problem defn & Define disease, use case, stakeholder, context & Gives pipeline clear purpose & Technically strong but academically weak \\
2 & Data source defn & Identify EEG source: Emotiv/IoT, dataset, hospital & Avoids ambiguity about origin & Data credibility weak \\
3 & Consent \& gov. & Patient consent, ethics, access, audit rules & Required for clinical/mobile use & Unsafe/non-compliant claim \\
4 & Patient onboard. & Registration, symptom intake, parent input & Connects human context to EEG & No human context for decisions \\
5 & Signal acquisition & Collect multichannel EEG & Core biomedical input & No biomedical evidence \\
6 & Signal quality & Check impedance, missing ch., noisy segments & Prevents bad data in pipeline & Noisy data undermines findings \\
7 & Artefact removal & Remove blink, muscle, motion, powerline & Improves reliability & Inflated error, unstable features \\
8 & Band-pass filter & Retain useful frequencies (0.5--45 Hz) & Standard preprocessing & Irrelevant noise remains \\
9 & Segment./epoch. & Divide EEG into windows (2--5 s) & Creates comparable units & Inconsistent analysis units \\
10 & Baseline corr. & Normalise reference/baseline activity & Improves comparability & Distorted channel interp. \\
11 & Time-domain & Hjorth parameters, variance, mean, RMS & Captures signal behaviour & Misses waveform-level info \\
12 & Freq.-domain & FFT, PSD/SPDD & Extracts spectral content & No band-power interp. \\
13 & Time-freq. & SWT, wavelet, scalogram & Analyses changing signal & Transient patterns missed \\
14 & Feature extr. & PSD, entropy, coherence, wavelets, deep & Converts to model-ready form & Raw data unusable \\
15 & Feature select. & PCA, mRMR, LASSO, importance ranking & Keeps strongest, least redundant & Overfitting, redundancy \\
16 & Normalisation & Z-score, min-max scaling & Standardises magnitude & Unstable training \\
17 & 1D to 2D conv. & Spectrogram, scalogram, topomap & Enables CNN/CV models & CV models unavailable \\
18 & Safe splitting & Subject-wise split, LOSO, CV & Prevents leakage & False perf.\ inflation \\
19 & Class imbalance & Weighted loss, oversampling, SMOTE & Reduces bias from skew & Misleading accuracy \\
20 & Model selection & ML, 1D CNN, 2D CNN, transformer, ensemble & Choose modelling family & Arbitrary arch.\ choice \\
21 & Hyperparam.\ tune & Grid search, Bayesian tuning & Optimise training & Weak/unstable performance \\
22 & Training & Fit model to data & Core modelling step & No predictive system \\
23 & Validation & Monitor development performance & Tuning feedback & Overfitting risk \\
24 & Testing & Final unbiased evaluation & Proves performance & No trustworthy claim \\
25 & Core evaluation & Confusion matrix, recall, acc., F1, AUC & Technical evidence & No hypothesis testing basis \\
26 & Threshold tune & ROC thresholding, calibration curve & Align to screening use & Unsafe FP/FN balance \\
27 & Robustness & Bootstrap, ablation, LOSO & Test stability & Results may be fragile \\
28 & Explainability & SHAP, saliency, attention, feature maps & Interpret model behaviour & Black-box untrusted \\
29 & Clinical plaus. & Compare features with domain knowledge & Plausible biomed.\ interp. & Explain.\ superficial \\
30 & RAG interp. & Retrieval-augmented generation & Evidence-backed summary & Expl.\ generic/unsupported \\
31 & Output format & Clinician, patient, admin views & Usable result presentation & Same output confuses users \\
32 & Decision logic & Risk rules, escalation logic & Review/monitor/escalate & Prediction $\neq$ action \\
33 & MCP/orchestr. & Model context protocol workflow & Coordinated execution & Disconnected components \\
34 & Mobile app & App onboarding, reports, monitoring & Patient-facing continuity & Weak B2C continuity \\
35 & Hospital integr. & EHR/HIS/dashboard/report API & Clinical workflow adoption & Hard to conceptually evaluate \\
36 & SOS/escalation & Alert engine, escalation workflow & Safety response pathway & High-risk cases missed \\
37 & HITL & Review dashboard, override mechanism & Safe decision-support use & Overclaim as autonomous \\
38 & Logging/audit & Audit logs, traceability layer & Governance evidence & Weak compliance \\
39 & ROI/workflow & Time-motion, throughput, effort & Business case for adoption & B2B argument incomplete \\
40 & Final decision & Adopt/pilot/defer matrix & Readiness conclusion & No usable conclusion \\
\end{xltabular}
\vspace{0.3em}\noindent\textit{Note.} SHAP = SHapley Additive exPlanations; RAG = Retrieval-Augmented Generation; HITL = Human-in-the-loop; EEG = Electroencephalography. See chapter prose for full context.
}%

\subsection{EEG Preprocessing Detail}\label{sec:appendix_ch3_eeg_preprocessing}
\noindent Table~\ref{tab:eeg_preprocessing} presents eEG Preprocessing and Signal Analysis Methods.


{\tablestripes\tablefontsize
\needspace{20\baselineskip}
\begin{xltabular}{\textwidth}{@{} >{\raggedright\arraybackslash}p{2.2cm} >{\raggedright\arraybackslash}p{2.4cm} L L @{}}
\caption[EEG Preprocessing and Signal Analysis Methods]{EEG Preprocessing and Signal Analysis Methods}\label{tab:eeg_preprocessing} \\
\toprule
\thead{Step} & \thead{Method} & \thead{Purpose} & \thead{Example Use} \\
\midrule
\endfirsthead
\endhead
\bottomrule
\endlastfoot
Noise removal & Band-pass filter & Remove irrelevant frequency components & Keep 0.5--40 Hz \\
Artefact handling & ICA, filtering, rejection & Improve usable signal quality & Cleaner epochs \\
Segmentation & Epoching/windowing & Create comparable units & 2 s or 5 s windows \\
Spectral analysis & FFT / PSD & Quantify alpha, beta, theta, delta power & Biomarker profiling \\
Time-frequency & SWT / wavelet & Detect transient and local changes & ASD/PD signal irregularities \\
Normalisation & Z-score, min-max & Standardise features & Stable model input \\
Representation & 1D to 2D mapping & Allow image-based modelling & Spectrogram, scalogram, topomap \\
\end{xltabular}}
\vspace{0.3em}\noindent\textit{Note.} ASD = Autism Spectrum Disorder; EEG = Electroencephalography. See chapter prose for full context.

\subsection{EEG Feature Types}\label{sec:appendix_ch3_eeg_feature_types}

\noindent\textbf{EEG Feature Extraction: Types and Applications.} Table~\ref{tab:eeg_feature_types} below presents the detailed analysis.

{\tablestripes\tablefontsize

\noindent Table~\ref{tab:eeg_feature_types} presents the eeg feature extraction.

\needspace{20\baselineskip}
\begin{xltabular}{\textwidth}{@{} >{\raggedright\arraybackslash}p{2cm} L L @{}}
\caption[EEG Feature Extraction]{EEG Feature Extraction: Types and Applications}\label{tab:eeg_feature_types} \\
\toprule
\thead{Feature Type} & \thead{Examples} & \thead{Why Useful} \\
\midrule
\endfirsthead
\endhead
\bottomrule
\endlastfoot
Time-domain & Mean, variance, Hjorth parameters & Simple signal characterisation \\
Frequency-domain & Band power, spectral entropy, peak frequency & Strong EEG biomarker use \\
Time-frequency & Wavelet coefficients, scalogram features & Capture dynamic signal changes \\
Connectivity & Coherence, correlation, phase-locking & Inter-channel relationships \\
Nonlinear & Entropy, fractal dimension & Complex signal behaviour \\
Deep features & CNN latent embeddings & Automatic learned representations \\
\end{xltabular}}
\vspace{0.3em}\noindent\textit{Note.} EEG = Electroencephalography. See chapter prose for full context.

\subsection{EEG Modelling Strategies}\label{sec:appendix_ch3_eeg_modelling_strategies}

\noindent\textbf{EEG Modelling Strategies: Input, Strength, and....} Table~\ref{tab:eeg_modelling_strategies} below presents the detailed analysis.

{\tablestripes\tablefontsize

\noindent Table~\ref{tab:eeg_modelling_strategies} presents the eeg modelling strategies.

\needspace{20\baselineskip}
\begin{xltabular}{\textwidth}{@{} >{\raggedright\arraybackslash}p{2.0cm} >{\raggedright\arraybackslash}p{2.0cm} L L @{}}
\caption[EEG Modelling Strategies]{EEG Modelling Strategies: Input, Strength, and Application}\label{tab:eeg_modelling_strategies} \\
\toprule
\thead{Strategy} & \thead{Input Type} & \thead{Best Use} & \thead{Strength} \\
\midrule
\endfirsthead
\endhead
\bottomrule
\endlastfoot
Classical ML & Extracted features & Smaller structured datasets & Interpretable baseline \\
1D CNN & Raw/filtered series & Temporal EEG patterns & Direct signal learning \\
LSTM / transf. & Sequences & Temporal dependency & Sequence-aware modelling \\
2D CNN & Spectrogram, scalogram & Image-like EEG repr. & Strong visual pattern learning \\
CV model & 2D converted EEG & Complex spatial-freq.\ patterns & After 1D to 2D conversion \\
Ensemble & Combined outputs & Improve robustness & Stable overall performance \\
\end{xltabular}}
\vspace{0.3em}\noindent\textit{Note.} EEG = Electroencephalography. See chapter prose for full context.

\subsection{EEG System Integration}\label{sec:appendix_ch3_eeg_system_integration}

\noindent\textbf{EEG Framework System Integration Components.} Table~\ref{tab:eeg_system_integration} below presents the detailed analysis.

{\tablestripes\tablefontsize

\noindent Table~\ref{tab:eeg_system_integration} presents the eeg framework system integration components.

\needspace{20\baselineskip}
\begin{xltabular}{\textwidth}{@{} >{\raggedright\arraybackslash}p{2.8cm} L L @{}}
\caption[EEG Framework System Integration Components]{EEG Framework System Integration Components}\label{tab:eeg_system_integration} \\
\toprule
\thead{Integration Area} & \thead{What It Does} & \thead{Value} \\
\midrule
\endfirsthead
\endhead
\bottomrule
\endlastfoot
Patient mobile app & Onboarding, symptom forms, report access, monitoring & Patient-side continuity \\
Parent/caregiver app & Behavioural input, observation logging, alerts & Primary-data capture \\
Hospital system & EHR/HIS integration, clinician dashboard, report storage & Workflow compatibility \\
MCP layer & Coordinates model, retrieval, records, tools, notifications & System orchestration \\
SOS/emergency & Trigger escalation when threshold exceeded & Safety and response continuity \\
\end{xltabular}}
\vspace{0.3em}\noindent\textit{Note.} EEG = Electroencephalography. See chapter prose for full context.

\subsection{Questionnaire Variable Mapping}\label{sec:appendix_ch3_questionnaire_variable_mapping}

\noindent Table~\ref{tab:questionnaire_variable_mapping} maps primary-Data Instrument: Questionnaire Variable Mapping.


{\tablestripes\tablefontsize
\needspace{20\baselineskip}
\begin{xltabular}
{\textwidth}{@{} >{\raggedright\arraybackslash}p{1.8cm} >{\raggedright\arraybackslash}p{1.1cm} >{\raggedright\arraybackslash}p{2.0cm} >{\raggedright\arraybackslash}p{2.3cm} >{\raggedright\arraybackslash}p{1.1cm} L @{}}
\caption[Primary-Data Instrument]{Primary-Data Instrument: Questionnaire Variable Mapping} \label{tab:questionnaire_variable_mapping} \\
\toprule
\thead{Section} & \thead{Resp.} & \thead{Focus} & \thead{Variable} & \thead{Type} & \thead{Analysis} \\
\midrule
\endfirsthead
\endhead
\bottomrule
\endlastfoot
Demographics & Par./Pat. & Age, gender, clin.\ history & demographic\_\allowbreak vars & Nom./ord. & Descriptive, subgrp \\
Social comm.\ (ASD) & Parent & Eye contact, reciprocity, peers & asd\_\allowbreak social\_\allowbreak score & Likert & Descriptive, reliab. \\
Repet.\ behav.\ (ASD) & Parent & Routines, repetitive acts, fixation & asd\_\allowbreak repetitive\_\allowbreak score & Likert & Descriptive, subgrp \\
Sensory (ASD) & Par./Pat. & Noise, touch, visual overload & asd\_\allowbreak sensory\_\allowbreak score & Likert & Severity scoring \\
Emot.\ reg.\ (ASD) & Par./Pat. & Meltdown, frustration, recovery & asd\_\allowbreak emotion\_\allowbreak score & Likert & Corr.\ w/ trust \\
% [ASD-ONLY] Motor  & Patient & Tremor, rigidity, gait, balance & pd\_\allowbreak motor\_\allowbreak score & Likert & Severity analysis \\
% [ASD-ONLY] Non-motor  & Patient & Fatigue, sleep, constipation & pd\_\allowbreak nonmotor\_\allowbreak score & Likert & Subgroup analysis \\
% [ASD-ONLY] Cognitive  & Patient & Memory, concentration, planning & pd\_\allowbreak cognitive\_\allowbreak score & Likert & Contextual support \\
% [ASD-ONLY] Emotional  & Patient & Anxiety, depression, apathy & pd\_\allowbreak emotional\_\allowbreak score & Likert & Corr.\ w/ usability \\
Trust in AI & All & Trust, clinical reasonableness & trust\_\allowbreak score & Likert & Mean, SD, $\alpha$ \\
Expl.\ clarity & All & Understandable, useful, clear & clarity\_\allowbreak score & Likert & Corr.\ w/ trust \\
Usability & Par./Pat. & Chatbot ease, question clarity & usability\_\allowbreak score & Likert & Onboarding assess. \\
Adoption & All & Would use again, recommend & adoption\_\allowbreak score & Likert & Threshold testing \\
Privacy & Par./Pat. & Data sharing comfort, privacy & privacy\_\allowbreak score & Likert & Governance analysis \\
Open-ended & Par./Pat. & Concerns, challenges, feedback & coded\_\allowbreak themes & Qual. & Thematic coding \\
\end{xltabular}
}
\vspace{0.3em}\noindent\textit{Note.} ASD = Autism Spectrum Disorder. See chapter prose for full context.
}%

\subsection{Primary Data Pipeline}\label{sec:appendix_ch3_primary_data_pipeline}

\noindent\textbf{Primary-Data Quantitative Analysis Pipeline.} Table~\ref{tab:primary_data_pipeline} below presents the detailed analysis.

{\tablestripes\tablefontsize

\noindent Table~\ref{tab:primary_data_pipeline} presents the primary-data quantitative analysis pipeline.

\needspace{20\baselineskip}
\begin{xltabular}{\textwidth}{@{} >{\raggedright\arraybackslash}p{1.2cm} >{\raggedright\arraybackslash}p{4cm} L @{}}
\caption[Primary-Data Quantitative Analysis Pipeline]{Primary-Data Quantitative Analysis Pipeline}\label{tab:primary_data_pipeline} \\
\toprule
\thead{Step} & \thead{Action} & \thead{Output} \\
\midrule
\endfirsthead
\endhead
\bottomrule
\endlastfoot
1 & Convert structured responses to numeric form & Analysis-ready dataset \\
2 & Build subscale scores & ASD behaviour score, trust score \\
3 & Test reliability & Cronbach's $\alpha$ / consistency evidence \\
4 & Summarise scores & Means, SD, frequencies, severity categories \\
5 & Compare groups or conditions & Differences across ASD subgroups or response groups \\
6 & Correlate with EEG/model outputs & Association between questionnaire burden and model evidence \\
7 & Evaluate threshold achievement & Trust $\geq$ threshold, usability $\geq$ threshold \\
8 & Integrate with qualitative findings & Joint interpretation via mixed-method display \\
\end{xltabular}}
\vspace{0.3em}\noindent\textit{Note.} ASD = Autism Spectrum Disorder; EEG = Electroencephalography. See chapter prose for full context.

\subsection{EEG Pipeline Diagram}\label{sec:appendix_ch3_eeg_pipeline_clean}
\needspace{24\baselineskip}
\begin{figure}[!htbp]
\centering
\resizebox{0.97\textwidth}{!}{%
\begin{tikzpicture}[
  stepbox/.style={
    rectangle,
    rounded corners=4pt,
    draw=ggunavy,
    fill=ggunavy!6,
    minimum width=2.8cm,
    minimum height=1.1cm,
    text width=2.4cm,
    align=center,
    font=\small
  },
  arr/.style={-{Stealth[length=4mm,width=3mm]}, thick}
]

\node[stepbox] (s1) at (0,0) {Acquire EEG};
\node[stepbox] (s2) at (3.2,0) {Remove Artefacts};
\node[stepbox] (s3) at (6.4,0) {Filter Signal};
\node[stepbox] (s4) at (9.6,0) {Segment Data};

\node[stepbox] (s5) at (1.6,-2.2) {Transform Signal};
\node[stepbox] (s6) at (4.8,-2.2) {Extract Features};
\node[stepbox] (s7) at (8.0,-2.2) {Normalise Input};
\node[stepbox] (s8) at (11.2,-2.2) {Run Inference};

\node[stepbox] (s9) at (4.0,-4.4) {Estimate Confidence};
\node[stepbox] (s10) at (8.0,-4.4) {Generate Explanation};

\draw[arr] (s1) -- (s2);
\draw[arr] (s2) -- (s3);
\draw[arr] (s3) -- (s4);

\draw[arr] (s4) -- (s5);
\draw[arr] (s5) -- (s6);
\draw[arr] (s6) -- (s7);
\draw[arr] (s7) -- (s8);

\draw[arr] (s8) -- (s9);
\draw[arr] (s9) -- (s10);

\end{tikzpicture}%
}
\caption{Secondary EEG Processing Pipeline}
\label{fig:eeg_pipeline_clean}
\vspace{0.5em}
\noindent\textit{Note.} Signal transformation includes FFT, PSD/SPDD, and SWT or wavelet-based analysis depending on domain requirements.
\end{figure}

\subsection{Analytical Pipeline Diagram}\label{sec:appendix_ch3_analytical_pipeline}
\needspace{24\baselineskip}
\begin{figure}[!htbp]
    \centering
    \resizebox{0.97\textwidth}{!}{%
    \begin{tikzpicture}[
        node distance=0.8cm,
        stagebox/.style={rectangle, rounded corners=4pt, draw=ggunavy,
            minimum width=8cm, minimum height=1.0cm, font=\figtitlefont,
            align=center, text=ggunavy},
        stagelabel/.style={circle, draw=ggugold!80!black, fill=ggugold!25,
            minimum size=0.7cm, font=\figtitlefont, text=ggunavy, inner sep=0pt},
        annot/.style={font=\fignodefont, text=ggunavy!80, align=left, anchor=west},
        arrow/.style={-{Stealth[length=5pt]}, thick, ggunavy}
    ]
        
\node[stagebox, fill=lightblue] (s1) {Stage 1: Data Ingestion};
        
\node[stagelabel, left=0.6cm of s1] (l1) {1};
        
\node[annot, right=0.6cm of s1] (a1) {KAU ASD dataset (KAU, PhysioNet,\\DEAP, SAM40, BCI-IV, PBS)};

        
\node[stagebox, fill=lightorange, below=of s1] (s2) {Stage 2: Preprocessing};
        
\node[stagelabel, left=0.6cm of s2] (l2) {2};
        
\node[annot, right=0.6cm of s2] (a2) {Filtering, ICA artifact removal,\
ormalization, epoching};

        
\node[stagebox, fill=lightteal, below=of s2] (s3) {Stage 3: Feature Engineering};
        
\node[stagelabel, left=0.6cm of s3] (l3) {3};
        
\node[annot, right=0.6cm of s3] (a3) {PSD, CWT, CSP, entropy,\\connectivity, GAN features};

        
\node[stagebox, fill=lightpurple, below=of s3] (s4) {Stage 4: Model Training \& Evaluation};
        
\node[stagelabel, left=0.6cm of s4] (l4) {4};
        
\node[annot, right=0.6cm of s4] (a4) {DL ensemble (VotingClassifier),\\baselines, DL compared};

        
\node[stagebox, fill=lightgreen, below=of s4] (s5) {Stage 5: Interpretation \& Governance};
        
\node[stagelabel, left=0.6cm of s5] (l5) {5};
        
\node[annot, right=0.6cm of s5] (a5) {SHAP feature importance,\\RAG explanations, audit trail};

        \draw[arrow] (s1) -- (s2);
        \draw[arrow] (s2) -- (s3);
        \draw[arrow] (s3) -- (s4);
        \draw[arrow] (s4) -- (s5);

        \draw[decorate, decoration={brace, amplitude=8pt, mirror}, thick, ggugold!80!black]
            ([xshift=-1.8cm]s1.north west) -- ([xshift=-1.8cm]s5.south west)
            node[midway, left=10pt, font=\figtitlefont, text=ggugold!80!black, align=center, rotate=90]
            {Governance Checkpoints};
        %% Extend bounding box to include left brace and label
        \path ([xshift=-3.2cm]s1.west) ++(0,0);
    \end{tikzpicture}%
    }% end resizebox
    \caption[Five-Stage Analytical Pipeline]{Five-Stage Analytical Pipeline for ASD Diagnosis}
    \label{fig:analytical_pipeline}
    \vspace{0.5em}
\noindent\textit{Note.} PSD = power spectral density; CWT = continuous wavelet transform; CSP = common spatial pattern; ICA = independent component analysis; DL = deep learning; SHAP = SHapley Additive exPlanations; RAG = retrieval-augmented generation.
\end{figure}

\subsection{C4 Container Architecture}\label{sec:appendix_ch3_c4_container}
\needspace{24\baselineskip}
\begin{figure}[!htbp]
\centering
\resizebox{0.97\textwidth}{!}{%
\begin{tikzpicture}[
    node distance=0.8cm and 1.2cm,
    container/.style={rectangle, rounded corners, draw=ch3header, fill=ch3light, font=\small, text width=2.5cm, align=center, minimum height=1.2cm},
    external/.style={rectangle, rounded corners, draw=ch3header!50, fill=white, font=\small, text width=2.5cm, align=center, minimum height=1cm, dashed},
    arrow/.style={-{Stealth[length=4mm, width=3mm]}, very thick, ggunavy},
    label/.style={font=\scriptsize\itshape, ch3header}
]
% External actors
\node[external] (patient) {Patient /\\Parent};
\node[external, right=4cm of patient] (clinician) {Clinician};
\node[external, right=4cm of clinician] (ehr) {Hospital\\EHR / HIS};

% Containers row 1
\node[container, below=1.5cm of patient] (app) {\textbf{Mobile App}\\Onboarding,\\Consent, Forms};
\node[container, right=of app] (device) {\textbf{Emotiv/IoT}\\EEG Capture};
\node[container, right=of device] (gateway) {\textbf{Gateway}\\API, Auth,\\Ingestion};
\node[container, right=of gateway] (backend) {\textbf{Backend}\\Preprocessing,\\Features, Model};

% Containers row 2
\node[container, below=1.5cm of app, xshift=1.5cm] (rag) {\textbf{RAG Engine}\\Retrieval,\\Explanation};
\node[container, right=of rag] (mcp) {\textbf{MCP / Orchestrator}\\Tool Coordination,\\Logging};
\node[container, right=of mcp] (dashboard) {\textbf{Dashboard}\\Clinician View,\\Review Queue};
\node[container, right=of dashboard] (monitor) {\textbf{Monitoring}\\Drift, KPIs,\\Alerts};

% Arrows
\draw[arrow] (patient) -- (app);
\draw[arrow] (app) -- (gateway);
\draw[arrow] (device) -- (gateway);
\draw[arrow] (gateway) -- (backend);
\draw[arrow] (backend) -- (mcp);
\draw[arrow] (mcp) -- (rag);
\draw[arrow] (mcp) -- (dashboard);
\draw[arrow] (mcp) -- (monitor);
\draw[arrow] (dashboard) -- (clinician);
\draw[arrow] (dashboard) -- (ehr);

% Labels
\node[label, above=0.1cm of app] {Patient-Side};
\node[label, above=0.1cm of backend] {Server-Side};
\node[label, above=0.1cm of dashboard] {Clinical-Side};

\end{tikzpicture}%
}
\caption[C4 Container Diagram: Mobile App, Gateway, Backend, RAG,]{C4 Container Diagram: Mobile App, Gateway, Backend, RAG, MCP, and Hospital Integration}
\label{fig:c4_container}
\end{figure}

\subsection{End-to-End Architecture Table}\label{sec:appendix_ch3_e2e_architecture}
\noindent Table~\ref{tab:e2e_architecture} specifies the complete 15-layer pipeline from patient intake through AI classification, RAG explanation, and clinical decision output.
\clearpage
{%
\tablefontsize
\begin{xltabular}
{\textwidth}{@{} r >{\raggedright\arraybackslash}p{2cm} L L L @{}}
\caption[End-to-End System Architecture: 15-Layer Pipeline from]{End-to-End System Architecture: 15-Layer Pipeline from Intake to Decision} \label{tab:e2e_architecture} \\
\toprule
\thead{\#} & \thead{Layer} & \thead{What Happens} & \thead{Key Controls} & \thead{Output} \\
\midrule
\endfirsthead
\endhead
\bottomrule
\endlastfoot
1 & Patient onboarding & Mobile/web intake, consent, symptom/behavioural capture & Identity, consent, RBAC, minimum data & Registered case \\
2 & Device session & Emotiv/IoT EEG pairs to app or gateway & Device auth, session token, signal-quality gate & Active recording session \\
3 & Data capture & EEG + forms/symptoms/behavioural inputs collected & Encrypted transport, PII separation & Raw EEG + primary data \\
4 & Gateway / ingestion & Backend receives stream/files & TLS, API auth, audit logging, schema validation & Accepted intake payload \\
5 & Preprocessing & Artefact handling, band-pass filter, epoching, normalisation & Reproducible rules, quality thresholds & Clean EEG segments \\
6 & Transformation & FFT, PSD/SPDD, SWT, feature extraction, 1D-to-2D & Versioned pipeline, leakage-safe processing & Model-ready inputs \\
7 & Model inference & ML/DL/CV model scores class/risk/confidence & Calibrated outputs, uncertainty handling, monitored version & Prediction + confidence \\
8 & Explainability & SHAP/saliency/attention computed & Method aligned to model type & Machine-level explanation \\
9 & RAG layer & Retrieve approved references, generate summaries & Controlled corpus, citation traceability & Clinician/patient explanations \\
10 & Orchestration & MCP coordinates model, retrieval, alerts, logs & Tool boundaries, access checks, traceability & Orchestrated workflow \\
11 & Decision support & Thresholding, review queue, escalation logic & Abstain/escalate rules & Review/monitor/escalate \\
12 & Hospital integration & Dashboard/EHR/HIS integration & Mapped identifiers, access control, audit trail & Clinician case summary \\
13 & Safety escalation & High-risk triggers SOS/escalation & Documented criteria, override path, notification log & Urgent review path \\
14 & Monitoring & Drift, latency, errors, trust/use metrics & KPI dashboard, alerts, version tracking & Operational governance \\
15 & Decision & Organisation decides evidence-maturity interpretation & Technical + trust + governance + illustrative financial scenario evidence & Evaluation recommendation \\
\end{xltabular}
\vspace{0.3em}\noindent\textit{Note.} SHAP = SHapley Additive exPlanations; RAG = Retrieval-Augmented Generation; EEG = Electroencephalography; KPI = Key Performance Indicator. See chapter prose for full context.
}%

\subsection{B2B Hospital Architecture}\label{sec:appendix_ch3_b2b_hospital_architecture}
\noindent Table~\ref{tab:b2b_hospital_architecture} presents north America B2B Hospital Architecture: Stages and Requirements.

{\tablestripes\tablefontsize
\needspace{20\baselineskip}
\begin{xltabular}{\textwidth}{@{} r L L @{}}
\caption[North America B2B Hospital Architecture]{North America B2B Hospital Architecture: Stages and Requirements}\label{tab:b2b_hospital_architecture} \\
\toprule
\thead{\#} & \thead{Hospital Architecture Stage} & \thead{Why This Is the Strongest North America Fit} \\
\midrule
\endfirsthead
\endhead
\bottomrule
\endlastfoot
1 & Patient onboarded by hospital staff through secure web intake & Reduces consumer misuse risk \\
2 & Emotiv EEG captured in supervised clinical workflow & Improves data quality and chain of custody \\
3 & Gateway sends data to hospital-controlled secure backend & Supports privacy and vendor governance \\
4 & Backend runs preprocessing, FFT/PSD/SWT, feature extraction, model inference & Keeps technical pipeline controlled \\
5 & Calibrated output with uncertainty and review logic & Safer for CDS-style use \\
6 & SHAP + RAG generates clinician-facing explanation & Improves interpretability \\
7 & Result goes only to clinician dashboard/EHR, not direct patient diagnosis & Aligns with safer CDS evaluation \\
8 & Clinician reviews, accepts, overrides, or escalates & Preserves human-in-the-loop responsibility \\
9 & Audit trail, access logs, model/version traceability maintained & Supports HIPAA/provincial compliance \\
10 & Monitoring tracks drift, latency, override rate, false negatives & Supports pilot governance \\
\end{xltabular}}
\vspace{2pt}
\noindent\textit{Note.} CDS = Clinical Decision Support; EHR = Electronic Health Record; FFT = Fast Fourier Transform; PSD = Power Spectral Density; SWT = Stationary Wavelet Transform; SHAP = SHapley Additive exPlanations; RAG = Retrieval-Augmented Generation.

\subsection{Thesis Objective Alignment Detail}\label{sec:appendix_ch3_thesis_objective_alignment}
\needspace{24\baselineskip}
\begingroup\tablefontsize
\noindent Table~\ref{tab:thesis_objective_alignment} presents thesis Objective Alignment Across Chapters.


\begin{longtable}
{@{} >{\raggedright\arraybackslash}p{3.0cm} >{\raggedright\arraybackslash}p{1.7cm} >{\raggedright\arraybackslash}p{1.8cm} >{\raggedright\arraybackslash}p{2.0cm} >{\raggedright\arraybackslash}p{1.8cm} >{\raggedright\arraybackslash}p{2.1cm} >{\raggedright\arraybackslash}p{2.6cm} @{}}
\caption{Thesis Objective Alignment Across Chapters}
\label{tab:thesis_objective_alignment} \\

\toprule
\thead{Objective Component} & \thead{Ch.~1} & \thead{Ch.~2} & \thead{Ch.~3} & \thead{Ch.~4} & \thead{Ch.~5} & \thead{Overall} \\
\midrule
\endfirsthead

\toprule
\thead{Objective Component} & \thead{Ch.~1} & \thead{Ch.~2} & \thead{Ch.~3} & \thead{Ch.~4} & \thead{Ch.~5} & \thead{Overall} \\
\midrule
\endhead

\bottomrule
\endfoot

Define biomedical AI problem clearly & Strong & Supportive & Indirect & Indirect & Indirect & Strong \\
Justify why literature and practice are insufficient & Moderate & Strong & Supportive & Indirect & Indirect & Strong \\
Develop unified ASD-focused framework & Framing only & Justification only & Strong & Tested empirically & Interpreted strategically & Strong \\
Integrate primary and secondary data & Previewed & Justified as gap & Strongly operationalised & Partially evidenced & Interpreted in evaluation & Strong \\
Evaluate technical EEG performance & Previewed early & Justified & Strongly designed & Strongly evidenced & Interpreted & Strong \\
Evaluate explainability and trust & Framed strongly & Justified strongly & Methodologically supported & Evidenced & Interpreted for adoption & Strong \\
Evaluate governance and safe evaluation & Framed strongly & Justified conceptually & Strongly designed & Partly evidenced & Heavily interpreted & Moderate--strong \\
Evaluate workflow and operational value & Framed strongly & Justified moderately & Methodologically included & Partly evidenced & Strongly interpreted & Moderate \\
Support B2B hospital decision making & Framed strongly & Justified moderately & Operationalised & Partly evidenced & Strongly interpreted & Strong \\
Bound B2C as extension or future pathway & Over-expanded & Moderately justified & Included, broad & Limited evidence & Interpreted, broad & Moderate \\
Produce evidence-maturity interpretation decision logic & Previewed early & Not central & Designed & Partially evidenced & Main destination & Moderate--strong \\

\end{longtable}
\endgroup

\subsection{Chapter vs Thesis Matrix}\label{sec:appendix_ch3_chapter_vs_thesis_objective_matrix}
\needspace{24\baselineskip}
\begingroup\tablefontsize
\noindent Table~\ref{tab:chapter_vs_thesis_objective_matrix} documents chapter Objective versus Thesis Objective Matrix.


\begin{longtable}
{@{} >{\raggedright\arraybackslash}p{1.4cm} >{\raggedright\arraybackslash}p{3.6cm} >{\raggedright\arraybackslash}p{3.6cm} >{\raggedright\arraybackslash}p{1.8cm} >{\raggedright\arraybackslash}p{3.2cm} @{}}
\caption{Chapter Objective versus Thesis Objective Matrix}
\label{tab:chapter_vs_thesis_objective_matrix} \\

\toprule
\thead{Chapter} & \thead{Chapter Objective} & \thead{Link to Thesis Objective} & \thead{Strength of Fit} & \thead{Main Problem} \\
\midrule
\endfirsthead

\toprule
\thead{Chapter} & \thead{Chapter Objective} & \thead{Link to Thesis Objective} & \thead{Strength of Fit} & \thead{Main Problem} \\
\midrule
\endhead

\bottomrule
\endfoot

Ch.\ 1 & Introduce problem, gap, objectives, hypotheses, scope & Defines why thesis exists & Moderate--strong & Too much method/architecture in introduction \\
Ch.\ 2 & Synthesise literature and identify gaps & Justifies why thesis is needed & Strong & Synthesis weakened by duplication \\
Ch.\ 3 & Explain framework and study design & Operationalises how thesis objective is tested & Strong & Excess architecture/protocol detail \\
Ch.\ 4 & Present primary, secondary, and integrated findings & Shows empirical support for thesis objective & Very strong & Overloaded with tables and displays \\
Ch.\ 5 & Interpret findings and make final decisions & Explains what thesis objective means in practice & Moderate--strong & Final decision diluted by adjacent frameworks \\

\end{longtable}
\endgroup

\subsection{Chapter SWOT Action Plan}\label{sec:appendix_ch3_chapter_strengths_weaknesses_actionplan}
\needspace{24\baselineskip}
\begingroup\tablefontsize
\noindent Table~\ref{tab:chapter_strengths_weaknesses_actionplan} presents chapter-Wise Strengths, Weaknesses, and Action Plan.


\begin{longtable}
{@{} >{\raggedright\arraybackslash}p{1.4cm} >{\raggedright\arraybackslash}p{3.8cm} >{\raggedright\arraybackslash}p{4cm} >{\raggedright\arraybackslash}p{4cm} @{}}
\caption{Chapter-Wise Strengths, Weaknesses, and Action Plan}
\label{tab:chapter_strengths_weaknesses_actionplan} \\

\toprule
\thead{Chapter} & \thead{Strengths} & \thead{Weaknesses} & \thead{Action Plan} \\
\midrule
\endfirsthead

\toprule
\thead{Chapter} & \thead{Strengths} & \thead{Weaknesses} & \thead{Action Plan} \\
\midrule
\endhead

\bottomrule
\endfoot

Ch.\ 1 & Strong problem framing, stakeholder relevance, ambitious scope & Overloaded; too many workflows and system details; methods creep & Cut operational detail, move method/system content to Ch.\ 3 \\
Ch.\ 2 & Broad evidence base, meaningful gaps, strong relevance to thesis rationale & Duplication, inventory-style writing, too many side frameworks & Tighten synthesis, remove duplicates, focus on gap-to-study logic \\
Ch.\ 3 & Strong methodological ambition, integrated design, good technical/governance coverage & Bloated, too many master tables, architecture/protocol overload & One method flow per strand, move protocols to appendix \\
Ch.\ 4 & Strongest empirical chapter, explicit hypothesis testing, strong evidence & Too many support tables, repeated displays, discussion leakage & One main table + one figure per hypothesis, move helpers to appendix \\
Ch.\ 5 & Rich interpretation, strong DBA orientation, practical relevance & Too many matrices, policy sprawl, final decision not sharp & Compress around domain readiness, trust, ROI, evidence-maturity interpretation \\

\end{longtable}
\endgroup

\subsection{Chapter Structure (APA)}\label{sec:appendix_ch3_chapter_structure_apa}
\needspace{24\baselineskip}
\begingroup\tablefontsize
\noindent Table~\ref{tab:chapter_structure_apa} presents chapter-Wise Objective, Process, and Decision Structure.


\begin{longtable}
{@{} >{\raggedright\arraybackslash}p{0.7cm} >{\raggedright\arraybackslash}p{1.8cm} >{\raggedright\arraybackslash}p{2cm} >{\raggedright\arraybackslash}p{2.8cm} >{\raggedright\arraybackslash}p{2.2cm} >{\raggedright\arraybackslash}p{2.2cm} >{\raggedright\arraybackslash}p{2.2cm} @{}}
\caption{Chapter-Wise Objective, Process, and Decision Structure}
\label{tab:chapter_structure_apa} \\

\toprule
\thead{Ch.} & \thead{Objective} & \thead{Goal} & \thead{Task} & \thead{Input} & \thead{Output} & \thead{Decision Role} \\
\midrule
\endfirsthead

\toprule
\thead{Ch.} & \thead{Objective} & \thead{Goal} & \thead{Task} & \thead{Input} & \thead{Output} & \thead{Decision Role} \\
\midrule
\endhead

\bottomrule
\endfoot

1 & Introduce study & Establish need and scope & Frame problem, gap, objectives, hypotheses, scope & Research problem, context, literature & Problem statement, study direction & Justifies what study should address \\
2 & Review literature & Justify study intellectually & Compare studies, identify gaps, build conceptual logic & Published studies, theories, benchmarks & Gap-driven literature synthesis & Whether study is justified \\
3 & Explain methods & Show how study tests claims & Define data, variables, preprocessing, modelling, validation, integration & Primary instruments, EEG data, method choices & Reproducible study design & Whether claims can be tested credibly \\
4 & Present findings & Show what evidence demonstrates & Report results, test hypotheses, integrate findings & Primary results, EEG outputs, robustness evidence & Empirical findings and verdicts & Whether framework is supported \\
5 & Interpret, conclude & Translate evidence into action & Assess readiness, implications, final recommendation & Ch.\ 4 findings, governance/adoption logic & Interpretation, evidence-maturity interpretation & What should be done next \\

\end{longtable}
\endgroup

\subsection{Five-Chapter Logic Diagram}\label{sec:appendix_ch3_five_chapter_logic}
\needspace{24\baselineskip}
\begin{figure}[!htbp]
\centering
\resizebox{0.97\textwidth}{!}{%
\begin{tikzpicture}[
  thesisstep/.style={
    rectangle,
    rounded corners=4pt,
    draw=ggunavy!70,
    fill=ggunavy!6,
    minimum width=3.0cm,
    minimum height=1.1cm,
    text width=2.8cm,
    align=center,
    font=\small
  },
  thesisarrow/.style={
    -{Stealth[length=4mm,width=3mm]},
    thick
  }
]

\node[thesisstep] (c1) at (0,0) {Chapter 1\\Frame Problem};
\node[thesisstep] (c2) at (4,0) {Chapter 2\\Justify Study};
\node[thesisstep] (c3) at (8,0) {Chapter 3\\Design Method};
\node[thesisstep] (c4) at (12,0) {Chapter 4\\Present Evidence};
\node[thesisstep] (c5) at (16,0) {Chapter 5\\Make Decision};

\draw[thesisarrow] (c1) -- (c2);
\draw[thesisarrow] (c2) -- (c3);
\draw[thesisarrow] (c3) -- (c4);
\draw[thesisarrow] (c4) -- (c5);

\end{tikzpicture}%
}
\caption{Five-Chapter Thesis Logic}
\label{fig:five_chapter_logic}
\vspace{0.5em}
\noindent\textit{Note.} The thesis moves from problem framing to justification, methodological design, empirical evidence, and final decision logic.
\end{figure}

\subsection{Detailed SMART Matrix}\label{sec:appendix_ch3_detailed_chapter_smart_matrix}

\noindent\textbf{SMART Assessment by Chapter.} Table~\ref{tab:detailed_chapter_smart_matrix} evaluates each chapter against the SMART criteria (Specific, Measurable, Achievable, Relevant, Time-bound), documenting the evidence produced, reliability level, and decision contribution.

{\tablestripes\tablefontsize

\noindent Table~\ref{tab:detailed_chapter_smart_matrix} presents the detailed chapter-wise smart, evidence, reliability, and.

\needspace{20\baselineskip}
\begin{xltabular}{\textwidth}{@{} >{\centering
\tablefontsize\arraybackslash}p{0.5cm} L >{\centering\arraybackslash}p{0.6cm} >{\centering\arraybackslash}p{0.6cm} >{\centering\arraybackslash}p{0.6cm} >{\centering\arraybackslash}p{0.6cm} >{\centering\arraybackslash}p{0.6cm} L L L @{}}
\caption[Detailed Chapter-Wise SMART, Evidence, Reliability, and]{Detailed Chapter-Wise SMART, Evidence, Reliability, and Decision Matrix. Each chapter is assessed against SMART criteria; the evidence and reliability columns confirm that every chapter contributes testable, measurable outputs to the evaluation decision.}
\label{tab:detailed_chapter_smart_matrix} \\

\toprule
\thead{Ch} & \thead{Objective} & \thead{S} & \thead{M} & \thead{A} & \thead{R} & \thead{T} & \thead{Evidence} & \thead{Reliability} & \thead{Decision} \\
\midrule
\endfirsthead
\endhead

\bottomrule
\endfoot

1 & Define problem, gap, objectives, hypotheses, scope & Yes & Problem clarity, alignment & Yes, if bounded & Yes & Yes & Problem, gap, stakeholders & Moderate--strong & What thesis addresses \\
2 & Synthesise literature, identify gaps & Yes & Theme coverage, gap clarity & Yes, if synthesis-focused & Yes & Yes & Review, gap matrix, theory & Strong if selective & Whether study justified \\
3 & Design full methodological framework & Yes & Variables, metrics, validation & Yes, if protocol controlled & Yes & Yes & Method tables, pipeline, KPIs & Strong if reproducible & Whether claims testable \\
4 & Present findings, test hypotheses & Yes & Accuracy, AUC, trust, robustness & Yes, if selective displays & Yes & Yes & Performance, SHAP, integration & Very strong if clean & Whether framework supported \\
5 & Interpret, conclude & Yes & Yes & Yes & Yes & Yes & Decision framework & Strong & Adopt/pilot/defer \\
\bottomrule
\end{xltabular}
}
\vspace{0.3em}\noindent\textit{Note.} SHAP = SHapley Additive exPlanations; AUC = Area Under (ROC) Curve. See chapter prose for full context.

\subsection{Chapter Decision Role Summary}\label{sec:appendix_ch3_chapter_decision_role_summary}

\noindent\textbf{Chapter-Wise Decision Role Summary.} Table~\ref{tab:chapter_decision_role_summary} below presents the detailed analysis.

{\tablestripes\tablefontsize

\noindent Table~\ref{tab:chapter_decision_role_summary} presents the chapter-wise decision role summary.

\needspace{20\baselineskip}
\begin{xltabular}{\textwidth}{@{} >{\raggedright\arraybackslash}p{1.5cm} L L @{}}
\caption[Chapter-Wise Decision Role Summary]{Chapter-Wise Decision Role Summary}\label{tab:chapter_decision_role_summary} \\
\toprule
\thead{Chapter} & \thead{Core Function} & \thead{Decision Role} \\
\midrule
\endfirsthead
\endhead
\bottomrule
\endlastfoot
1 & Frame the problem and define scope & Decides what the thesis must address \\
2 & Justify the study through literature and gaps & Decides why the study is needed \\
3 & Design the method and evaluation logic & Decides how the claims will be tested \\
4 & Report evidence and hypothesis results & Decides whether the claims are supported \\
5 & Interpret results and set evaluation stance & Decides what should be done in practice \\
\end{xltabular}}
\vspace{0.3em}\noindent\textit{Note.} Source: dissertation analysis; abbreviations defined in chapter prose.

\subsection{Chapter Decision Flow Diagram}\label{sec:appendix_ch3_chapter_decision_flow}
\needspace{24\baselineskip}
\begin{figure}[!htbp]
\centering
\resizebox{0.97\textwidth}{!}{%
\begin{tikzpicture}[
  thesisstep/.style={
    rectangle,
    rounded corners=4pt,
    draw=ggunavy!70,
    fill=ggunavy!6,
    minimum width=3.0cm,
    minimum height=1.0cm,
    text width=2.8cm,
    align=center,
    font=\small
  },
  thesisarrow/.style={-{Stealth[length=4mm,width=3mm]}, thick}
]
\node[thesisstep] (a) at (0,0) {Chapter 1\\Problem and Scope};
\node[thesisstep] (b) at (4,0) {Chapter 2\\Literature and Gaps};
\node[thesisstep] (c) at (8,0) {Chapter 3\\Method and Evaluation};
\node[thesisstep] (d) at (12,0) {Chapter 4\\Evidence and Verdicts};
\node[thesisstep] (e) at (16,0) {Chapter 5\\Interpretation and Decision};

\draw[thesisarrow] (a) -- (b);
\draw[thesisarrow] (b) -- (c);
\draw[thesisarrow] (c) -- (d);
\draw[thesisarrow] (d) -- (e);
\end{tikzpicture}%
}
\caption{Chapter-Wise Thesis Decision Flow}
\label{fig:chapter_decision_flow}
\vspace{0.5em}
\noindent\textit{Note.} The thesis progresses from problem framing to literature justification, methodological design, empirical evidence, and final decision logic.
\end{figure}


%% =============================================================================
%% APPENDIX C — EXTENDED TECHNICAL SUPPLEMENT (2026-05-12)
%% Subsections moved out of Chapter 3 per scope-discipline:
%%   C-S1. Signal Standardization Protocol (Part B technical detail)
%%   C-S2. PBS Image Preprocessing Pipeline (Part B technical detail)
%%   C-S3. Deep Learning Architectures (Part B technical detail)
%% Chapter 3 retains one-paragraph pointers to each section here.
%% =============================================================================

\addcontentsline{toc}{section}{Extended Technical Supplement} % [TOC-appendix-section-insertion]
\section*{Extended Technical Supplement}
\addcontentsline{toc}{section}{Extended Technical Supplement}
\label{sec:appc_extended_supplement}

\noindent The following three subsections were relocated from Chapter~3
(Research Methods) to preserve Chapter~3's role as methodology \emph{argument}
rather than methodology \emph{specification}. Each is a technical preprocessing
or architectural detail that supports the empirical evaluation methodology in
Chapter~3~Part~B but is not load-bearing for the methodology argument itself.
The cross-reference labels (\texttt{tab:signal\_\allowbreak filter\_\allowbreak params}, \texttt{fig:pbs\_\allowbreak pipeline},
etc.) are retained unchanged so Chapter~3 and downstream references resolve to
these locations.


%% ======================================================================
%% C-S1. Signal Standardization Protocol
%% Moved from Chapter 3 — 2026-05-12
%% ======================================================================
\addcontentsline{toc}{section}{C-S1. Signal Standardization Protocol} % [TOC-appendix-section-insertion]
\section*{C-S1. Signal Standardization Protocol}
\addcontentsline{toc}{section}{C-S1. Signal Standardization Protocol}
\label{sec:appc_signal_std}

\iffalse %% Engineering detail moved to Appendix C --- sec1_signal_standardisation
%% SUPPRESSED FOR WORD COUNT
Table~\ref{tab:eeg_bands} documents the five canonical EEG frequency bands, their clinical significance, and their relevance to the five clinical domains in this study.


\needspace{24\baselineskip}
\iffalse % AUTO-SUPPRESS non-ASD
\noindent\textbf{EEG Frequency Bands: Neurophysiological Basis and.}

\needspace{24\baselineskip}
\begin{table}[H]
\tablestripes
\caption[EEG Frequency Bands]{EEG Frequency Bands: Neurophysiological Basis and Clinical Domain Relevance}
\tablefontsize
\begin{tabularx}{\textwidth}{@{} >{\raggedright\arraybackslash}p{2cm} >{\raggedright\arraybackslash}p{1.8cm} >{\raggedright\arraybackslash}p{2.5cm} L L @{}}
\toprule
\thead{Band} & \thead{Range} & \thead{Brain State} & \thead{Clinical Significance} & \thead{Domain Relevance} \\
\midrule
Delta ($\delta$) & 0.5--4 Hz & Deep sleep & Brain damage; encephalopathy & Schizophrenia (sleep-EEG); Epilepsy \\
\addlinespace
Theta ($\theta$) & 4--8 Hz & Drowsiness; memory & Cognitive decline; ADHD & ASD ($\theta$/$\beta$ ratio); Stress \\
\addlinespace
Alpha ($\alpha$) & 8--13 Hz & Relaxed wakefulness & Eyes-closed baseline; asymmetry & Stress (frontal asymmetry); BCI ($\mu$ rhythm) \\
\addlinespace
Beta ($\beta$) & 13--30 Hz & Active thinking; motor & Motor planning; concentration & ; BCI \\
\addlinespace
Gamma ($\gamma$) & 30--45 Hz & Cognitive processing & Perception; consciousness & ASD-focused (high-level cognition) \\
\bottomrule
\end{tabularx}
\vspace{0.5em}
\noindent\textit{Note.} Hz = hertz; ADHD = attention deficit hyperactivity disorder; $\mu$ rhythm = sensorimotor rhythm in the alpha range (8--12\,Hz) specific to motor cortex. The 0.5--45\,Hz band-pass filter (Figure~\ref{fig:eeg_preproc}) retains all five bands while excluding high-frequency muscle artifacts ($>$45\,Hz) and DC drift ($<$0.5\,Hz). Gamma range is truncated at 45\,Hz to avoid contamination from power-line interference (50/60\,Hz). Clinical significance informed by~\cite{niedermeyer2005eeg}.
\end{table}
\fi % AUTO-SUPPRESS non-ASD

%% MOVED TO SUPPLEMENTARY VOLUME — EEG data attributes table, 10-20 electrode system
%% table, and Fourier transform family table (Tables tab:eeg_data_attributes,
%% tab:1020_system, tab:fourier_family)
\subsubsection{EEG Signal Characteristics and Acquisition Standards}

%% SUPPRESSED FOR WORD COUNT

The signal processing pipeline rests on four mathematical foundations: the Continuous Wavelet Transform (Equation~\ref{eq:cwt}), the Short-Time Fourier Transform (Equation~\ref{eq:stft}), Sample Entropy (Equation~\ref{eq:sample_entropy}), and Phase-Locking Value (Equation~\ref{eq:plv}). Each formula is introduced at its point of use in the subsections below; Appendix~\ref{app:ch3_support} provides full derivations and parameter sensitivity analyses.

\paragraph{Continuous Wavelet Transform (CWT).}
The CWT decomposes the EEG signal $x(t)$ using a Morlet wavelet:
\begin{equation}
W(a,b) = \frac{1}{\sqrt{a}} \int_{-\infty}^{\infty} x(t) \, \psi^{*}\!\left(\frac{t-b}{a}\right) dt
\end{equation}
where $\psi(t) = \pi^{-1/4} e^{i\omega_0 t} e^{-t^2/2}$ is the Morlet wavelet ($\omega_0 = 6$), $a$ is the scale parameter, and $b$ is the translation parameter. The framework uses 64 logarithmically spaced frequency bins spanning 0.5--45~Hz with 128 time samples per epoch, producing 64$\times$128 scalogram images per channel.

\paragraph{Short-Time Fourier Transform (STFT).}
The STFT applies a sliding Hann window ($N = 256$ samples) with 75\% overlap and 512-point FFT:
\begin{equation}
X(m, \omega) = \sum_{n=0}^{N-1} x(n + mH) \, w(n) \, e^{-j\omega n}
\end{equation}
where $w(n)$ is the Hann window, $m$ is the frame index, and $H = N/4$ is the hop size. This produces 129$\times T$ spectrogram images, where $T$ depends on epoch length.

Table~\ref{tab:cwt_stft_specs} compares the two representations and their domain assignments.


\noindent\textbf{Time-Frequency Transformation Specifications.}

\needspace{24\baselineskip}
\begin{table}[!htbp]
\tablestripes
\caption{Time-Frequency Transformation Specifications}
\tablefontsize
\begin{tabularx}{\textwidth}{@{} >{\raggedright\arraybackslash}p{2.5cm} L L @{}}
\toprule
\thead{Parameter} & \thead{CWT (Morlet)} & \thead{STFT (Hann)} \\
\midrule
Wavelet/window & Morlet ($\omega_0 = 6$) & Hann ($N = 256$) \\
\addlinespace
Frequency range & 0.5--45 Hz (64 log bins) & 0--128 Hz (129 linear bins) \\
\addlinespace
Time resolution & 128 samples & Depends on epoch length \\
\addlinespace
Output shape & 64 $\times$ 128 per channel & 129 $\times$ $T$ per channel \\
\addlinespace
Overlap & N/A (continuous) & 75\% (hop = 64 samples) \\
\addlinespace
FFT size & N/A & 512 points \\
\addlinespace
Assigned model & 2D-CNN (scalograms; compared model) & Transformer (spectrograms; compared model) \\
\addlinespace
Clinical use & Multi-scale oscillatory patterns; alpha/beta/gamma band analysis & Drowsiness-specific alpha/theta transitions; steady-state evoked potentials \\
\bottomrule
\end{tabularx}
\vspace{0.5em}
\noindent\textit{Note.} CWT = continuous wavelet transform; STFT = short-time Fourier transform; FFT = fast Fourier transform; Hz = hertz; CNN = convolutional neural network. CWT provides superior frequency resolution at low frequencies; STFT provides uniform time-frequency resolution.
\end{table}

Figure~\ref{fig:eeg_preproc} details the eight-step pipeline. Procedural descriptions are in Appendix~\ref{app:ch3_support}.


\noindent\textbf{Eight-Step EEG Preprocessing Pipeline.}

\needspace{24\baselineskip}
\begin{figure}[!htbp]
\centering
\resizebox{0.97\textwidth}{!}{%
\begin{tikzpicture}[
  every node/.style={font=\fignodefont},
  pstep/.style={draw=ggunavy, fill=ggunavy!8, rounded corners=3pt, text width=13.5cm, minimum height=0.9cm, font=\small, align=left, inner sep=6pt},
  arr/.style={-{Stealth[length=4mm, width=3mm]}, very thick, ggunavy},
  node distance=0.8cm
]

\node[pstep, fill=lightblue] (s1) {\textbf{Step 1: Quality Check} --- Identify corrupted/missing recordings {\small (MNE-Python visual inspection)}};

\node[pstep, fill=lightorange, below=of s1] (s2) {\textbf{Step 2: Band-pass Filtering} --- Retain 0.5--45\,Hz; remove drift and noise {\small (Butterworth 4th-order)}};

\node[pstep, fill=lightteal, below=of s2] (s3) {\textbf{Step 3: Notch Filter} --- Remove power-line interference (50/60\,Hz) {\small (FIR notch filter)}};

\node[pstep, fill=lightpurple, below=of s3] (s4) {\textbf{Step 4: ICA Artifact Removal} --- Separate and remove ocular/muscular artifacts {\small (FastICA algorithm; ICLabel classification, $>$80\% non-brain rejection)}};

\node[pstep, fill=lightgreen, below=of s4] (s5) {\textbf{Step 5: Baseline Correction} --- Remove residual DC offset {\small (Epoch-mean subtraction)}};

\node[pstep, fill=lightyellow, below=of s5] (s6) {\textbf{Step 6: Normalisation} --- Standardise amplitude distributions {\small (Z-score normalisation)}};

\node[pstep, fill=lightpink, below=of s6] (s7) {\textbf{Step 7: Segmentation} --- Create fixed-length epochs (1--4\,s) {\small (Sliding window, 50\% overlap)}};

\node[pstep, fill=lightsky, below=of s7] (s8) {\textbf{Step 8: Transformation} --- Generate spectrograms, scalograms, CSP features {\small (STFT, CWT, CSP)}};
\draw[arr] (s1) -- (s2);
\draw[arr] (s2) -- (s3);
\draw[arr] (s3) -- (s4);
\draw[arr] (s4) -- (s5);
\draw[arr] (s5) -- (s6);
\draw[arr] (s6) -- (s7);
\draw[arr] (s7) -- (s8);
\end{tikzpicture}%
}% end resizebox
\caption{Eight-Step EEG Preprocessing Pipeline}
\vspace{0.5em}
\noindent\textit{Note.} ICA = independent component analysis; FIR = finite impulse response; CSP = common spatial pattern; STFT = short-time Fourier transform; CWT = continuous wavelet transform; Hz = hertz; s = seconds.
\end{figure}

\paragraph{ICA Component Rejection Criteria.} Independent components were classified using ICLabel~\cite{pion2019iclabel}; those with $>$80\% probability of non-brain origin were rejected across five artifact categories. This automated classification reduces subjective bias in manual artifact rejection and ensures reproducibility across datasets.

\begin{itemize}[leftmargin=*, itemsep=3pt]
  \item \textbf{Ocular artifacts:} Eye blinks and saccades identified by frontal topography and characteristic low-frequency time course.
  \item \textbf{Muscular artifacts:} Electromyographic contamination identified by broadband high-frequency spectral profile.
  \item \textbf{Cardiac artifacts:} Heart-related components identified by periodic QRS-complex morphology.
  \item \textbf{Line noise:} Power-line interference at 50/60~Hz identified by narrowband spectral peaks.
  \item \textbf{Channel noise:} Single-channel artifacts from electrode malfunction identified by spatially focal topography.
\end{itemize}


\noindent Domain-specific ICA implementations varied by dataset: FastICA for ASD; 20-component ICA for PD; wavelet denoising complementing ICA for Stress; and amplitude thresholding as the primary rejection method for BCI (due to the motor-imagery paradigm's distinct artifact profile).

%% ============================================================================
%% EEG SIGNAL PROCESSING PIPELINE (from EEG-RAG paper)
%% ============================================================================
%% MOVED TO APPENDIX: Redundant EEG Signal Processing Pipeline (fig:eeg_signal_pipeline, tab:eeg_filter_params) — overlaps with fig:eeg_preproc

%% SUPPRESSED — redundant with fig:eeg_preproc (TikZ version preferred over PNG)
\iffalse

\noindent To provide an integrated view of how raw EEG recordings are transformed into model-ready feature vectors, Figure~\ref{fig:eeg_pipeline_visual} presents the end-to-end preprocessing and feature extraction pipeline. This pipeline operationalises the artefact removal methods described above and connects them to the feature engineering stages that follow.


\needspace{24\baselineskip}
\begin{figure}[H]
    \centering
    \includegraphics[width=\textwidth,height=0.7\textheight,keepaspectratio]{eeg_pipeline.png}
    \caption{EEG Signal Preprocessing and Feature Extraction Pipeline}
    \label{fig:eeg_pipeline_visual}
    \vspace{0.5em}
\noindent\textit{Note.} ICA = independent component analysis; CAR = common average reference; SNR = signal-to-noise ratio; SMOTE = synthetic minority oversampling technique; PSD = power spectral density; CWT = continuous wavelet transform; PLV = phase-locking value; CNN = convolutional neural network. \textit{Source:} Author's own construction.
\end{figure}
\fi

\fi %% End suppressed sec1_signal_standardisation



\noindent All EEG signals were standardised to a common 256\,Hz sampling rate, bandpass-filtered (0.5--45\,Hz), and segmented into 4-second epochs with 50\% overlap. Table~\ref{tab:preproc_param_summary} consolidates parameters and reproducibility commitments; full derivations are in Appendix~C (Section~\ref{app:ch3_support}).

\paragraph{EEG Preprocessing Parameter and.}


\clearpage

{\tablestripes\tablefontsize
\begin{xltabular}{\textwidth}{@{} >{\raggedright\arraybackslash}p{2.8cm} >{\raggedright\arraybackslash}p{2.8cm} L >{\raggedright\arraybackslash}p{3.0cm} @{}}
\caption[EEG Preprocessing Parameter and Reproducibility Summary]{EEG Preprocessing Parameter and Reproducibility Summary. Single-page summary of each preprocessing stage with method, parameters, and validation threshold; full derivations in Appendix~C.}\label{tab:preproc_param_summary} \\
\toprule
\thead{Stage} & \thead{Method} & \thead{Parameters} & \thead{Validation Threshold} \\
\midrule
\endfirsthead
\endhead
\bottomrule
\endlastfoot
1. Acquisition & EDF/CSV via MNE-Python & 19--64 channels (10--20 system); standardised to 256\,Hz & File integrity checksum; consent confirmed; provenance logged \\
\addlinespace
2. Quality check & Visual + automated SNR & Per-channel SNR over 4\,s windows & SNR $\geq$ 10\,dB; channel-rejection $<$ 20\% \\
\addlinespace
3. Band-pass filter & Butterworth 4th-order, zero-phase & 0.5--45\,Hz & Stop-band attenuation $\geq$ 40\,dB \\
\addlinespace
4. Notch filter & FIR notch & 50\,Hz (EU/India) or 60\,Hz (US) & Residual line-noise $\leq$ $-$30\,dB \\
\addlinespace
5. Artefact removal & FastICA + ICLabel & Reject $\geq$ 80\% non-brain across 5 categories; PD: 20-component ICA; Stress: wavelet denoising; BCI: amplitude threshold & $\geq$ 60\% clean epochs retained \\
\addlinespace
6. Baseline correction & Epoch-mean subtraction & 200\,ms pre-stimulus baseline & Residual DC $\leq$ 1\,$\mu$V \\
\addlinespace
7. Normalisation & Z-score per channel/epoch & $\mu = 0$, $\sigma = 1$ & Channel mean $|\mu| \leq 0.05$ \\
\addlinespace
8. Segmentation & Sliding window & 4\,s epochs, 50\% overlap & Class-balance ratio $<$ 10:1 \\
\addlinespace
9. CWT transform & Continuous wavelet, Morlet & $\omega_0 = 6$; 64 log bins (0.5--45\,Hz); 64$\times$128 scalogram & Cone-of-influence masking \\
\addlinespace
10. STFT transform & Hann window & $N = 256$, 75\% overlap, 512-pt FFT; 129$\times T$ spectrogram & Spectral leakage $\leq$ $-$60\,dB \\
\addlinespace
11. Reproducibility seed & Deterministic random state & \texttt{seed=42} across NumPy, SciPy, scikit-learn, PyTorch & Bit-exact replay across 10-fold CV \\
\end{xltabular}}
\vspace{0.5em}{\small\noindent\textit{Note.} Reference implementation: \texttt{agenticfinder/eeg\_\allowbreak pipeline/preprocessing.py} with tests at \texttt{agenticfinder/tests/test\_\allowbreak preprocessing.py}; library versions pinned via project \texttt{requirements.txt}. SNR = signal-to-noise ratio; FIR = finite impulse response; ICA = independent component analysis; CWT/STFT = continuous wavelet/short-time Fourier transform.}


\iffalse %% Moved to Appendix C: EEG Pipeline Diagram (engineering detail)
\noindent\textbf{Secondary EEG Processing Pipeline.}

\needspace{24\baselineskip}
\begin{figure}[!htbp]
\centering
\resizebox{0.97\textwidth}{!}{%
\begin{tikzpicture}[
  every node/.style={font=\fignodefont},
  stepbox/.style={
    rectangle,
    rounded corners=4pt,
    draw=ggunavy,
    fill=ggunavy!6,
    minimum width=2.8cm,
    minimum height=1.1cm,
    text width=2.4cm,
    align=center,
    font=\small
  },
  arr/.style={-{Stealth[length=4mm, width=3mm]}, very thick}
]

\node[stepbox] (s1) at (0,0) {Acquire EEG};
\node[stepbox] (s2) at (3.6,0) {Remove Artefacts};
\node[stepbox] (s3) at (7.2,0) {Filter Signal};
\node[stepbox] (s4) at (10.8,0) {Segment Data};

\node[stepbox] (s5) at (1.6,-2.2) {Transform Signal};
\node[stepbox] (s6) at (4.8,-2.2) {Extract Features};
\node[stepbox] (s7) at (8.0,-2.2) {Normalise Input};
\node[stepbox] (s8) at (11.2,-2.2) {Run Inference};

\node[stepbox] (s9) at (4.0,-4.4) {Estimate Confidence};
\node[stepbox] (s10) at (8.0,-4.4) {Generate Explanation};

\draw[arr] (s1) -- (s2);
\draw[arr] (s2) -- (s3);
\draw[arr] (s3) -- (s4);

\draw[arr] (s4) -- (s5);
\draw[arr] (s5) -- (s6);
\draw[arr] (s6) -- (s7);
\draw[arr] (s7) -- (s8);

\draw[arr] (s8) -- (s9);
\draw[arr] (s9) -- (s10);

\end{tikzpicture}%
}
\caption{Secondary EEG Processing Pipeline}
\label{fig:eeg_pipeline_clean}
\vspace{0.5em}
\noindent\textit{Note.} Signal transformation includes FFT, PSD/SPDD, and SWT or wavelet-based analysis depending on domain requirements.
\end{figure}
\fi
\noindent Full details are provided in Appendix~C, Section~\ref{sec:appendix_ch3_eeg_pipeline_clean}.




%% ======================================================================
%% C-S2. PBS Image Preprocessing Pipeline
%% Moved from Chapter 3 — 2026-05-12
%% ======================================================================
\addcontentsline{toc}{section}{C-S2. PBS Image Preprocessing Pipeline} % [TOC-appendix-section-insertion]
\section*{C-S2. PBS Image Preprocessing Pipeline}
\addcontentsline{toc}{section}{C-S2. PBS Image Preprocessing Pipeline}
\label{sec:appc_pbs_preproc}

\iffalse %% Engineering detail moved to Appendix C --- sec2_pbs_preproc
%% SUPPRESSED FOR WORD COUNT
Unlike EEG preprocessing, which addresses temporal artefacts, the cancer domain requires spatial normalisation of microscopy images. Table~\ref{tab:img_preproc} summarises the five-step PBS image preprocessing pipeline, designed to standardise staining variation and class imbalance before classification.


\noindent\textbf{Five-Step PBS Image Preprocessing Pipeline.}

\needspace{24\baselineskip}
\begin{table}[H]
\tablestripes
\tablefontsize
\caption{Five-Step PBS Image Preprocessing Pipeline}
\begin{tabularx}{\textwidth}{@{} >{\raggedright\arraybackslash}p{0.8cm} L L @{}}
\toprule
\thead{Step} & \thead{Process} & \thead{Purpose} \\
\midrule
1 & Resize to 512$\times$512 & Standardise spatial dimensions across images \\
2 & HSV colour thresholding & Segment cell regions from background; normalise stain intensity \\
3 & Normalisation to [0,1] & Scale pixel intensities for neural network input \\
4 & GAN augmentation & Generate synthetic training samples to balance classes and improve generalisation \\
5 & Data splitting & Stratified 80/20 train/test with 10-fold cross-validation \\
\bottomrule
\end{tabularx}
\vspace{0.5em}
\noindent\textit{Note.} HSV = hue--saturation--value colour space; GAN = generative adversarial network; PBS = peripheral blood smear.
\end{table}


\noindent The five-step sequence is ordered deliberately: spatial standardisation (Steps~1--2) precedes intensity normalisation (Step~3), which precedes augmentation (Step~4), ensuring that GAN-generated images inherit the same normalised properties as real samples. This ordering prevents data leakage by applying augmentation only after the train/test split (Step~5), a safeguard that directly supports the internal validity claims tested in Chapter~\ref{chap:analysis}.

%% MOVED TO SUPPLEMENTARY VOLUME — Computer vision pipeline table (tab:cv_pipeline)
%% and signal vs image comparison table (tab:signal_vs_image)
\subsubsection{PBS Classification Pipeline}

PBS microscopy images undergo parallel preprocessing: resizing to 512$\times$512, HSV colour thresholding for stain normalisation, pixel intensity scaling to [0,1], and GAN-based augmentation. The GAN generator (ConvTranspose2d layers) synthesises realistic cell images, while the discriminator provides auxiliary classification, and a ResNet-18 classifier performs the final 4-class ALL classification. Full pipeline specifications are in the Supplementary Volume.

%% ============================================================================
%% GAN ARCHITECTURE FOR CANCER DETECTION (from Cancer paper)
%% ============================================================================

Figure~\ref{fig:gan_pipeline} illustrates the end-to-end GAN-based PBS cancer classification pipeline, from input image preprocessing through GAN-based augmentation to the final four-class ALL output. The GAN addresses class imbalance by synthesising realistic cell images (generator) and provides auxiliary classification signals through the discriminator's learned features.


\needspace{24\baselineskip}
\iffalse % AUTO-SUPPRESS non-ASD
\noindent\textbf{GAN-Based PBS Cancer Classification Pipeline.}

\needspace{24\baselineskip}
\begin{figure}[!htbp]
\centering
\resizebox{0.97\textwidth}{!}{%
\begin{tikzpicture}[
  every node/.style={font=\fignodefont},
  pbox/.style={rectangle, draw=ggunavy, fill=ggunavy!8, rounded corners=3pt,
    text width=2.8cm, minimum height=1.6cm, align=center, font=\small, inner sep=5pt},
  gbox/.style={rectangle, draw=ggunavy, fill=ggunavy!18, rounded corners=3pt,
    text width=3.2cm, minimum height=1.6cm, align=center, font=\small, inner sep=5pt},
  obox/.style={rectangle, draw=ggunavy, fill=ggunavy!15, rounded corners=3pt,
    text width=2.6cm, minimum height=1.6cm, align=center, font=\small\bfseries, inner sep=5pt},
  arr/.style={-{Stealth[length=4mm, width=3mm]}, very thick, ggunavy},
  node distance=0.8cm
]
% Row 1: Input processing

\node[pbox, fill=lightblue] (input) {\textbf{\color{ggunavy}Input PBS}\\[2pt]20,000 images\\512$\times$512 RGB};

\node[pbox, fill=lightorange, right=of input] (resize) {\textbf{\color{ggunavy}Resize}\\[2pt]512$\times$512\\standardise};

\node[pbox, fill=lightteal, right=of resize] (norm) {\textbf{\color{ggunavy}Normalise}\\[2pt]$[0,1]$ range\\HSV threshold};

% Row 1: GAN

\node[gbox, fill=lightpurple, right=of norm] (gan) {\textbf{\color{ggunavy}GAN}\\[2pt]Generator:\\151,808 params\\Discriminator:\\262,297 params};

% Row 2: Classification

\node[pbox, fill=lightgreen, right=of gan] (resnet) {\textbf{\color{ggunavy}ResNet-18}\\[2pt]Classifier\\6,149,220 params};

\node[obox, fill=lightyellow, right=of resnet] (output) {\textbf{\color{ggunavy}4-Class Output}\\[2pt]Benign\\Early Pre-B\\Pre-B, Pro-B};

% Arrows
\draw[arr] (input) -- (resize);
\draw[arr] (resize) -- (norm);
\draw[arr] (norm) -- (gan);
\draw[arr] (gan) -- (resnet);
\draw[arr] (resnet) -- (output);
\end{tikzpicture}%
}
\caption{GAN-Based PBS Cancer Classification Pipeline}
\vspace{0.5em}
\noindent\textit{Note.} PBS = peripheral blood smear; GAN = generative adversarial network; RGB = red--green--blue; HSV = hue--saturation--value. The generator uses ConvTranspose2d layers to synthesise realistic cell images; the discriminator provides auxiliary classification features. Total pipeline: 6,563,325 trainable parameters. \textit{Source:}~\cite{asthana2025cancer}.
\end{figure}
\fi % AUTO-SUPPRESS non-ASD

Table~\ref{tab:gan_model_params} documents the parameter counts for each pipeline component. The ResNet-18 classifier accounts for 93.7\% of the total parameters, reflecting the transfer learning strategy where ImageNet-pretrained convolutional features are fine-tuned on PBS microscopy images. The GAN generator and discriminator together comprise only 414,105 parameters (6.3\%), making the augmentation component computationally lightweight.


\needspace{24\baselineskip}
\iffalse % AUTO-SUPPRESS non-ASD
\noindent\textbf{GAN Cancer Classification Pipeline: Model.}

\needspace{24\baselineskip}
\begin{table}[H]
\tablestripes
\caption{GAN Cancer Classification Pipeline: Model Component Parameters}
\tablefontsize
\begin{tabularx}{\textwidth}{@{} >{\raggedright\arraybackslash}p{3.5cm} r >{\raggedright\arraybackslash}p{2cm} L @{}}
\toprule
\thead{Component} & \thead{Parameters} & \thead{Share (\%)} & \thead{Function} \\
\midrule
Generator & 151,808 & 2.3 & Synthesise realistic PBS cell images via ConvTranspose2d upsampling layers to augment training data \\
\addlinespace
Discriminator & 262,297 & 4.0 & Distinguish real from generated images; provide auxiliary classification features for the downstream classifier \\
\addlinespace
ResNet-18 classifier & 6,149,220 & 93.7 & Four-class ALL classification using ImageNet-pretrained features fine-tuned on PBS microscopy \\
\midrule
\textbf{Total} & \textbf{6,563,325} & \textbf{100.0} & --- \\
\bottomrule
\end{tabularx}
\vspace{0.5em}
\noindent\textit{Note.} ALL = acute lymphoblastic leukaemia; PBS = peripheral blood smear; GAN = generative adversarial network. Parameter counts from~\cite{asthana2025cancer}. The generator uses five ConvTranspose2d layers (64$\to$512 channels); the discriminator mirrors this with five Conv2d layers. ResNet-18 is initialised with ImageNet weights and fine-tuned with a learning rate of $10^{-4}$.
\end{table}
\fi % AUTO-SUPPRESS non-ASD

%% ============================================================================
%% GAN INTERNAL ARCHITECTURE (from published cancer paper: Asthana et al.)
%% ============================================================================
Figure~\ref{fig:gan_architecture} presents the internal architecture of the three GAN components---Generator, Discriminator, and Classifier---showing the layer-by-layer progression from latent noise vector to four-class ALL classification output.


\needspace{24\baselineskip}
\iffalse % AUTO-SUPPRESS non-ASD
\noindent\textbf{GAN Internal Architecture: Generator,.}

\needspace{24\baselineskip}
\begin{figure}[!htbp]
\centering
\resizebox{0.97\textwidth}{!}{%
\begin{tikzpicture}[
  every node/.style={font=\fignodefont},
  genbox/.style={rectangle, draw=ggunavy, fill=lightpurple, rounded corners=3pt,
    thick, minimum width=4.2cm, minimum height=1.0cm,
    align=center, font=\small, inner sep=3pt},
  disbox/.style={rectangle, draw=ggunavy, fill=lightorange, rounded corners=3pt,
    thick, minimum width=4.2cm, minimum height=1.0cm,
    align=center, font=\small, inner sep=3pt},
  clsbox/.style={rectangle, draw=ggunavy, fill=lightgreen, rounded corners=3pt,
    thick, minimum width=4.2cm, minimum height=1.0cm,
    align=center, font=\small, inner sep=3pt},
  hdr/.style={rectangle, draw=ggunavy, fill=ggunavy!15, rounded corners=3pt,
    thick, minimum width=4.4cm, minimum height=0.55cm,
    align=center, font=\small\bfseries, inner sep=3pt},
  arr/.style={-{Stealth[length=4mm, width=3mm]}, very thick, ggunavy},
  node distance=0.8cm,
]

% ===== GENERATOR (left column) =====

\node[hdr] (gh) at (0,0) {Generator (G) --- 151,808 params};

\node[genbox, below=0.6cm of gh] (g1) {Input: Latent $z \sim \mathcal{N}(0,1)$, dim=100};

\node[genbox, below=of g1] (g2) {ConvTranspose2d (latent$\to$512), 4$\times$4};

\node[genbox, below=of g2] (g3) {BatchNorm2d(512) + ReLU};

\node[genbox, below=of g3] (g4) {ConvTranspose2d (512$\to$256), stride 2};

\node[genbox, below=of g4] (g5) {ConvTranspose2d (256$\to$128), stride 2};

\node[genbox, below=of g5] (g6) {ConvTranspose2d (128$\to$64), stride 2};

\node[genbox, below=of g6] (g7) {ConvTranspose2d (64$\to$3), Tanh};

\node[genbox, fill=lightblue, below=0.15cm of g7] (gout) {Output: 64$\times$64$\times$3 synthetic PBS image};
\draw[arr] (g1) -- (g2);
\draw[arr] (g3) -- (g4);
\draw[arr] (g5) -- (g6);
\draw[arr] (g6) -- (g7);
\draw[arr] (g7) -- (gout);

% ===== DISCRIMINATOR (middle column) =====

\node[hdr] (dh) at (5.5,0) {Discriminator (D) --- 262,297 params};

\node[disbox, below=0.6cm of dh] (d1) {Input: 64$\times$64$\times$3 image};

\node[disbox, below=of d1] (d2) {Conv2d (3$\to$64), LeakyReLU(0.2)};

\node[disbox, below=of d2] (d3) {Conv2d (64$\to$128), BN, LeakyReLU};

\node[disbox, below=of d3] (d4) {Conv2d (128$\to$256), BN, LeakyReLU};

\node[disbox, below=of d4] (d5) {Conv2d (256$\to$512), BN, LeakyReLU};

\node[disbox, below=of d5] (d6) {Adversarial Head: Conv2d (512$\to$1)};

\node[disbox, fill=lightyellow, below=0.15cm of d6] (dout) {Output: Real/Fake (Sigmoid)};
\draw[arr] (d1) -- (d2);
\draw[arr] (d3) -- (d4);
\draw[arr] (d4) -- (d5);
\draw[arr] (d5) -- (d6);
\draw[arr] (d6) -- (dout);

% ===== CLASSIFIER (right column) =====

\node[hdr] (ch) at (11,0) {ResNet-18 Classifier --- 6,149,220 params};

\node[clsbox, below=0.6cm of ch] (c1) {Input: 64$\times$64$\times$3 image};

\node[clsbox, below=of c1] (c2) {Initial Conv2d (3$\to$64), BN, ReLU};

\node[clsbox, below=of c2] (c3) {Layer1 (BasicBlock $\times$ 2)};

\node[clsbox, below=of c3] (c4) {Layer2 (BasicBlock $\times$ 2)};

\node[clsbox, below=of c4] (c5) {Layer3 (BasicBlock $\times$ 2)};

\node[clsbox, below=of c5] (c6) {Layer4 (BasicBlock $\times$ 2)};

\node[clsbox, below=of c6] (c7) {AdaptiveAvgPool2d + Flatten};

\node[clsbox, below=of c7] (c8) {Fully Connected (512$\to$4)};

\node[clsbox, fill=lightblue, below=0.15cm of c8] (cout) {Output: 4-class (Benign/Early/Pre/Pro)};
\draw[arr] (c1) -- (c2);
\draw[arr] (c3) -- (c4);
\draw[arr] (c5) -- (c6);
\draw[arr] (c7) -- (c8);
\draw[arr] (c8) -- (cout);

% ===== Inter-component arrows =====
\draw[arr, dashed, ggunavy!60] (gout.east) -- ++(0.4,0) |- (d1.west)
  node[pos=0.25, right, font=\small\itshape\color{ggunavy!60}] {generated};
\draw[arr, dashed, ggunavy!60] (dout.east) -- ++(0.4,0) |- ([yshift=-0.15cm]c1.west)
  node[pos=0.25, right, font=\small\itshape\color{ggunavy!60}] {features};

\end{tikzpicture}%
}
\caption[GAN Internal Architecture]{GAN Internal Architecture: Generator, Discriminator, and ResNet-18 Classifier for PBS Cancer Classification}
\vspace{0.5em}
\noindent\textit{Note.} The Generator upsamples a 100-dimensional latent vector $z$ through five transposed convolutional layers to produce 64$\times$64 synthetic PBS images. The Discriminator distinguishes real from generated images using five convolutional layers with LeakyReLU activation. The ResNet-18 Classifier performs 4-class ALL subtype classification using ImageNet-pretrained weights fine-tuned with cross-entropy loss. Total parameters: 6,563,325. BN = batch normalisation. \textit{Source:} Asthana, P., Lalawat, R.\,S., Gond, S.\,S., \& Anand, S.\,K. (2025). A GAN-Based Framework for the Classification of Hematological Malignancies Using Peripheral Blood Smear Images.
\end{figure}
\fi % AUTO-SUPPRESS non-ASD

%% ============================================================================
%% 15-STEP PROCESSING FLOWCHART (from published cancer paper: Asthana et al.)
%% ============================================================================
Figure~\ref{fig:gan_15step} presents the complete 15-step processing sequence for GAN-based PBS image classification, showing the numbered flow from raw image loading through adversarial training to the final four-class output.


\needspace{24\baselineskip}
\iffalse % AUTO-SUPPRESS non-ASD
\noindent\textbf{Complete 15-Step Processing Flowchart for.}

\needspace{24\baselineskip}
\begin{figure}[!htbp]
\centering
\resizebox{0.97\textwidth}{!}{%
\begin{tikzpicture}[
  every node/.style={font=\fignodefont},
  sbox/.style={rectangle, draw=ggunavy, fill=ggunavy!8, rounded corners=3pt,
    thick, minimum width=2.4cm, minimum height=1.1cm,
    align=center, font=\small, inner sep=4pt},
  decision/.style={diamond, draw=ggunavy, fill=lightyellow, aspect=2,
    thick, align=center, font=\small, inner sep=2pt},
  terminal/.style={rectangle, draw=ggunavy, fill=ggunavy!20, rounded corners=8pt,
    thick, minimum width=1.6cm, minimum height=1.0cm,
    align=center, font=\small\bfseries, inner sep=3pt},
  num/.style={circle, draw=ggunavy, fill=ggunavy, text=white,
    font=\small\bfseries, minimum size=0.4cm, inner sep=0pt},
  arr/.style={-{Stealth[length=4mm, width=3mm]}, very thick, ggunavy},
  node distance=0.8cm
]

% Row 1: Steps 1-5 (preprocessing)

\node[terminal] (start) {Start};

\node[sbox, right=of start] (s1) {Load PBS\\Images};

\node[num, above left=0.01cm and 0.01cm of s1.north east] {1};

\node[sbox, right=of s1] (s2) {Resize to\\512$\times$512};

\node[num, above left=0.01cm and 0.01cm of s2.north east] {2};

\node[sbox, right=of s2] (s3) {Normalise\\$[0,1]$};

\node[num, above left=0.01cm and 0.01cm of s3.north east] {3};

\node[sbox, right=of s3] (s4) {Data\\Augment.};

\node[num, above left=0.01cm and 0.01cm of s4.north east] {4};

\node[sbox, right=of s4] (s5) {Train/Val\\Split};

\node[num, above left=0.01cm and 0.01cm of s5.north east] {5};

\draw[arr] (start) -- (s1);
\draw[arr] (s1) -- (s2);
\draw[arr] (s2) -- (s3);
\draw[arr] (s3) -- (s4);
\draw[arr] (s4) -- (s5);

% Row 2: Steps 6-10 (GAN training)

\node[sbox, below=0.8cm of s1] (s6) {Sample\\$z \sim \mathcal{N}(0,1)$};

\node[num, above left=0.01cm and 0.01cm of s6.north east] {6};

\node[sbox, fill=lightpurple, right=of s6] (s7) {Generator\\$G(z)$};

\node[num, above left=0.01cm and 0.01cm of s7.north east] {7};

\node[sbox, fill=lightorange, right=of s7] (s8) {Discriminator\\$D(x)$};

\node[num, above left=0.01cm and 0.01cm of s8.north east] {8};

\node[decision, right=0.6cm of s8] (s9) {Adv.\\Loss?};

\node[num, above=0.01cm of s9.north] {9};

\node[sbox, right=0.6cm of s9] (s10) {Update\\$G$ \& $D$};

\node[num, above left=0.01cm and 0.01cm of s10.north east] {10};

\draw[arr] (s5.south) -- ++(0,-0.3) -| (s6.north);
\draw[arr] (s6) -- (s7);
\draw[arr] (s7) -- (s8);
\draw[arr] (s8) -- (s9);
\draw[arr] (s9) -- node[above, font=\small] {Yes} (s10);
\draw[arr, dashed] (s10.south) -- ++(0,-0.3) -| (s7.south);

% Row 3: Steps 11-15 (classification)

\node[sbox, fill=lightgreen, below=0.8cm of s6] (s11) {Classifier\\$C(x)$};

\node[num, above left=0.01cm and 0.01cm of s11.north east] {11};

\node[sbox, right=of s11] (s12) {Softmax\\Output};

\node[num, above left=0.01cm and 0.01cm of s12.north east] {12};

\node[sbox, right=of s12] (s13) {Predict\\Class};

\node[num, above left=0.01cm and 0.01cm of s13.north east] {13};

\node[sbox, fill=lightblue, right=of s13] (s14) {Output:\\Benign/Early/\\Pre/Pro};

\node[num, above left=0.01cm and 0.01cm of s14.north east] {14};

\node[terminal, right=of s14] (s15) {End};

\node[num, above=0.01cm of s15.north] {15};

\draw[arr] (s9.south) -- node[right, font=\small] {No} ++(0,-0.5) -| (s11.north);
\draw[arr] (s11) -- (s12);
\draw[arr] (s12) -- (s13);
\draw[arr] (s13) -- (s14);
\draw[arr] (s14) -- (s15);

\end{tikzpicture}%
}
\caption[GAN Processing Flowchart]{Complete 15-Step Processing Flowchart for GAN-Based PBS Cancer Classification}
\vspace{0.5em}
\noindent\textit{Note.} Steps 1--5 handle preprocessing (resize, normalise, augment, split). Steps 6--10 perform GAN adversarial training (sample latent vector, generate synthetic images, discriminate, compute adversarial loss, update weights). Steps 11--15 perform classification (ResNet-18 classifier, softmax, predict 4-class ALL output). The adversarial loop (steps 7--10) iterates until convergence (50 epochs). PBS = peripheral blood smear. \textit{Source:}~\cite{asthana2025cancer}.
\end{figure}
\fi % AUTO-SUPPRESS non-ASD

\fi %% End suppressed sec2_pbs_preproc



The complete imaging preprocessing protocol is documented in Appendix~C (Section~\ref{app:ch3_support}).


%% ======================================================================
%% C-S3. Deep Learning Architectures
%% Moved from Chapter 3 — 2026-05-12
%% ======================================================================
\addcontentsline{toc}{section}{C-S3. Deep Learning Architectures} % [TOC-appendix-section-insertion]
\section*{C-S3. Deep Learning Architectures}
\addcontentsline{toc}{section}{C-S3. Deep Learning Architectures}
\label{sec:appc_dl_arch}
\iffalse %% Engineering detail moved to Appendix C --- sec5_dl_arch

The framework evaluates ten deep learning architectures as \textit{compared models} for EEG domains and as the \textit{best-performing model} for imaging domains . For EEG-based classification, DL architectures were systematically compared against the VotingClassifier ensemble (Section~\ref{sec:voting_ensemble}) but did not outperform it on the engineered feature sets used in RGAIG. DL remains the deployed model for cancer imaging, where spatial feature hierarchies in histopathological images favour convolutional architectures. Table~\ref{tab:dl_architecture_comparison} summarises the ten architectures, their input modalities, primary use cases, and evaluation status.


\needspace{24\baselineskip}
\iffalse % AUTO-SUPPRESS non-ASD
\noindent\textbf{Deep Learning Architecture Comparison: Input.}

\needspace{24\baselineskip}
\begin{table}[!htbp]
\tablestripes
\caption[Deep Learning Architecture Comparison]{Deep Learning Architecture Comparison: Input Type, Primary Use, and Evaluation Status}
\tablefontsize
\begin{tabularx}{\textwidth}{@{} >{\raggedright\arraybackslash}p{0.4cm} >{\raggedright\arraybackslash}p{2.8cm} >{\raggedright\arraybackslash}p{2.5cm} L >{\raggedright\arraybackslash}p{2.2cm} @{}}
\toprule
\thead{\#} & \thead{Architecture} & \thead{Input Type} & \thead{Primary Use} & \thead{Deployment Status} \\
\midrule
1 & 1D-CNN & Raw EEG time-series & Temporal pattern extraction from single-channel signals & Compared \\
\addlinespace
2 & 2D-CNN & CWT scalograms & Spectro-temporal image classification for EEG & Compared \\
\addlinespace
3 & CNN-LSTM (bi-directional + self-attention) & STFT spectrograms & Captures both spatial and temporal EEG dynamics & Compared \\
\addlinespace
4 & DeepConvNet~\cite{schirrmeister2017deep} & Raw multi-channel EEG & End-to-end deep convolutional decoding of EEG & Compared \\
\addlinespace
5 & EEGNet~\cite{lawhern2018eegnet} & Raw multi-channel EEG & Compact, generalisable EEG classification across paradigms & Compared \\
\addlinespace
6 & GAN generator (ConvTranspose2d) & PBS microscopy images & Synthetic image generation for class balancing & Deployed  \\
\addlinespace
7 & ResNet-18 classifier & PBS microscopy images & 4-class ALL leukaemia classification from cell images & Deployed  \\
\addlinespace
8 & Vision Transformer (ViT) & Patch-based image tokens & Global self-attention over image patches for classification & Compared \\
\addlinespace
9 & Transformer encoder & EEG feature sequences & Self-attention over temporal EEG segments & Compared \\
\addlinespace
10 & MLP meta-learner (Ultra Stacking) & Stacked base-model predictions & Multi-disease screening via meta-learned ensemble & Deployed (multi-disease) \\
\bottomrule
\end{tabularx}
\vspace{0.5em}
\noindent\textit{Note.} CNN = convolutional neural network; LSTM = long short-term memory; CWT = continuous wavelet transform; STFT = short-time Fourier transform; GAN = generative adversarial network; PBS = peripheral blood smear; ALL = acute lymphoblastic leukaemia; ViT = vision transformer; MLP = multilayer perceptron. ``Compared'' = evaluated but not deployed (VotingClassifier outperformed on EEG features); ``Deployed'' = selected as best-performing model for the indicated domain. Full architecture specifications are in Appendix~C.
\end{table}
\fi % AUTO-SUPPRESS non-ASD

%% ENGINEERING DETAIL SUPPRESSED — DBA summary retained; full specifications in Appendix~C
%% (Architecture parameters, equations, layer configurations)
\iffalse
%% BEGIN SUPPRESSED ENGINEERING: DL Ensemble + Ultra Stacking Architecture Details
\subsubsection{Three-Model DL Ensemble}
Three DL architectures evaluated: EEGNet (2,548 params), 2D-CNN (148,226 params), Transformer (502,018 params).
Ensemble prediction: $P_{\text{ensemble}} = \sum_{i=1}^{3} w_i \cdot p_i$. Total: 652,792 params.
\subsubsection{Ultra Stacking Ensemble (Multi-Disease Classification)}
Three-layer stacking: Layer 1 (15+ base classifiers), Layer 2 (mutual information feature selection, top 300), Layer 3 (MLP meta-learner 128-64-32).
%% END SUPPRESSED ENGINEERING
\fi


The framework's model selection strategy follows an adaptive decision principle: the best-performing architecture is selected per clinical domain based on empirical evidence, not architectural preference. Three ensemble strategies were evaluated (Table~\ref{tab:ensemble_comparison}):
\begin{enumerate}[leftmargin=*, itemsep=1pt, label=(\arabic*)]
    \item VotingClassifier for single-condition EEG detection.
    \item A three-model DL ensemble as a compared alternative.
    \item Ultra Stacking for multi-disease screening.
\end{enumerate}
\noindent Classification results are reported in Chapter~\ref{chap:analysis}. This adaptive selection represents a governance-aligned decision: organisations conceptually consider the architecture that maximises diagnostic accuracy for their specific clinical use case rather than defaulting to the most complex model.

Table~\ref{tab:ensemble_comparison} compares the three ensemble strategies and their evaluation contexts.


\noindent\textbf{Ensemble Architecture Comparison:.}

\needspace{24\baselineskip}
\begin{table}[!htbp]
\tablestripes
\caption[Ensemble Architecture Comparison]{Ensemble Architecture Comparison: VotingClassifier vs Three-Model DL vs Ultra Stacking}
\tablefontsize
\begin{tabularx}{\textwidth}{@{} >{\raggedright\arraybackslash}p{2.3cm} L >{\raggedright\arraybackslash}p{3cm} L @{}}
\toprule
\thead{Aspect} & \thead{VotingClassifier
ewline(Deployed)} & \thead{Three-Model DL
ewline(Compared)} & \thead{Ultra Stacking
ewline(Multi-Disease)} \\
\midrule
Base models & ExtraTrees (300) + RF (300) & EEGNet, 2D-CNN, Transformer & 15+ classifiers (tree, neural, kernel) \\
\addlinespace
Input & 140+ engineered features & Raw EEG, CWT scalograms, STFT spectrograms & 47 engineered features per channel \\
\addlinespace
Fusion & Soft voting (probability averaging) & Weighted probability averaging & Stacking with MLP meta-learner \\
\addlinespace
Validation & 5-fold stratified CV & LOSO-CV & 5-fold stratified CV \\
\addlinespace
Use case & Single-condition EEG detection & Alternative EEG architecture & Multi-disease screening \\
\addlinespace
Inference & $<$5 ms (CPU) & 32 ms (GPU) & 15 ms (CPU) \\
\addlinespace
Status & \textbf{Best-performing; deployed} & Compared; not deployed & Deployed for multi-disease \\
\bottomrule
\end{tabularx}
\vspace{0.5em}
\noindent\textit{Note.} RF = Random Forest; LOSO = leave-one-subject-out; CV = cross-validation; DL = deep learning; CWT = continuous wavelet transform; STFT = short-time Fourier transform; MLP = multilayer perceptron. The VotingClassifier is the final deployed model for all EEG-based domains (ASD, PD, Stress). DL architectures were systematically compared but did not outperform the ensemble on engineered features.
\end{table}

The key managerial decision is model selection: the VotingClassifier ensemble is deployed for all EEG domains (ASD, PD, Stress) because it outperforms DL alternatives on engineered features, while GAN+ResNet-18 is deployed for cancer imaging where spatial feature hierarchies favour convolutional architectures. This evidence-based selection tests H2: whether governance-driven model selection---choosing the best-performing architecture per domain rather than defaulting to complexity---yields higher classification accuracy.

%% ENGINEERING FIGURE SUPPRESSED — CNN-LSTM architecture diagram moved to Appendix C
\iffalse

\needspace{24\baselineskip}
\begin{figure}[!htbp]
    \centering
    \begin{tikzpicture}[
  every node/.style={font=\fignodefont},
  % Local styles for this diagram
        cnnblock/.style={draw=ggunavy, fill=ggunavy!8, rounded corners=3pt,
            thick, minimum width=2.1cm, minimum height=1.6cm,
            align=center, font=\small, inner sep=4pt},
        procblock/.style={draw=ggunavy, fill=ggunavy!8, rounded corners=3pt,
            thick, minimum width=2.1cm, minimum height=1.2cm,
            align=center, font=\small, inner sep=4pt},
        ioblock/.style={draw=ggunavy, fill=ggunavy!15, rounded corners=3pt,
            thick, minimum width=2.1cm, minimum height=1.2cm,
            align=center, font=\small\bfseries, inner sep=4pt},
        grplabel/.style={font=\small\bfseries\color{ggunavy}, anchor=south},
        arr/.style={-{Stealth[length=4mm, width=3mm]}, very thick, ggunavy},
        node distance=0.45cm,
    ]

    % === Row 1: Input → Conv Blocks → Flatten (left to right) ===
    
\node[ioblock] (input) {Input\\{\small 128\,$\times$\,128\,$\times$\,3}\\{\small CWT Scalogram}};

    % Conv Block 1
    
\node[cnnblock, right=0.6cm of input] (conv1) {Conv2D\\{\small 32 filters, 3\,$\times$\,3}\\{\small BN\,$\to$\,ReLU}\\{\small MaxPool 2\,$\times$\,2}};

    % Conv Block 2
    
\node[cnnblock, right=0.45cm of conv1] (conv2) {Conv2D\\{\small 64 filters, 3\,$\times$\,3}\\{\small BN\,$\to$\,ReLU}\\{\small MaxPool 2\,$\times$\,2}};

    % Conv Block 3
    
\node[cnnblock, right=0.45cm of conv2] (conv3) {Conv2D\\{\small 128 filters, 3\,$\times$\,3}\\{\small BN\,$\to$\,ReLU}\\{\small MaxPool 2\,$\times$\,2}};

    % Flatten + Reshape
    
\node[procblock, right=0.6cm of conv3] (flatten) {Flatten\\+\\Reshape};

    % === Grouping box around Conv Blocks ===
    \begin{scope}[on background layer]
        
\node[draw=ggunavy, dashed, rounded corners=5pt, thick,
              fill=ggunavy!3, inner sep=8pt,
              fit=(conv1)(conv2)(conv3),
              label={[grplabel]above:Convolutional Feature Extraction}] (convgroup) {};
    \end{scope}

    % === Row 2: Bi-LSTM → Attention → Dense → Dropout → Output ===
    % Position below flatten, then flow right to left
    
\node[procblock, below=1.6cm of flatten, minimum width=2.3cm] (bilstm)
        {Bi-LSTM\\{\small 128 units}};

    
\node[procblock, left=0.45cm of bilstm] (attention)
        {Self-\\Attention};

    
\node[procblock, left=0.45cm of attention] (dense)
        {Dense\\{\small 64 units, ReLU}};

    
\node[procblock, left=0.45cm of dense] (dropout)
        {Dropout\\{\small $p = 0.5$}};

    
\node[ioblock, left=0.6cm of dropout] (output)
        {Output\\{\small Sigmoid / Softmax}};

    % === Grouping box around sequence processing ===
    \begin{scope}[on background layer]
        
\node[draw=ggunavy, dashed, rounded corners=5pt, thick,
              fill=ggunavy!3, inner sep=8pt,
              fit=(bilstm)(attention),
              label={[grplabel]below:Temporal Sequence Learning}] (seqgroup) {};
    \end{scope}

    % === Arrows: Row 1 ===
    \draw[arr] (input) -- (conv1);
    \draw[arr] (conv1) -- (conv2);
    \draw[arr] (conv2) -- (conv3);
    \draw[arr] (conv3) -- (flatten);

    % === Arrow: Row 1 → Row 2 (vertical turn) ===
    \draw[arr] (flatten.south) -- (bilstm.north);

    % === Arrows: Row 2 (right to left) ===
    \draw[arr] (bilstm) -- (attention);
    \draw[arr] (attention) -- (dense);
    \draw[arr] (dense) -- (dropout);
    \draw[arr] (dropout) -- (output);

    \end{tikzpicture}
    \vspace{0.5em}
    \caption[CNN-LSTM Hybrid Architecture]{CNN-LSTM Hybrid Architecture for EEG-Based Biomedical Diagnosis (Compared Model)}
    \label{fig:cnn_lstm_arch}
    \vspace{0.5em}
\noindent\textit{Note.} BN = batch normalization; ReLU = rectified linear unit; Bi-LSTM = bidirectional long short-term memory; CWT = continuous wavelet transform. Input dimensions correspond to 128$\times$128$\times$3 CWT scalogram representations. \textit{Source:} Author's own construction.
\end{figure}
\fi
%% Phantom label for cross-references

%% ============================================================================
%% CNN-ASD MODEL ARCHITECTURE (from published ICSIT 2025 paper: Asthana et al.)
%% ============================================================================

\noindent\textbf{Proposed CNN-ASD Model Architecture for EEG-Based.}

\needspace{24\baselineskip}
\begin{figure}[!htbp]
    \centering
    \resizebox{0.97\textwidth}{!}{%
    \begin{tikzpicture}[
  every node/.style={font=\fignodefont},
  block/.style={draw=ggunavy, fill=ggunavy!8, rounded corners=3pt,
            thick, minimum width=7cm, minimum height=1.0cm,
            align=center, font=\small, inner sep=4pt},
        convblock/.style={draw=ggunavy, fill={#1}, rounded corners=3pt,
            thick, minimum width=7cm, minimum height=1.0cm,
            align=center, font=\small, inner sep=4pt},
        noteblock/.style={font=\small\color{ggunavy!60}, align=center},
        arr/.style={-{Stealth[length=4mm, width=3mm]}, very thick, ggunavy},
        node distance=0.8cm,
    ]

    % Input
    
\node[block, fill=ggunavy!15, font=\small\bfseries] (input) {Input (224\,$\times$\,224\,$\times$\,3)};

    % Conv Block 1
    
\node[convblock=orange!15, below=0.6cm of input] (conv1) {Conv2D (32 filters, 3\,$\times$\,3, ReLU, same padding)};
    
\node[noteblock, below=0.01cm of conv1] (note1) {\textit{3\,$\times$\,3 kernel $\mid$ ReLU $\mid$ Same padding}};
    
\node[convblock=orange!10, below=0.01cm of note1] (bn1) {BatchNorm};
    
\node[convblock=yellow!15, below=0.12cm of bn1] (mp1) {MaxPool (2\,$\times$\,2)};

    % Conv Block 2
    
\node[convblock=green!12, below=0.6cm of mp1] (conv2) {Conv2D (64 filters, 3\,$\times$\,3, ReLU, same padding)};
    
\node[noteblock, below=0.01cm of conv2] (note2) {\textit{3\,$\times$\,3 kernel $\mid$ ReLU $\mid$ Same padding}};
    
\node[convblock=green!8, below=0.01cm of note2] (bn2) {BatchNorm};
    
\node[convblock=yellow!15, below=0.12cm of bn2] (mp2) {MaxPool (2\,$\times$\,2)};

    % Conv Block 3
    
\node[convblock=red!10, below=0.6cm of mp2] (conv3) {Conv2D (128 filters, 3\,$\times$\,3, ReLU, same padding)};
    
\node[noteblock, below=0.01cm of conv3] (note3) {\textit{3\,$\times$\,3 kernel $\mid$ ReLU $\mid$ Same padding}};
    
\node[convblock=red!8, below=0.01cm of note3] (bn3) {BatchNorm};
    
\node[convblock=yellow!15, below=0.12cm of bn3] (mp3) {MaxPool (2\,$\times$\,2)};

    % Flatten + Dense + Output
    
\node[convblock=blue!10, below=0.6cm of mp3] (flatten) {Flatten};
    
\node[noteblock, below=0.01cm of flatten] (noteflt) {\textit{100,352-dim vector}};
    
\node[convblock=blue!15, below=0.01cm of noteflt] (dense) {Dense (128 units, ReLU)};
    
\node[convblock=orange!12, below=0.12cm of dense] (dropout) {Dropout (0.5)};
    
\node[block, fill=ggunavy!15, font=\small\bfseries, below=0.6cm of dropout] (output) {Output (2 classes: Normal / ASD)};

    % Arrows
    \draw[arr] (input) -- (conv1);
    \draw[arr] (bn1) -- (mp1);
    \draw[arr] (mp1) -- (conv2);
    \draw[arr] (bn2) -- (mp2);
    \draw[arr] (mp2) -- (conv3);
    \draw[arr] (bn3) -- (mp3);
    \draw[arr] (mp3) -- (flatten);
    \draw[arr] (dense) -- (dropout);
    \draw[arr] (dropout) -- (output);

    % Grouping bracket labels
    
\node[font=\small\bfseries\color{ggunavy}, rotate=90, anchor=south] at ([xshift=-4cm]conv1.west |- bn1) {Block 1};
    
\node[font=\small\bfseries\color{ggunavy}, rotate=90, anchor=south] at ([xshift=-4cm]conv2.west |- bn2) {Block 2};
    
\node[font=\small\bfseries\color{ggunavy}, rotate=90, anchor=south] at ([xshift=-4cm]conv3.west |- bn3) {Block 3};

    \end{tikzpicture}%
    }
    \vspace{0.5em}
    \caption[CNN-ASD Model Architecture]{Proposed CNN-ASD Model Architecture for EEG-Based Autism Spectrum Disorder Classification}
    \label{fig:cnn_asd_arch}
    \vspace{0.5em}
\noindent\textit{Note.} Three convolutional blocks with progressively increasing filter counts (32, 64, 128) extract hierarchical features from 224\,$\times$\,224\,$\times$\,3 CWT scalogram inputs. Dropout ($p$\,=\,0.5) prevents overfitting on the small KAU dataset ($N$\,=\,16). The model outperforms BiLSTM+Attention and ResNet50 baselines on the ASD binary classification task. \textit{Source:}~\cite{asthana2025asd}.
\end{figure}

%% ============================================================================
%% ASD CLASSIFICATION PIPELINE (from published ICSIT 2025 paper)
%% ============================================================================

\noindent\textbf{EEG-Based ASD Classification Pipeline: Signal.}

\needspace{24\baselineskip}
\begin{figure}[!htbp]
    \centering
    \begin{tikzpicture}[
  every node/.style={font=\fignodefont},
  stage/.style={draw=ggunavy, fill=ggunavy!8, rounded corners=4pt,
            thick, minimum width=2.6cm, minimum height=1.8cm,
            align=center, font=\small, inner sep=5pt},
        classifier/.style={draw=ggunavy, fill=#1, rounded corners=3pt,
            thick, minimum width=2.2cm, minimum height=0.65cm,
            align=center, font=\small, inner sep=3pt},
        result/.style={draw=ggunavy, fill=#1, rounded corners=10pt,
            thick, minimum width=1.4cm, minimum height=0.65cm,
            align=center, font=\small\bfseries, inner sep=3pt},
        arr/.style={-{Stealth[length=4mm, width=3mm]}, very thick, ggunavy},
        darr/.style={-{Stealth[length=4mm, width=3mm]}, very thick, ggunavy, dashed},
    ]

    % Stage 1: EEG Acquisition
    
\node[stage] (acq) {EEG Signal\\Acquisition\\{\small 16 ch, 256 Hz}\\{\small KAU dataset}};

    % Stage 2: Pre-processing
    
\node[stage, right=0.8cm of acq] (preproc) {Pre-Processing\\{\small 12th-order IIR}\\{\small (40 Hz cutoff)}\\{\small Z-score norm.}};

    % Stage 3: Feature Extraction
    
\node[stage, right=0.8cm of preproc] (feat) {Feature\\Extraction\\{\small CWT (Morlet)}\\{\small 4-sec segments}};

    % Stage 4: Classifiers
    
\node[classifier=cyan!15, right=1.2cm of feat, yshift=1.0cm] (bilstm) {Bi-LSTM + AM};
    
\node[classifier=green!15, right=1.2cm of feat, yshift=0cm] (resnet) {ResNet50};
    
\node[classifier=orange!20, right=1.2cm of feat, yshift=-1.0cm] (cnnasd) {\textbf{CNN-ASD}\\{\small (Proposed)}};

    % Classifier group label
    
\node[font=\small\bfseries\color{ggunavy}, above=0.15cm of bilstm] {Classifiers};

    % Stage 5: Output
    
\node[result=green!20, right=1.2cm of bilstm, yshift=0.5cm] (normal) {Normal};
    
\node[result=red!20, right=1.2cm of cnnasd, yshift=-0.5cm] (asd) {ASD};

    % Arrows: main pipeline
    \draw[arr] (acq) -- (preproc);
    \draw[arr] (preproc) -- (feat);
    \draw[arr] (feat) -- (bilstm);
    \draw[arr] (feat) -- (resnet);
    \draw[arr] (feat) -- (cnnasd);

    % Arrows: classifiers to output
    \draw[darr] (bilstm.east) -- ++(0.4,0) |- (normal.west);
    \draw[darr] (resnet.east) -- ++(0.3,0);
    \draw[darr] (cnnasd.east) -- ++(0.4,0) |- (asd.west);

    \end{tikzpicture}
    \vspace{0.5em}
    \caption[ASD Classification Pipeline]{EEG-Based ASD Classification Pipeline: Signal Acquisition to Binary Diagnosis}
    \label{fig:asd_classification_pipeline}
    \vspace{0.5em}
\noindent\textit{Note.} Raw EEG signals from 16 channels at 256\,Hz undergo bandpass filtering (12th-order IIR, 40\,Hz cutoff), z-score normalisation, and CWT scalogram conversion before being classified by three architectures. The proposed CNN-ASD model is benchmarked against Bi-LSTM with Attention Mechanism (AM) and ResNet50. \textit{Source:}~\cite{asthana2025asd}.
\end{figure}
\fi %% End suppressed sec5_dl_arch



\noindent Ten deep learning architectures were evaluated, selected for their established performance in biomedical signal classification. Architecture specifications are in Appendix~C.


%% ============================================================================
%% AGENTIC AI DISEASE FINDER — SYSTEM ARCHITECTURE (from IEEE TBME paper)
%% ============================================================================
\paragraph{Agentic AI Disease Finder} % [Part B]
The Agentic AI Disease Finder is included here only as a methodological precedent demonstrating how preprocessing, feature extraction, classification, and explanation can be orchestrated within one analytical pipeline. In the final thesis logic, this material functions as supporting implementation context rather than as a primary method block for the ASD validation study.

Accordingly, the detailed architecture diagram, representative ROC/precision-recall curves, and full per-dataset LOSO performance table are moved to Appendix~C. The main methodological point retained here is that the prior system used a staged EEG pipeline with artefact removal, dual-path feature extraction, ensemble classification, and RAG-supported explanation generation under subject-level validation constraints.

%% ======================================================================
%% C-S5. GenAI Output Evaluation Framework
%% Moved from Chapter 3 — 2026-05-12 (round 2)
%% ======================================================================
\addcontentsline{toc}{section}{C-S5. GenAI Output Evaluation Framework} % [TOC-appendix-section-insertion]
\section*{C-S5. GenAI Output Evaluation Framework}
\addcontentsline{toc}{section}{C-S5. GenAI Output Evaluation Framework}
\label{sec:appc_genai_eval}
\iffalse
Three frameworks: DeepEval (14 metrics), RAGAS (8 metrics), G-Eval (LLM-based scoring). Detailed comparison in Supplementary Volume.
\fi

RGAIG combines three evaluation frameworks (the corresponding table): DeepEval for faithfulness, RAGAS for retrieval quality, and G-Eval for narrative coherence.

\noindent The main methodological point retained here is triangulation: retrieval quality, faithfulness, and narrative quality were assessed through separate but complementary evaluation frameworks. Detailed framework-by-framework metric inventories are retained outside the main chapter.


\paragraph{RGAIG Output Quality Metrics.}

The output quality metrics and acceptance thresholds applied during RGAIG validation testing ($n$=500 generated explanations across five clinical domains) are reported in Chapter~\ref{chap:analysis}. Predefined acceptance thresholds are: faithfulness $\geq$0.85, answer relevancy $\geq$0.80, context recall $\geq$0.80, and hallucination rate $\leq$5\%.


%% MOVED TO SUPPLEMENTARY VOLUME — RAG Governance 15 Assurance Layers figure
%% (fig:rag_governance_layers)
\paragraph{RAG Governance: Fifteen Assurance Layers.}

The RAG pipeline implements 15 automated assurance layers (A1--A15) spanning faithfulness validation (DeepEval + RAGAS), input sanitisation (prompt/SQL injection blocking), role-based access control, contradiction detection (NLI-based), hallucination monitoring (citation accuracy $\geq$90\%), lineage versioning (hash-based immutable registry), and regulatory compliance verification (HIPAA \S164.312, GDPR Art.~30). Nine layers execute synchronously per query; six operate asynchronously on scheduled monitoring. Detailed layer specifications are in the Supplementary Volume.

%% ---------------------------------------------------------------------------
%% THREE ARCHITECTURE DIAGRAMS from Agentic AI Disease Finder Master Analysis
%% Source: eeg_complete_master.pdf (Asthana, 2025)
%% ---------------------------------------------------------------------------

the corresponding figure presents the C4 four-layer system architecture~\cite{asthana2025eegmaster}.
The implementation also follows a layered system architecture separating interface, orchestration, analytical, and infrastructure responsibilities. In the final thesis, this is treated as supplementary implementation evidence rather than a required part of the core research-method description.


the corresponding figure illustrates the seven-stage ML pipeline architecture from signal acquisition to clinical explanation.


\noindent\textbf{End-to-End ML Pipeline Architecture with Seven.}

\needspace{24\baselineskip}
\begin{figure}[!htbp]
\centering
\resizebox{0.99\textwidth}{!}{%
\begin{tikzpicture}[
  every node/.style={font=\fignodefont},
  stage/.style={draw=ggunavy, thick, rounded corners=4pt, fill=#1,
    minimum width=2.65cm, minimum height=3.0cm, inner sep=5pt, font=\fignodefont},
  stagetitle/.style={font=\small\bfseries\sffamily\color{ggunavy}},
  stageitem/.style={font=\small\sffamily, text width=2.2cm, align=center},
  arrow/.style={-stealth, ggunavy, very thick, >=latex, shorten >=2pt, shorten <=2pt, font=\fignodefont}]

% Stage 1: Input

\node[stage=lightyellow!50] (s1) at (0,0) {};

\node[stagetitle] at ([yshift=-4pt]s1.north) {1. Input};

\node[stageitem] at (s1.center) {Raw EEG\\256\,Hz, 22\,Ch};

\node[stageitem, font=\small\sffamily, text=gray] at ([yshift=4pt]s1.south) {MNE-Python\\EDF/BDF\\Montage};

% Stage 2: Preprocessing

\node[stage=lightblue!50] (s2) at (3.9,0) {};

\node[stagetitle] at ([yshift=-4pt]s2.north) {2. Preprocess};

\node[stageitem] at (s2.center) {Filtered\\Cleaned\\Re-referenced};

\node[stageitem, font=\small\sffamily, text=gray] at ([yshift=4pt]s2.south) {Notch/BPF\\Bandpass\\ICA};

% Stage 3: Feature Engineering

\node[stage=lightgreen!50] (s3) at (7.8,0) {};

\node[stagetitle] at ([yshift=-4pt]s3.north) {3. Feature Eng.};

\node[stageitem] at (s3.center) {Spectral\\Connectivity\\Bio-markers};

\node[stageitem, font=\small\sffamily, text=gray] at ([yshift=4pt]s3.south) {PSD\\CWT\\Connectivity};

% Stage 4: Modeling

\node[stage=lightpurple!50] (s4) at (11.7,0) {};

\node[stagetitle] at ([yshift=-4pt]s4.north) {4. Modeling};

\node[stageitem] at (s4.center) {CNN/Trans\\Riemannian\\Ensemble};

\node[stageitem, font=\small\sffamily, text=gray] at ([yshift=4pt]s4.south) {EEGNet\\TDN\\ViT};

% Stage 5: Inference

\node[stage=lightorange!50] (s5) at (15.6,0) {};

\node[stagetitle] at ([yshift=-4pt]s5.north) {5. Inference};

\node[stageitem] at (s5.center) {Prediction\\Probability\\Class};

\node[stageitem, font=\small\sffamily, text=gray] at ([yshift=4pt]s5.south) {Softmax\\Calibration\\Threshold};

% Stage 6: RAG Explainability

\node[stage=lightcoral!30] (s6) at (19.5,0) {};

\node[stagetitle] at ([yshift=-4pt]s6.north) {6. RAG};

\node[stageitem] at (s6.center) {Retrieve\\Synthesis\\Evidence-grounded};

\node[stageitem, font=\small\sffamily, text=gray] at ([yshift=4pt]s6.south) {ChromaDB\\llama3.2:3b\\SHAP+Citations};

% Stage 7: Output

\node[stage=lightpink!50] (s7) at (23.4,0) {};

\node[stagetitle] at ([yshift=-4pt]s7.north) {7. Output};

\node[stageitem] at (s7.center) {Screening\\Explanation\\Report};

\node[stageitem, font=\small\sffamily, text=gray] at ([yshift=4pt]s7.south) {B2B/B2C\\PDF report\\Guardrail};

% Arrows between stages
\draw[arrow] (s1.east) -- (s2.west);
\draw[arrow] (s2.east) -- (s3.west);
\draw[arrow] (s3.east) -- (s4.west);
\draw[arrow] (s4.east) -- (s5.west);
\draw[arrow] (s5.east) -- (s6.west);
\draw[arrow] (s6.east) -- (s7.west);

\end{tikzpicture}%
}% end resizebox
\caption[ML Pipeline Architecture (7 Stages)]{End-to-End ML Pipeline Architecture with Seven Sequential Stages Including RAG Explainability. Each stage enforces input/output contracts and quality gates before passing data to the next stage.}

\vspace{0.5em}
\noindent\textit{Note.} EEG = electroencephalography; Hz = hertz; Ch = channels; BPF = bandpass filter; ICA = independent component analysis; PSD = power spectral density; CWT = continuous wavelet transform; CNN = convolutional neural network; Trans = Transformer; TDN = temporal difference network; ViT = Vision Transformer; SHAP = SHapley Additive exPlanations; RAG = retrieval-augmented generation; FAISS = Facebook AI Similarity Search; LLM = large language model. Stage~6 (RAG) retrieves evidence from 50{,}000+ PubMed abstracts via ChromaDB/FAISS and synthesises clinical explanations using llama3.2:3b. B2B = business-to-business; B2C = business-to-consumer.
\end{figure}

the corresponding figure depicts the data flow through six processing stages with complete lineage.


\noindent\textbf{Data flow diagram showing six processing stages.}

\needspace{24\baselineskip}
\begin{figure}[!htbp]
\centering
\resizebox{0.97\textwidth}{!}{%
\begin{tikzpicture}[
  every node/.style={font=\fignodefont},
  process/.style={draw=ggunavy, thick, circle, fill=lightblue!40,
    minimum size=1.6cm, font=\small\sffamily\bfseries, align=center},
  entity/.style={draw=ggunavy, thick, fill=lightpink!50, rounded corners=2pt,
    minimum width=2cm, minimum height=0.9cm, font=\small\sffamily\bfseries, align=center},
  datastore/.style={draw=ggunavy, thick, fill=lightyellow!50,
    minimum width=2.2cm, minimum height=1.0cm, font=\small\sffamily\bfseries, align=center},
  flowlabel/.style={font=\small\sffamily, fill=white, inner sep=1pt},
  arrow/.style={-stealth, ggunavy, thick, shorten >=2pt, shorten <=2pt, font=\fignodefont}]

% External entities

\node[entity] (eeg) at (-1, 5) {EEG Device};

\node[entity] (clin) at (15, 5) {Clinician};

% Processes (circles)

\node[process] (p1) at (2, 3.5) {1.0\\Receive};

\node[process] (p2) at (5, 3.5) {2.0\\Preprocess};

\node[process] (p3) at (8, 3.5) {3.0\\Extract};

\node[process] (p4) at (11, 3.5) {4.0\\Predict};

\node[process, fill=lightpurple!40] (p5) at (11, 0.4) {5.0\\RAG\\Explain};

\node[process] (p6) at (5, 0.4) {6.0\\Store};

% Data stores

\node[datastore] (d1) at (2, 0.4) {D1: Raw EEG};

\node[datastore] (d2) at (15, 0.4) {D2: Results};

\node[datastore] (db) at (8, -1.2) {Patient DB};

\node[datastore, fill=lightpurple!20] (kb) at (14, -1.2) {D3: Knowledge Base\\(50K abstracts)};

% Arrows: External to Process
\draw[arrow] (eeg) -- node[flowlabel, above, sloped] {EEG Signal} (p1);
\draw[arrow] (p4) -- node[flowlabel, above, sloped] {Diagnosis} (clin);

% Arrows: Process to Process
\draw[arrow] (p1) -- node[flowlabel, above] {Raw Data} (p2);
\draw[arrow] (p2) -- node[flowlabel, above] {Clean Signal} (p3);
\draw[arrow] (p3) -- node[flowlabel, above] {Features} (p4);
\draw[arrow] (p4) -- node[flowlabel, right] {Prediction} (p5);
\draw[arrow] (p5) -- node[flowlabel, above] {Explanation} (p6);

% Arrows: Process to Data Store
\draw[arrow] (p1) -- node[flowlabel, left] {Store Raw} (d1);
\draw[arrow] (p5) -- node[flowlabel, above, sloped] {Results} (d2);
\draw[arrow] (p6) -- node[flowlabel, left] {Archive} (db);

% Arrows: Data Store to Process (read-back)
\draw[arrow, dashed] (d1) -- node[flowlabel, right] {Retrieve} (p2.south west);
\draw[arrow, dashed] (db) -- node[flowlabel, right] {Patient Info} (p5.south);
\draw[arrow, dashed, purple!70!black] (kb) -- node[flowlabel, above, sloped] {FAISS Top-5} (p5.south east);

\end{tikzpicture}%
}% end resizebox
\caption[Data Flow Diagram: Six Processing Stages]{Data flow diagram showing six processing stages (circles), two external entities (rounded rectangles), and three data stores (rectangles). Solid arrows indicate primary data flow; dashed arrows indicate secondary look-up queries.}

\vspace{0.5em}
\noindent\textit{Note.} EEG = electroencephalography; DB = database; D1, D2 = data stores for raw signals and computed results respectively. The six processes map to the pipeline stages in the corresponding figure. RAG-based explanation (Process 5.0) retrieves context from the Patient DB before generating clinician-facing narratives.
\end{figure}

%% ============================================================================

%% ======================================================================
%% C-S4. End-to-End IoT-Enabled EEG System Architecture and Evaluation Methodology
%% Moved from Chapter 3 — 2026-05-12 (round 2)
%% ======================================================================
\addcontentsline{toc}{section}{C-S4. End-to-End IoT EEG Architecture and Deployment} % [TOC-appendix-section-insertion]
\section*{C-S4. End-to-End IoT EEG Architecture and Deployment}
\addcontentsline{toc}{section}{C-S4. End-to-End IoT-Enabled EEG System Architecture and Evaluation Methodology}
\label{sec:appc_iot_arch}

The system architecture evaluated in this study was designed as an end-to-end decision-support pipeline linking patient-side data capture, wearable EEG acquisition, backend signal analysis, explainable model inference, and hospital-facing workflow integration. The objective of this architecture was not merely to produce an isolated classification output, but to support a clinically reviewable and governance-aware assessment process suitable for North American evaluation contexts.

At the patient-side layer, onboarding was conducted through a mobile or web-based interface supporting identity capture, consent collection, symptom reporting, behavioural and psychological input, and device session initiation. EEG data were acquired through an IoT-enabled wearable device such as Emotiv, transmitting signals either directly to an application layer or through a gateway service responsible for secure ingestion and session handling.

Following acquisition, EEG data passed through a preprocessing pipeline consisting of signal-quality checking, artefact handling, band-pass filtering, segmentation, and normalisation. Signal transformation was performed using frequency-domain and time-frequency techniques, including Fourier-based analysis, power spectral density estimation, and stationary wavelet transform procedures. Feature extraction, selection, and optional one-dimensional to two-dimensional representation construction followed for computer-vision-based modelling.

The modelling layer generated classification outputs, confidence estimates, and decision-support signals. An explainability layer including feature-attribution techniques and retrieval-augmented generation transformed model output into grounded explanatory summaries for clinicians and patients. A model-context or orchestration layer coordinated interactions among the inference service, retrieval components, logging systems, and external integrations.

The evaluation methodology explicitly incorporated human-in-the-loop review. System outputs were treated as decision-support recommendations rather than autonomous diagnoses. Governance controls including encryption, access control, audit logging, data minimisation, and retention management were included as architectural requirements. This end-to-end methodology allowed evaluation of technical performance, explainability, workflow fit, governance readiness, and adoption potential within one integrated architecture.

\iffalse %% Moved to Appendix C: C4 Container Architecture (engineering detail)
\noindent\textbf{C4 Container Diagram: Mobile App, Gateway,.}

\needspace{24\baselineskip}
\begin{figure}[!htbp]
\centering
\resizebox{0.97\textwidth}{!}{%
\begin{tikzpicture}[
  every node/.style={font=\fignodefont},
  node distance=0.8cm and 1.2cm,
    container/.style={rectangle, rounded corners, draw=ch3header, fill=ch3light, font=\small, text width=2.5cm, align=center, minimum height=1.2cm},
    external/.style={rectangle, rounded corners, draw=ch3header!50, fill=white, font=\small, text width=2.5cm, align=center, minimum height=1cm, dashed},
    arrow/.style={-{Stealth[length=4mm, width=3mm]}, very thick, ggunavy},
    label/.style={font=\scriptsize\itshape, ch3header}
]
% External actors
\node[external] (patient) {Patient /\\Parent};
\node[external, right=4cm of patient] (clinician) {Clinician};
\node[external, right=4cm of clinician] (ehr) {Hospital\\EHR / HIS};

% Containers row 1
\node[container, below=1.5cm of patient] (app) {\textbf{Mobile App}\\Onboarding,\\Consent, Forms};
\node[container, right=of app] (device) {\textbf{Emotiv/IoT}\\EEG Capture};
\node[container, right=of device] (gateway) {\textbf{Gateway}\\API, Auth,\\Ingestion};
\node[container, right=of gateway] (backend) {\textbf{Backend}\\Preprocessing,\\Features, Model};

% Containers row 2
\node[container, below=1.5cm of app, xshift=1.5cm] (rag) {\textbf{RAG Engine}\\Retrieval,\\Explanation};
\node[container, right=of rag] (mcp) {\textbf{MCP / Orchestrator}\\Tool Coordination,\\Logging};
\node[container, right=of mcp] (dashboard) {\textbf{Dashboard}\\Clinician View,\\Review Queue};
\node[container, right=of dashboard] (monitor) {\textbf{Monitoring}\\Drift, KPIs,\\Alerts};

% Arrows
\draw[arrow] (patient) -- (app);
\draw[arrow] (app) -- (gateway);
\draw[arrow] (device) -- (gateway);
\draw[arrow] (gateway) -- (backend);
\draw[arrow] (backend) -- (mcp);
\draw[arrow] (mcp) -- (rag);
\draw[arrow] (mcp) -- (dashboard);
\draw[arrow] (mcp) -- (monitor);
\draw[arrow] (dashboard) -- (clinician);
\draw[arrow] (dashboard) -- (ehr);

% Labels
\node[label, above=0.1cm of app] {Patient-Side};
\node[label, above=0.1cm of backend] {Server-Side};
\node[label, above=0.1cm of dashboard] {Clinical-Side};

\end{tikzpicture}%
}
\caption[C4 Container Diagram: Mobile App, Gateway, Backend, RAG,]{C4 Container Diagram: Mobile App, Gateway, Backend, RAG, MCP, and Hospital Integration}
\label{fig:c4_container}
\end{figure}
\fi
\noindent Full details are provided in Appendix~C, Section~\ref{sec:appendix_ch3_c4_container}.


\iffalse %% Moved to Appendix C: End-to-End Architecture Table (engineering detail)
\clearpage
{%
\tablefontsize

\noindent Table~\ref{tab:e2e_architecture} presents end-to-End System Architecture: 15-Layer Pipeline from Intake to Decision.

\begin{xltabular}
{\textwidth}{@{} r >{\raggedright\arraybackslash}p{2cm} L L L @{}}
\caption[End-to-End System Architecture: 15-Layer Pipeline from]{End-to-End System Architecture: 15-Layer Pipeline from Intake to Decision} \label{tab:e2e_architecture} \\
\toprule
\thead{\#} & \thead{Layer} & \thead{What Happens} & \thead{Key Controls} & \thead{Output} \\
\midrule
\endfirsthead
\endhead
\bottomrule
\endlastfoot
1 & Patient onboarding & Mobile/web intake, consent, symptom/behavioural capture & Identity, consent, RBAC, minimum data & Registered case \\
2 & Device session & Emotiv/IoT EEG pairs to app or gateway & Device auth, session token, signal-quality gate & Active recording session \\
3 & Data capture & EEG + forms/symptoms/behavioural inputs collected & Encrypted transport, PII separation & Raw EEG + primary data \\
4 & Gateway / ingestion & Backend receives stream/files & TLS, API auth, audit logging, schema validation & Accepted intake payload \\
5 & Preprocessing & Artefact handling, band-pass filter, epoching, normalisation & Reproducible rules, quality thresholds & Clean EEG segments \\
6 & Transformation & FFT, PSD/SPDD, SWT, feature extraction, 1D-to-2D & Versioned pipeline, leakage-safe processing & Model-ready inputs \\
7 & Model inference & ML/DL/CV model scores class/risk/confidence & Calibrated outputs, uncertainty handling, monitored version & Prediction + confidence \\
8 & Explainability & SHAP/saliency/attention computed & Method aligned to model type & Machine-level explanation \\
9 & RAG layer & Retrieve approved references, generate summaries & Controlled corpus, citation traceability & Clinician/patient explanations \\
10 & Orchestration & MCP coordinates model, retrieval, alerts, logs & Tool boundaries, access checks, traceability & Orchestrated workflow \\
11 & Decision support & Thresholding, review queue, escalation logic & Abstain/escalate rules & Review/monitor/escalate \\
12 & Hospital integration & Dashboard/EHR/HIS integration & Mapped identifiers, access control, audit trail & Clinician case summary \\
13 & Safety escalation & High-risk triggers SOS/escalation & Documented criteria, override path, notification log & Urgent review path \\
14 & Monitoring & Drift, latency, errors, trust/use metrics & KPI dashboard, alerts, version tracking & Operational governance \\
15 & Decision & Organisation decides evidence-maturity interpretation & Technical + trust + governance + illustrative financial scenario evidence & Evaluation recommendation \\
\end{xltabular}
}%
\fi
\noindent Full details are provided in Appendix~C, Section~\ref{sec:appendix_ch3_e2e_architecture}.


\noindent\textit{Note.} Layers~1--6 handle data acquisition and preparation; layers~7--10 perform inference and explanation; layers~11--15 support clinical decision-making and governance. All layers log to an immutable audit trail.

\clearpage
{%
\tablefontsize

\noindent Table~\ref{tab:na_compliance} lists the North America compliance checklist mapping each privacy / security requirement to the corresponding implementation in this study.

\paragraph{North America Compliance Checklist.}




\needspace{20\baselineskip}
\begin{xltabular}{\textwidth}{@{} >{\raggedright\arraybackslash}p{2.2cm} L L @{}}
\caption[North America Compliance Checklist for Clinical AI]{North America Compliance Checklist for Clinical AI Deployment} \label{tab:na_compliance} \\
\toprule
\thead{Area} & \thead{U.S. / Canada Requirement} & \thead{Implementation} \\
\midrule
\endfirsthead
\endhead
\bottomrule
\endlastfoot
Privacy law mapping & Determine HIPAA or Canadian provincial health privacy rules & Jurisdiction map by customer/site \\
Consent & Collect consent, document lawful purpose & Explicit onboarding consent flow, versioned notices \\
Data minimisation & Collect/use only data needed & Role-scoped views, separate intake fields \\
Encryption & Protect data in transit and at rest & TLS in transit, encryption at rest, managed keys \\
Access control & Limit access to authorised users & RBAC, MFA, least privilege \\
Audit controls & Record access, changes, overrides, disclosures & Immutable audit logs, review workflow \\
Integrity / traceability & Preserve data and model traceability & Model versioning, feature pipeline versioning \\
Breach handling & Maintain breach records and notification procedures & Incident plan, breach register \\
Human-in-the-loop & Avoid unsafe autonomous clinical use & Clinician review queue, override capability \\
CDS / SaMD classification & Decide whether software is non-device CDS or regulated & Regulatory assessment, intended-use statement \\
AI/ML change control & Control model updates post-deployment & Approval workflow, retraining governance \\
Monitoring & Monitor drift, latency, failures, quality & KPI dashboard, threshold alerts, rollback plan \\
Data retention / deletion & Define retention and secure deletion rules & Retention schedule, deletion jobs \\
Cross-border data & Check data residency requirements & Data residency plan, vendor review \\
Third-party / vendor risk & Cloud, app, gateway vendor compliance & Vendor due diligence, contracts \\
\end{xltabular}
}%

\noindent\textit{Note.} HIPAA = Health Insurance Portability and Accountability Act; CDS = Clinical Decision Support; SaMD = Software as a Medical Device; RBAC = Role-Based Access Control; MFA = Multi-Factor Authentication; TLS = Transport Layer Security.

\iffalse %% Moved to Appendix C: B2B Hospital Architecture (engineering detail)
\noindent\textbf{North America B2B Hospital Architecture: Stages and Requirem.}

\noindent Table~\ref{tab:b2b_hospital_architecture} presents north America B2B Hospital Architecture: Stages and Requirements.

{\tablestripes\tablefontsize
\needspace{20\baselineskip}
\begin{xltabular}{\textwidth}{@{} r L L @{}}
\caption[North America B2B Hospital Architecture]{North America B2B Hospital Architecture: Stages and Requirements}\label{tab:b2b_hospital_architecture} \\
\toprule
\thead{\#} & \thead{Hospital Architecture Stage} & \thead{Why This Is the Strongest North America Fit} \\
\midrule
\endfirsthead
\endhead
\bottomrule
\endlastfoot
1 & Patient onboarded by hospital staff through secure web intake & Reduces consumer misuse risk \\
2 & Emotiv EEG captured in supervised clinical workflow & Improves data quality and chain of custody \\
3 & Gateway sends data to hospital-controlled secure backend & Supports privacy and vendor governance \\
4 & Backend runs preprocessing, FFT/PSD/SWT, feature extraction, model inference & Keeps technical pipeline controlled \\
5 & Calibrated output with uncertainty and review logic & Safer for CDS-style use \\
6 & SHAP + RAG generates clinician-facing explanation & Improves interpretability \\
7 & Result goes only to clinician dashboard/EHR, not direct patient diagnosis & Aligns with safer CDS evaluation \\
8 & Clinician reviews, accepts, overrides, or escalates & Preserves human-in-the-loop responsibility \\
9 & Audit trail, access logs, model/version traceability maintained & Supports HIPAA/provincial compliance \\
10 & Monitoring tracks drift, latency, override rate, false negatives & Supports pilot governance \\
\end{xltabular}}
\vspace{2pt}
\noindent\textit{Note.} CDS = Clinical Decision Support; EHR = Electronic Health Record; FFT = Fast Fourier Transform; PSD = Power Spectral Density; SWT = Stationary Wavelet Transform; SHAP = SHapley Additive exPlanations; RAG = Retrieval-Augmented Generation.
\fi
\noindent Full details are provided in Appendix~C, Section~\ref{sec:appendix_ch3_b2b_hospital_architecture}.




\noindent\textit{\color{currentaccent}Part C: the following sections detail the B2B hospital assessment workflow, patient questionnaire framework, and clinical forms.}


%% ======================================================================
%% C-S9. Reference and Governance Integration — Full Per-Framework Detail
%% Condensation companion to Chapter 3 tab:ch3_partd_reference (4-row summary).
%% Chapter table = decision-level group summary; this appendix table = full
%% per-framework breakdown across all four governance/compliance/quality/
%% ethical dimensions. Migrated 2026-05-13 per main-chapter table-row
%% reduction discipline (max 2--4 rows in body, full enumeration in appendix).
%% ======================================================================
\addcontentsline{toc}{section}{C-S9. Reference and Governance Integration --- Full Per-Framework Breakdown} % [TOC-appendix-section-insertion]
\section*{C-S9. Reference and Governance Integration --- Full Per-Framework Breakdown}
\addcontentsline{toc}{section}{C-S9. Reference and Governance Integration}
\label{sec:appc_partd_reference_full}

\noindent This section provides the full per-framework breakdown summarised in Chapter~3, Table~\ref{tab:ch3_partd_reference}. Each of the four dimension groups (Governance, Compliance, Quality \& Architecture, Ethical / Responsible AI) is decomposed into its constituent frameworks and standards, with the application to Parts A, B, and C documented per row.

\noindent Table~\ref{tab:appc_partd_reference_full} summarises reference and governance integration for appendix review.

{\tablestripes\tablefontsize

\noindent Table~\ref{tab:appc_partd_reference_full} below presents Reference and Governance Integration: Full Per-Framework Breakdown, laying out each row in structured form to help examiners trace the source, applicability, and downstream use at a glance. % [auto-intro-strict inserted]
\paragraph{Reference And Governance Integration Full.} % [auto-heading inserted]
\noindent Table~\ref{tab:appc_partd_reference_full} below presents Reference and Governance Integration: Full Per-Framework Breakdown, laying out each row in structured form to help examiners trace the source, applicability, and downstream use at a glance. % [auto-intro-after-heading]



\needspace{20\baselineskip}
\begin{xltabular}{\textwidth}{@{} >{\raggedright\arraybackslash}p{3.2cm} p{5.2cm} p{6.0cm} @{}}
\caption[Reference and Governance Integration: Full Per-Framework Breakdown]{Reference and Governance Integration: full per-framework breakdown across Parts A--C. This is the supporting evidence behind the 4-row decision-level summary in Table~\ref{tab:ch3_partd_reference}; the chapter table groups these into Governance, Compliance, Quality \& Architecture, and Ethical / Responsible AI.}\label{tab:appc_partd_reference_full} \\
\toprule
\thead{Dimension} & \thead{Framework / Standard} & \thead{Application to Parts A--C} \\
\midrule
\endfirsthead
\endhead
\bottomrule
\endlastfoot
\multicolumn{3}{@{}l}{\cellcolor{currentheader!15}\textbf{Governance Frameworks}} \\
\addlinespace
RGAIG 5-Layer & L1--L5 governance architecture & Applied to all pipeline stages across A, B, C \\
NIST AI RMF & Govern, Map, Measure, Manage & Risk assessment for EEG pipeline and survey data \\
ISO/IEC 42001 & AI management system & Quality management for model lifecycle \\
CMM Maturity & 5-level capability model & Maturity assessment for evidence-maturity interpretation \\
\midrule
\multicolumn{3}{@{}l}{\cellcolor{currentheader!15}\textbf{Compliance Requirements}} \\
\addlinespace
HIPAA & Protected health information & B2B hospital EEG data (Part C) \\
GDPR / PIPEDA & Consumer data protection & B2C patient data (Parts A, B) \\
FDA SaMD & Software as medical device & Classification pipeline (Parts B, C) \\
EU AI Act & High-risk AI classification & All AI components (Parts A, B, C) \\
IRB / Ethics & Research ethics approval & Survey instruments (Part A), EEG datasets (Part B) \\
\midrule
\multicolumn{3}{@{}l}{\cellcolor{currentheader!15}\textbf{Quality and Architecture}} \\
\addlinespace
SMART Goals & Specific, Measurable, Achievable, Relevant, Time-bound & Each Part has SMART table \\
CRISP-DM & Data mining lifecycle & EEG pipeline stages (Part B) \\
ISO 25010 & Software quality model & System reliability, usability, security \\
HL7 FHIR R4 & Healthcare interoperability & Hospital integration (Part C) \\
525-Module Taxonomy & Quality assurance modules & Governance validation across A, B, C \\
\midrule
\multicolumn{3}{@{}l}{\cellcolor{currentheader!15}\textbf{Ethical and Responsible AI}} \\
\addlinespace
Explainability & SHAP + RAG narratives & Model outputs (Part B), clinician reports (Part C) \\
Fairness & Demographic subgroup audit & $\pm 2\%$ accuracy variance across subgroups \\
Transparency & Audit trail, decision logging & All pipeline stages \\
Accountability & RACI matrix, named owners & Every governance artefact \\
Privacy & De-identification, encryption & All patient data (Parts A, B, C) \\
Trust & Confidence calibration, human-in-the-loop & Clinician review (Part C), patient alerts (Part A) \\
\end{xltabular}}
\vspace{0.4em}
\noindent\textit{Note.} Twenty rows aggregated into four chapter-level group summaries in Table~\ref{tab:ch3_partd_reference}. The grouping logic (Governance / Compliance / Quality / Ethical) follows the AI governance maturity literature (NIST AI RMF, ISO/IEC 42001, EU AI Act). Cross-cutting application is enforced via the 525-module quality taxonomy (Appendix~G). HITL = human-in-the-loop; RACI = Responsible, Accountable, Consulted, Informed.

%% ======================================================================
%% C-S8. Research Survey Instruments
%% Moved from Chapter 3 — 2026-05-12 (round 2)
%% ======================================================================
\addcontentsline{toc}{section}{C-S8. Research Survey Instruments} % [TOC-appendix-section-insertion]
\section*{C-S8. Research Survey Instruments}
\addcontentsline{toc}{section}{C-S8. Research Survey Instruments}
\label{sec:appc_survey_instruments}

The RGAIG framework evaluation employs three validated survey instruments targeting distinct stakeholder groups: clinicians (B2B), consumers (B2C), and domain experts. Each instrument uses a 5-point Likert scale (1 = Strongly Disagree to 5 = Strongly Agree) unless otherwise specified. the corresponding table--N/A present the complete item sets with theoretical grounding.

Trust in AI diagnostics is measured through a six-item construct adapted from Lee and See~\cite{lee2004trust} and Madsen and Gregor~\cite{madsen2000measuring}, achieving internal consistency of $\alpha = 0.91$ and convergent validity AVE $= 0.69$ in pilot testing. Items assess perceived reliability, predictability, and faith in the system's diagnostic outputs. These six items (C5, C6, C11, C12, C13, C14 in the corresponding table) collectively operationalise the trust construct that mediates between governance antecedents and adoption intention in the structural model (SH1--SH7). High internal consistency confirms that clinicians interpret trust as a coherent latent variable rather than a set of disparate perceptions---a prerequisite for the mediation analysis reported in Chapter~\ref{chap:analysis}.

%% Clinician Survey (B2B)

\noindent\textbf{Clinician Survey Instrument (B2B): AI Trust,.}

\needspace{24\baselineskip}
\noindent The table below presents the clinician Survey Instrument (B2B): AI Trust, Adoption, and Governance.

\begin{table}[p]
\tablestripes
\tablefontsize
\caption{Clinician Survey Instrument (B2B): AI Trust, Adoption, and Governance}
\begin{tabularx}{\textwidth}{@{} >{\raggedright\arraybackslash}p{0.5cm} >{\raggedright\arraybackslash}p{2.2cm} L >{\raggedright\arraybackslash}p{1.5cm} @{}}
\toprule
\thead{\#} & \thead{Construct} & \thead{Survey Item} & \thead{Theory} \\
\midrule
C1 & Perceived Usefulness & The AI-assisted EEG interpretation system would improve the accuracy of my diagnostic decisions & TAM \\
\addlinespace
C2 & Perceived Usefulness & The system would reduce the time I spend on routine EEG analysis & TAM \\
\addlinespace
C3 & Ease of Use & I would find the AI system easy to integrate into my current clinical workflow & TAM \\
\addlinespace
C4 & Ease of Use & The system interface is intuitive and requires minimal training & TAM \\
\addlinespace
C5 & Cognitive Trust & I trust the AI system's outputs because they are based on validated clinical evidence & Trust Theory \\
\addlinespace
C6 & Affective Trust & I feel confident relying on the AI system for patient screening recommendations & Trust Theory \\
\addlinespace
C7 & Explainability & The AI explanations (SHAP values, RAG citations) help me understand why a recommendation was made & XAI \\
\addlinespace
C8 & Explainability & The literature citations in the AI report are relevant and clinically meaningful & XAI \\
\addlinespace
C9 & Perceived Risk & I am concerned about legal liability if the AI system produces an incorrect recommendation & Risk Theory \\
\addlinespace
C10 & Perceived Risk & I worry that patients might be harmed by false positive or false negative AI results & Risk Theory \\
\addlinespace
C11 & Governance & I believe the AI system has adequate safety controls and governance mechanisms & RGAIG \\
\addlinespace
C12 & Governance & The audit trail and decision logging give me confidence in system accountability & RGAIG \\
\addlinespace
C13 & Human Control & I feel I retain appropriate control over final diagnostic decisions when using the AI system & HITL \\
\addlinespace
C14 & Human Control & The system appropriately escalates uncertain cases to human review & HITL \\
\addlinespace
C15 & Social Influence & My colleagues' positive experiences with AI tools would influence my adoption decision & UTAUT \\
\addlinespace
C16 & Adoption Intent & I would recommend this AI system for use in my department & UTAUT \\
\addlinespace
C17 & Adoption Intent & I intend to use this AI system regularly if it becomes available in my hospital & UTAUT \\
\addlinespace
C18 & Cognitive Load & The AI system reduces my cognitive burden during complex multi-patient assessments & CLT \\
\addlinespace
C19 & Data Quality & I trust the quality of EEG data captured by the wearable device for clinical screening & DQ \\
\addlinespace
C20 & Overall Satisfaction & Overall, the AI-assisted EEG system meets my expectations for clinical decision support & Overall \\
\bottomrule
\end{tabularx}
\vspace{0.5em}
\noindent\textit{Note.} 5-point Likert scale: 1 = Strongly Disagree, 2 = Disagree, 3 = Neutral, 4 = Agree, 5 = Strongly Agree. TAM = Technology Acceptance Model; UTAUT = Unified Theory of Acceptance and Use of Technology; XAI = Explainable AI; HITL = Human-in-the-Loop; CLT = Cognitive Load Theory; DQ = Data Quality; RGAIG = Responsible Generative AI Governance.
\end{table}

%% Consumer Survey (B2C)

\noindent\textbf{Consumer Survey Instrument (B2C): Usability, Trust, and Deci.}

\needspace{24\baselineskip}
\noindent The table below presents the consumer Survey Instrument (B2C): Usability, Trust, and Decision Support.

\begin{table}[p]
\tablestripes
\tablefontsize
\caption[Consumer Survey Instrument (B2C): Usability, Trust, and]{Consumer Survey Instrument (B2C): Usability, Trust, and Decision Support}
\begin{tabularx}{\textwidth}{@{} >{\raggedright\arraybackslash}p{0.5cm} >{\raggedright\arraybackslash}p{2.2cm} L >{\raggedright\arraybackslash}p{1.5cm} @{}}
\toprule
\thead{\#} & \thead{Construct} & \thead{Survey Item} & \thead{Theory} \\
\midrule
U1 & Understanding & I can understand the health insights the app provides about my EEG data & Literacy \\
\addlinespace
U2 & Understanding & The app explains medical terms in language I can comprehend & Literacy \\
\addlinespace
U3 & Trust & I trust the app to provide accurate information about my brain health & Trust \\
\addlinespace
U4 & Trust & I believe the app uses my EEG data responsibly and securely & Trust \\
\addlinespace
U5 & Explainability & The app clearly explains why it generated a particular health alert & XAI \\
\addlinespace
U6 & Explainability & I understand the confidence level shown with each AI recommendation & XAI \\
\addlinespace
U7 & Decision Support & The app helps me decide when I should consult a healthcare professional & Decision \\
\addlinespace
U8 & Decision Support & The app provides clear next steps after showing me my results & Decision \\
\addlinespace
U9 & Privacy & I am aware of how my EEG data is stored and who can access it & Privacy \\
\addlinespace
U10 & Privacy & The app gives me control over my data sharing preferences & Privacy \\
\addlinespace
U11 & Alert Quality & The alerts I receive are relevant and not excessive (no alert fatigue) & Usability \\
\addlinespace
U12 & Alert Quality & The severity levels of alerts help me prioritise which ones need attention & Usability \\
\addlinespace
U13 & Personalisation & The app provides insights tailored to my individual health patterns & Personal \\
\addlinespace
U14 & Behavioural Safety & The app does not cause me unnecessary anxiety or stress about my health & Safety \\
\addlinespace
U15 & Behavioural Safety & I feel calm and informed (not panicked) after reading the app's results & Safety \\
\addlinespace
U16 & Engagement & I use the app regularly because it provides valuable health insights & Engage \\
\addlinespace
U17 & Device Usability & The wearable EEG headset is comfortable and easy to use at home & Device \\
\addlinespace
U18 & Healthcare Link & If the app recommends seeing a doctor, I know how to connect with one through the app & Escalation \\
\addlinespace
U19 & Adoption Intent & I would recommend this wearable EEG system to family and friends & Adoption \\
\addlinespace
U20 & Overall Satisfaction & Overall, the wearable EEG app helps me better understand and manage my brain health & Overall \\
\bottomrule
\end{tabularx}
\vspace{0.5em}
\noindent\textit{Note.} 5-point Likert scale: 1 = Strongly Disagree to 5 = Strongly Agree. Constructs mapped to: Health Literacy Theory, Trust Theory, XAI perception, Decision Support, Privacy Awareness, Usability (ISO 9241), Personalisation, Behavioural Safety Design, Engagement, and Technology Adoption.
\end{table}

%% Expert Validation Survey


\noindent\textbf{Expert Validation Survey: RGAIG Framework Evaluation.} The following table presents the detailed analysis.

\needspace{24\baselineskip}
\noindent The table below presents the expert Validation Survey: RGAIG Framework Evaluation.

\begin{table}[p]
\tablestripes
\tablefontsize
\caption{Expert Validation Survey: RGAIG Framework Evaluation}
\begin{tabularx}{\textwidth}{@{} >{\raggedright\arraybackslash}p{0.5cm} >{\raggedright\arraybackslash}p{2.2cm} L >{\raggedright\arraybackslash}p{1.5cm} @{}}
\toprule
\thead{\#} & \thead{Dimension} & \thead{Survey Item} & \thead{RGAIG Layer} \\
\midrule
E1 & Architecture & The five-tier RGAIG architecture is well-structured for ASD AI & L1--L5 \\
\addlinespace
E2 & Architecture & The modular design allows effective domain extension to new conditions & L1--L5 \\
\addlinespace
E3 & Explainability & The RAG-based explanation generation provides clinically meaningful outputs & L3 \\
\addlinespace
E4 & Explainability & The SHAP/LIME feature attribution adds value to clinical decision-making & L3 \\
\addlinespace
E5 & Safety & The six-gate guardrail system provides adequate protection against harmful outputs & L4 \\
\addlinespace
E6 & Safety & The confidence threshold and human-fallback mechanisms are appropriately designed & L4 \\
\addlinespace
E7 & Trust & The framework's trust-building mechanisms (audit trail, transparency) are sufficient for clinical adoption & L4 \\
\addlinespace
E8 & Governance & The 525-module quality taxonomy provides comprehensive coverage of AI governance & L4 \\
\addlinespace
E9 & Governance & The responsible AI compliance scoring demonstrates adequate governance maturity & L4 \\
\addlinespace
E10 & Accuracy & The classification accuracy (ASD ensemble composite) is acceptable for clinical screening & L2 \\
\addlinespace
E11 & Accuracy & The ASD-focused validation approach provides adequate evidence of generalisability & L2 \\
\addlinespace
E12 & Clinical Utility & The framework would be useful in real hospital settings for EEG-based screening & Clinical \\
\addlinespace
E13 & Clinical Utility & The RAG-generated reports are suitable for clinician review and patient communication & Clinical \\
\addlinespace
E14 & Scalability & The architecture can scale to handle enterprise-level patient volumes & Technical \\
\addlinespace
E15 & Regulatory & The framework addresses key regulatory requirements (FDA SaMD, EU AI Act, HIPAA) & Compliance \\
\addlinespace
E16 & Innovation & The RGAIG framework represents a meaningful contribution to healthcare AI governance & Novelty \\
\addlinespace
E17 & Completeness & The framework adequately addresses the B2B (hospital) and B2C (consumer) use cases & Scope \\
\addlinespace
E18 & Limitations & The framework honestly acknowledges its limitations and boundary conditions & Rigour \\
\bottomrule
\end{tabularx}
\vspace{0.5em}
\noindent\textit{Note.} 5-point Likert scale: 1 = Strongly Disagree to 5 = Strongly Agree. Expert panel: minimum 5 evaluators from clinical neurology, AI/ML engineering, healthcare governance, and regulatory affairs. Inter-rater reliability measured by Krippendorff's $\alpha$ (target $\geq$0.80). RGAIG layers: L1 = Data, L2 = Model, L3 = Explanation, L4 = Governance, L5 = Interface.
\end{table}

%===============================================================================

%% ======================================================================
%% C-S7. Clinical Assessment Forms
%% Moved from Chapter 3 — 2026-05-12 (round 2)
%% ======================================================================
\addcontentsline{toc}{section}{C-S7. Clinical Assessment Forms} % [TOC-appendix-section-insertion]
\section*{C-S7. Clinical Assessment Forms}
\addcontentsline{toc}{section}{C-S7. Clinical Assessment Forms}
\label{sec:appc_clinical_forms}

In addition to the research surveys, the RGAIG system uses structured clinical assessment forms for each target condition. These forms capture patient-specific data that contextualises the AI model's EEG analysis. the corresponding table--N/A present the complete assessment instruments.

%% ASD Clinical Assessment Form

\noindent\textbf{ASD Clinical Assessment Form: Structured Data.}

\needspace{24\baselineskip}
\noindent The table below presents the aSD Clinical Assessment Form: Structured Data Capture.

\begin{table}[p]
\tablestripes
\tablefontsize
\caption{ASD Clinical Assessment Form: Structured Data Capture}
\begin{tabularx}{\textwidth}{@{} >{\raggedright\arraybackslash}p{0.5cm} >{\raggedright\arraybackslash}p{2cm} L >{\raggedright\arraybackslash}p{2cm} @{}}
\toprule
\thead{\#} & \thead{Domain} & \thead{Assessment Item} & \thead{Response Type} \\
\midrule
A1 & Demographics & Patient age, sex, handedness, developmental history & Numeric/Select \\
\addlinespace
A2 & Social & Eye contact frequency during interaction (none / occasional / frequent) & 3-point scale \\
\addlinespace
A3 & Social & Response to name when called (no response / delayed / immediate) & 3-point scale \\
\addlinespace
A4 & Social & Interest in peer interaction (avoidant / passive / active) & 3-point scale \\
\addlinespace
A5 & Communication & Speech development stage (non-verbal / single words / phrases / fluent) & 4-point scale \\
\addlinespace
A6 & Communication & Presence of echolalia (none / occasional / frequent) & 3-point scale \\
\addlinespace
A7 & Communication & Non-verbal communication use (gestures, pointing, showing) & Yes/No + notes \\
\addlinespace
A8 & Behaviour & Repetitive behaviours observed (list: hand flapping, rocking, spinning, lining up) & Checklist \\
\addlinespace
A9 & Behaviour & Insistence on sameness / routine disruption reaction & 5-point severity \\
\addlinespace
A10 & Sensory & Sensory sensitivities (auditory, visual, tactile, olfactory, gustatory) & Checklist \\
\addlinespace
A11 & Sensory & Sensory seeking behaviours observed & Yes/No + notes \\
\addlinespace
A12 & Cognitive & Estimated cognitive level (below / at / above age expectations) & 3-point scale \\
\addlinespace
A13 & Emotional & Emotional regulation difficulties (frequency and severity) & 5-point severity \\
\addlinespace
A14 & EEG Quality & EEG signal quality rating (poor / acceptable / good / excellent) & 4-point scale \\
\addlinespace
A15 & EEG Findings & Mu-rhythm suppression observed during social stimuli & Yes/No \\
\addlinespace
A16 & EEG Findings & Abnormal spectral power in frontal/temporal regions & Yes/No + band \\
\addlinespace
A17 & Prior Diagnosis & Previous ASD screening results (ADOS-2, ADI-R, M-CHAT scores) & Numeric \\
\addlinespace
A18 & Medications & Current medications and supplements & Free text \\
\addlinespace
A19 & Comorbidities & Co-occurring conditions (ADHD, anxiety, epilepsy, sleep disorders) & Checklist \\
\addlinespace
A20 & Clinician Notes & Overall clinical impression and recommendation & Free text \\
\bottomrule
\end{tabularx}
\vspace{0.5em}
\noindent\textit{Note.} This form is completed by the clinician during or after the EEG session. Items A14--A16 are populated automatically by the RGAIG system. ADOS-2 = Autism Diagnostic Observation Schedule; ADI-R = Autism Diagnostic Interview-Revised; M-CHAT = Modified Checklist for Autism in Toddlers.
\end{table}

%% PD Clinical Assessment Form

\needspace{24\baselineskip}
\iffalse % AUTO-SUPPRESS non-ASD
\noindent\textbf{Parkinson's Disease Clinical Assessment Form: Structured Dat.}

\needspace{24\baselineskip}
\noindent The table below presents the parkinson's Disease Clinical Assessment Form: Structured Data Capture.

\begin{table}[p]
\tablestripes
\tablefontsize
\caption{Parkinson's Disease Clinical Assessment Form: Structured Data Capture}
\begin{tabularx}{\textwidth}{@{} >{\raggedright\arraybackslash}p{0.5cm} >{\raggedright\arraybackslash}p{2cm} L >{\raggedright\arraybackslash}p{2cm} @{}}
\toprule
\thead{\#} & \thead{Domain} & \thead{Assessment Item} & \thead{Response Type} \\
\midrule
P1 & Demographics & Patient age, sex, age at symptom onset, family PD history & Numeric/Select \\
\addlinespace
P2 & Motor: Tremor & Resting tremor presence and affected limbs (L/R hand, L/R leg, jaw) & Checklist + side \\
\addlinespace
P3 & Motor: Bradykinesia & Slowness of movement severity (finger tapping, hand movements) & 5-point MDS \\
\addlinespace
P4 & Motor: Rigidity & Limb and axial rigidity on passive movement & 5-point MDS \\
\addlinespace
P5 & Motor: Gait & Gait pattern (normal / mildly impaired / shuffling / freezing / unable) & 5-point scale \\
\addlinespace
P6 & Motor: Balance & Postural stability (pull test result) & 5-point MDS \\
\addlinespace
P7 & Non-Motor: Sleep & Sleep quality, REM behaviour disorder presence, daytime sleepiness & 5-point + Yes/No \\
\addlinespace
P8 & Non-Motor: Cognitive & MoCA score or equivalent cognitive screening result & Numeric (0--30) \\
\addlinespace
P9 & Non-Motor: Mood & Depression (PHQ-9 score), anxiety (GAD-7 score), apathy & Numeric \\
\addlinespace
P10 & Non-Motor: Autonomic & Orthostatic hypotension, constipation, urinary symptoms & Checklist \\
\addlinespace
P11 & Speech & Voice volume, monotone quality, articulation clarity & 5-point scale \\
\addlinespace
P12 & Function & Hoehn and Yahr stage (1--5); Schwab and England ADL score & Stage + \% \\
\addlinespace
P13 & Medication & Current PD medications, dosages, timing, on/off fluctuations & Free text \\
\addlinespace
P14 & Medication & Levodopa equivalent daily dose (LEDD) calculation & Numeric (mg) \\
\addlinespace
P15 & EEG Quality & Signal quality rating; electrode impedance check & 4-point scale \\
\addlinespace
P16 & EEG Findings & Beta-band power changes (frontal/parietal regions) & Auto-populated \\
\addlinespace
P17 & EEG Findings & Cortical slowing pattern detected; event-related desynchronisation & Auto-populated \\
\addlinespace
P18 & Disease Duration & Years since diagnosis; years since first symptom & Numeric \\
\addlinespace
P19 & Comorbidities & Co-occurring conditions (dementia, depression, falls, osteoporosis) & Checklist \\
\addlinespace
P20 & Clinician Notes & MDS-UPDRS Part III total score; overall impression; follow-up plan & Numeric + text \\
\bottomrule
\end{tabularx}
\vspace{0.5em}
\noindent\textit{Note.} Motor items P2--P6 align with MDS-UPDRS Part III. Items P15--P17 are auto-populated by RGAIG. MDS = Movement Disorder Society; UPDRS = Unified Parkinson's Disease Rating Scale; MoCA = Montreal Cognitive Assessment; PHQ-9 = Patient Health Questionnaire; GAD-7 = Generalised Anxiety Disorder scale; LEDD = Levodopa Equivalent Daily Dose; ADL = Activities of Daily Living.
\end{table}
\fi % AUTO-SUPPRESS non-ASD

%% Epilepsy Clinical Assessment Form


\noindent\textbf{Epilepsy Clinical Assessment Form: Structured Data....} The following table presents the detailed analysis.

\needspace{24\baselineskip}
\noindent The table below presents the epilepsy Clinical Assessment Form: Structured Data Capture.

\begin{table}[p]
\tablestripes
\tablefontsize
\caption{Epilepsy Clinical Assessment Form: Structured Data Capture}
\begin{tabularx}{\textwidth}{@{} >{\raggedright\arraybackslash}p{0.5cm} >{\raggedright\arraybackslash}p{2cm} L >{\raggedright\arraybackslash}p{2cm} @{}}
\toprule
\thead{\#} & \thead{Domain} & \thead{Assessment Item} & \thead{Response Type} \\
\midrule
EP1 & Demographics & Patient age, sex, age at first seizure, family epilepsy history & Numeric/Select \\
\addlinespace
EP2 & Seizure Type & Classification: focal aware / focal impaired / tonic-clonic / absence / myoclonic / atonic & Select \\
\addlinespace
EP3 & Frequency & Seizure frequency (daily / weekly / monthly / yearly / seizure-free) & Select + count \\
\addlinespace
EP4 & Duration & Typical seizure duration (seconds / minutes) & Numeric \\
\addlinespace
EP5 & Triggers & Known triggers (sleep deprivation, stress, photosensitivity, alcohol, illness, missed medication) & Checklist \\
\addlinespace
EP6 & Aura & Aura presence and type (visual, auditory, olfactory, gustatory, emotional, d\'ej\`a vu) & Checklist \\
\addlinespace
EP7 & Ictal & Seizure description: motor activity, consciousness level, automatisms, lateralisation & Free text \\
\addlinespace
EP8 & Post-Ictal & Post-ictal symptoms: confusion (duration), headache, fatigue, Todd's paralysis & Checklist + time \\
\addlinespace
EP9 & Status & History of status epilepticus (prolonged seizure $>$5 minutes) & Yes/No + count \\
\addlinespace
EP10 & Medication & Current anti-seizure medications, doses, levels, adherence & Free text \\
\addlinespace
EP11 & Medication & Drug-resistant epilepsy: failed $\geq$2 appropriate medications at adequate doses & Yes/No \\
\addlinespace
EP12 & Cognitive & Cognitive function: memory, attention, processing speed, language & 5-point scale \\
\addlinespace
EP13 & Psychosocial & Impact on driving, employment, education, relationships & 5-point severity \\
\addlinespace
EP14 & Safety & Seizure-related injuries: falls, burns, drowning risk, nocturnal seizure precautions & Checklist \\
\addlinespace
EP15 & EEG Quality & Signal quality; recording duration; sleep/wake state during recording & 4-point + time \\
\addlinespace
EP16 & EEG Findings & Interictal epileptiform discharges: spike/sharp waves, location, frequency & Auto-populated \\
\addlinespace
EP17 & EEG Findings & Background activity: normal / mildly slow / moderately slow / severely abnormal & Auto-populated \\
\addlinespace
EP18 & EEG Findings & Seizure captured during recording (ictal pattern documented) & Yes/No \\
\addlinespace
EP19 & Imaging & MRI findings: normal / lesional (specify) / pending & Select + text \\
\addlinespace
EP20 & Clinician Notes & Epilepsy syndrome classification; treatment plan; surgery candidacy; follow-up & Free text \\
\bottomrule
\end{tabularx}
\vspace{0.5em}
\noindent\textit{Note.} Items EP16--EP18 are auto-populated by the RGAIG system's automated spike detection; classification accuracy is reported in Chapter~\ref{chap:analysis}. Seizure classification follows ILAE 2017 criteria. Drug-resistant epilepsy defined per Kwan et al.\ (2010). ILAE = International League Against Epilepsy.
\end{table}

%% ============================================================================
%% VARIABLE CLASSIFICATION AND MEASUREMENT MODEL
%% ============================================================================


%% ======================================================================
%% C-S6. Enterprise AI Pipeline Taxonomy and Lifecycle Integration
%% Moved from Chapter 3 — 2026-05-12 (round 2)
%% ======================================================================
\addcontentsline{toc}{section}{C-S6. Enterprise AI Pipeline Taxonomy} % [TOC-appendix-section-insertion]
\section*{C-S6. Enterprise AI Pipeline Taxonomy}
\addcontentsline{toc}{section}{C-S6. Enterprise AI Pipeline Taxonomy and Lifecycle Integration}
\label{sec:appc_ai_pipeline}

Production-grade biomedical AI involves more lifecycle stages than are necessary to explain the dissertation's core research method. The main chapter therefore notes only that ingestion, preprocessing, modeling, explanation, monitoring, and feedback exist as lifecycle categories; the full enterprise pipeline taxonomy is retained outside the core chapter.


Table~\ref{tab:pipeline_framework_mapping} maps each RGAIG framework layer to the pipeline categories that operationalise it, demonstrating that the architecture is not a conceptual abstraction but a deployable system with concrete implementation components.


\noindent\textbf{RGAIG Framework Layer to Pipeline Mapping.}

\needspace{24\baselineskip}
\begin{table}[H]
\tablestripes
\tablefontsize
\caption{RGAIG Framework Layer to Pipeline Mapping}
\begin{tabularx}{\textwidth}{@{} >{\raggedright\arraybackslash}p{2.5cm} L L @{}}
\toprule
\thead{Framework Layer} & \thead{Supporting Pipelines} & \thead{Governance Control} \\
\midrule
Data Layer (L1) & Ingestion, Quality, Validation & Schema validation; consent verification; lineage tracking \\
\addlinespace
Processing Layer (L2) & ETL, Training, Inference & Reproducibility controls; version pinning; seed logging \\
\addlinespace
Validation Layer (L3) & Quality, RAG, Guardrails & Cross-fold consistency; hallucination detection; citation accuracy \\
\addlinespace
Governance Layer (L4) & Policy, Audit, Security, IAM & RACI enforcement; access logs; retention policy; bias monitoring \\
\addlinespace
Decision Layer (L5) & Monitoring, Alerting, CI/CD, Incident & Drift detection; escalation triggers; rollback; KPI dashboards \\
\bottomrule
\end{tabularx}
\vspace{0.5em}
\noindent\textit{Note.} Each layer requires at least two pipeline categories to be operational. The governance layer (L4) spans the most pipelines because compliance controls must intercept data at every stage.
\end{table}

Table~\ref{tab:pipeline_risk} quantifies the operational risk when any pipeline category is absent, grounding the architecture's completeness claim in consequence analysis rather than aspirational design.


\noindent\textbf{Operational Risk of Missing Pipeline Categories.}

\needspace{24\baselineskip}
\begin{table}[!htbp]
\tablestripes
\tablefontsize
\caption{Operational Risk of Missing Pipeline Categories}
\begin{tabularx}{\textwidth}{@{} >{\raggedright\arraybackslash}p{2.1cm} X >{\raggedright\arraybackslash}p{1.6cm} X @{}}
\toprule
\thead{Pipeline Category} & \thead{Risk if Missing} & \thead{Severity} & \thead{Historical Precedent} \\
\midrule
Data Quality & Wrong predictions from corrupted input; undetected distribution shift & Critical & IBM Watson oncology: unsafe recommendations from weak training basis \\
\addlinespace
Governance & Regulatory non-compliance; undetectable bias; no audit trail & Critical & EU AI Act: risk management is mandatory for high-risk AI \\
\addlinespace
Monitoring & Silent model degradation; delayed incident response & High & Zillow iBuying: model drift caused major write-down \\
\addlinespace
Explainability & Black-box decisions; clinician rejection; regulatory failure & High & FDA SaMD: transparency required for clinical decision support \\
\addlinespace
Drift Detection & Accuracy erosion over time; stale predictions on new populations & High & COVID-19 diagnostic models degraded within months without retraining \\
\addlinespace
Security & Data breach; unauthorised access to patient brain data & Critical & HIPAA: major penalties for data breaches \\
\addlinespace
Incident Response & Extended outage; no fallback; patient safety risk & Medium & Clinical AI downtime without manual fallback delays diagnosis \\
\bottomrule
\end{tabularx}
\vspace{0.5em}
\noindent\textit{Note.} Severity ratings follow the RGAIG risk classification: Critical = safety or regulatory exposure; High = operational disruption or accuracy degradation; Medium = efficiency loss without immediate safety impact.
\end{table}


% ╔══════════════════════════════════════════════════════════════════════════════╗
% ║  PART D: REFERENCE, GOVERNANCE, AND CROSS-CUTTING INTEGRATION              ║
% ║  Frameworks, compliance, architecture, ethics, quality — applies to A+B+C  ║
% ╚══════════════════════════════════════════════════════════════════════════════╝
%%======NEW SECTIONS APPENDED 2026-05-13======

%% ======================================================================
%% C-S10. Part B SMART/RGAIG Alignment — Full Per-Row Breakdown
%% Condensation companion to Chapter 3 tab:ch3_partb_smart (4-row group
%% summary). Full 17-row 3-section decomposition migrated 2026-05-13
%% per main-chapter table-row reduction discipline.
%% ======================================================================
\addcontentsline{toc}{section}{C-S10. Part B SMART/RGAIG Alignment --- Full Per-Row Breakdown} % [TOC-appendix-section-insertion]
\section*{C-S10. Part B SMART/RGAIG Alignment --- Full Per-Row Breakdown}
\addcontentsline{toc}{section}{C-S10. Part B SMART/RGAIG Alignment}
\label{sec:appc_partb_smart_full}

\noindent This section provides the full per-row breakdown summarised in Chapter~3, Table~\ref{tab:ch3_partb_smart}. Each of the four chapter-level framework-lens rows (SMART, RGAIG 5 Layers, Quality Dimensions, Cross-cutting) is decomposed into its original 17 rows across three sections.

\noindent Table~\ref{tab:appc_partb_smart_full} summarises part b smart/rgaig framework alignment for appendix review.

{\tablestripes\tablefontsize

\noindent Table~\ref{tab:appc_partb_smart_full} below presents Part B SMART/RGAIG Alignment: Full Per-Row Breakdown, laying out each row in structured form to help examiners trace the source, applicability, and downstream use at a glance. % [auto-intro-strict inserted]
\paragraph{Part B Smart Rgaig Alignment.} % [auto-heading inserted]
\noindent Table~\ref{tab:appc_partb_smart_full} below presents Part B SMART/RGAIG Alignment: Full Per-Row Breakdown, laying out each row in structured form to help examiners trace the source, applicability, and downstream use at a glance. % [auto-intro-after-heading]



\needspace{20\baselineskip}
\begin{xltabular}{\textwidth}{@{} >{\raggedright\arraybackslash}p{2.0cm} L @{}}
\caption[Part B SMART/RGAIG Alignment: Full Per-Row Breakdown]{Part B SMART/RGAIG Framework Alignment: full per-row breakdown. Supporting evidence behind the 4-row decision-level summary in Table~\ref{tab:ch3_partb_smart}.}\label{tab:appc_partb_smart_full} \\
\toprule
\thead{Framework} & \thead{Part B Alignment} \\
\midrule
\endfirsthead
\endhead
\bottomrule
\endlastfoot
\multicolumn{2}{@{}l}{\cellcolor{currentheader!15}\textbf{SMART Goal Alignment}} \\
\addlinespace
Specific & Classify ASD EEG recordings with $\geq 85\%$ accuracy using ensemble classification on the KAU dataset \\
Measurable & Accuracy, AUC, F1, sensitivity, specificity, Cohen's $d$, bootstrap 95\% CI \\
Achievable & Literature baseline 78.35\% (Asthana et al., 2025); target 85\% is within demonstrated range \\
Relevant & Addresses G1 (no unified framework) and H1 (diagnostic accuracy threshold) \\
Time-bound & Data collection and model training completed within the doctoral research period (2025--2026) \\
\midrule
\multicolumn{2}{@{}l}{\cellcolor{currentheader!15}\textbf{RGAIG Five-Layer Governance Alignment}} \\
\addlinespace
L1: Data Foundation & EEG acquisition, quality gates, preprocessing, filtering, artefact removal \\
L2: Analytics Engine & Feature extraction, 1D$\to$2D conversion, ensemble classification, ablation study \\
L3: Trust \& Accountability & SHAP explainability, RAG narrative generation, confidence calibration \\
L4: Decision Orchestration & Threshold-based accept/reject, human-in-the-loop escalation, report generation \\
L5: Strategic Governance & 525-module quality taxonomy, compliance checklist, audit trail, ethical review \\
\midrule
\multicolumn{2}{@{}l}{\cellcolor{currentheader!15}\textbf{AI Quality Dimensions Addressed}} \\
\addlinespace
Explainability & SHAP feature attribution + RAG literature-grounded narratives \\
Responsibility & Bias audit (demographic subgroup accuracy $\pm 2\%$), fairness constraints \\
Performance & Classification accuracy, latency ($< 100$\,ms), throughput \\
Architecture & Modular 16-stage pipeline, HL7 FHIR integration, Emotiv IoT capture \\
Ethics & IRB exemption (public data), de-identification, informed consent protocol \\
Trust & Confidence scores, uncertainty quantification, human-fallback mechanism \\
Debuggability & Ablation study (component removal impact), error analysis (847 misclassified samples) \\
\end{xltabular}}
\vspace{0.4em}
\noindent\textit{Note.} Seventeen rows aggregated into four chapter-level framework-lens summaries in Table~\ref{tab:ch3_partb_smart}. Three original sections preserved: SMART criteria $\to$ chapter ``SMART'' row; RGAIG L1--L5 $\to$ chapter ``RGAIG'' row; 7 quality dimensions $\to$ chapter ``Quality Dimensions'' row; cross-cutting integration $\to$ chapter ``Cross-cutting'' row. HITL = human-in-the-loop; FHIR = Fast Healthcare Interoperability Resources.

%% ======================================================================
%% C-S11. Part A SMART/RGAIG Alignment — Full Per-Row Breakdown
%% Condensation companion to Chapter 3 tab:ch3_parta_smart (4-row group
%% summary). Full 15-row 3-section decomposition migrated 2026-05-13
%% per main-chapter table-row reduction discipline.
%% ======================================================================
\addcontentsline{toc}{section}{C-S11. Part A SMART/RGAIG Alignment --- Full Per-Row Breakdown} % [TOC-appendix-section-insertion]
\section*{C-S11. Part A SMART/RGAIG Alignment --- Full Per-Row Breakdown}
\addcontentsline{toc}{section}{C-S11. Part A SMART/RGAIG Alignment}
\label{sec:appc_parta_smart_full}

\noindent This section provides the full per-row breakdown summarised in Chapter~3, Table~\ref{tab:ch3_parta_smart}. The four chapter-level framework-lens rows (SMART, RGAIG 5 Layers, Quality Dimensions, Cross-cutting) decompose into the original 15-row 3-section enumeration for Part~A (B2C primary surveys).

\noindent Table~\ref{tab:appc_parta_smart_full} summarises part a smart/rgaig framework alignment for appendix review.

\clearpage
{\tablestripes\tablefontsize

\noindent Table~\ref{tab:appc_parta_smart_full} below presents Part A SMART/RGAIG Alignment: Full Per-Row Breakdown, laying out each row in structured form to help examiners trace the source, applicability, and downstream use at a glance. % [auto-intro-strict inserted]
\paragraph{Part A Smart Rgaig Alignment.} % [auto-heading inserted]
\noindent Table~\ref{tab:appc_parta_smart_full} below presents Part A SMART/RGAIG Alignment: Full Per-Row Breakdown, laying out each row in structured form to help examiners trace the source, applicability, and downstream use at a glance. % [auto-intro-after-heading]



\needspace{20\baselineskip}
\begin{xltabular}{\textwidth}{@{} >{\raggedright\arraybackslash}p{2.0cm} L @{}}
\caption[Part A SMART/RGAIG Alignment: Full Per-Row Breakdown]{Part A SMART/RGAIG Framework Alignment: full per-row breakdown. Supporting evidence behind the 4-row decision-level summary in Table~\ref{tab:ch3_parta_smart}.}\label{tab:appc_parta_smart_full} \\
\toprule
\thead{Framework} & \thead{Part A Alignment} \\
\midrule
\endfirsthead
\endhead
\bottomrule
\endlastfoot
\multicolumn{2}{@{}l}{\cellcolor{currentheader!15}\textbf{SMART Goal Alignment}} \\
\addlinespace
Specific & Measure clinician trust ($\geq 3.5/5$) and governance acceptance via four ASD-focused surveys (S1--S4) \\
Measurable & Likert scores (1--5), Cronbach's $\alpha \geq 0.70$, pre--post $\Delta$ with Cohen's $d$, SEM path $\beta$ \\
Achievable & Pilot testing achieved $\alpha = 0.91$ for trust construct; expert panel ($n = 12$) already recruited \\
Relevant & Addresses H4 (clinician trust) and H5 (governance convergence) directly \\
Time-bound & Survey administration and analysis within doctoral timeline (Q1--Q2 2026) \\
\midrule
\multicolumn{2}{@{}l}{\cellcolor{currentheader!15}\textbf{RGAIG Five-Layer Governance Alignment}} \\
\addlinespace
L1: Data Foundation & Survey data collection (online + paper), data screening, de-identification \\
L2: Analytics Engine & Descriptive stats, reliability, CFA, SEM path analysis \\
L3: Trust \& Accountability & Trust measurement (cognitive + affective), pre--post comparison, adoption barriers \\
L4: Decision Orchestration & Threshold-based evidence-maturity interpretation from aggregated trust scores \\
L5: Strategic Governance & Expert validation ($n = 12$, Krippendorff's $\alpha$), governance satisfaction assessment \\
\midrule
\multicolumn{2}{@{}l}{\cellcolor{currentheader!15}\textbf{AI Quality Dimensions Addressed}} \\
\addlinespace
Explainability & Clinician perception of SHAP/RAG explanation quality (C7--C8 survey items) \\
Responsibility & Informed consent, IRB compliance, respondent anonymity, data governance \\
Performance & Workflow readiness score, time-to-adoption, cognitive load reduction \\
Trust & Six-item trust construct ($\alpha = 0.91$), cognitive and affective trust dimensions \\
Ethics & Voluntary participation, right to withdraw, no clinical risk from survey \\
\end{xltabular}}
\vspace{0.4em}
\noindent\textit{Note.} Fifteen rows aggregated into four chapter-level framework-lens summaries in Table~\ref{tab:ch3_parta_smart}. SEM = structural equation modelling; CFA = confirmatory factor analysis; IRB = Institutional Review Board.

%% ======================================================================
%% C-S12. Part C SMART/RGAIG Alignment — Full Per-Row Breakdown
%% Condensation companion to Chapter 3 tab:ch3_partc_smart (4-row group
%% summary). Full 13-row 3-section decomposition migrated 2026-05-13
%% per main-chapter table-row reduction discipline.
%% ======================================================================
\addcontentsline{toc}{section}{C-S12. Part C SMART/RGAIG Alignment --- Full Per-Row Breakdown} % [TOC-appendix-section-insertion]
\section*{C-S12. Part C SMART/RGAIG Alignment --- Full Per-Row Breakdown}
\addcontentsline{toc}{section}{C-S12. Part C SMART/RGAIG Alignment}
\label{sec:appc_partc_smart_full}

\noindent This section provides the full per-row breakdown summarised in Chapter~3, Table~\ref{tab:ch3_partc_smart}. The four chapter-level framework-lens rows decompose into the original 13-row 3-section enumeration for Part~C (B2B hospital integration).

\noindent Table~\ref{tab:appc_partc_smart_full} summarises part c smart/rgaig framework alignment for appendix review.

{\tablestripes\tablefontsize

\noindent Table~\ref{tab:appc_partc_smart_full} below presents Part C SMART/RGAIG Alignment: Full Per-Row Breakdown, laying out each row in structured form to help examiners trace the source, applicability, and downstream use at a glance. % [auto-intro-strict inserted]
\paragraph{Part C Smart Rgaig Alignment.} % [auto-heading inserted]
\noindent Table~\ref{tab:appc_partc_smart_full} below presents Part C SMART/RGAIG Alignment: Full Per-Row Breakdown, laying out each row in structured form to help examiners trace the source, applicability, and downstream use at a glance. % [auto-intro-after-heading]



\needspace{20\baselineskip}
\begin{xltabular}{\textwidth}{@{} >{\raggedright\arraybackslash}p{2.0cm} L @{}}
\caption[Part C SMART/RGAIG Alignment: Full Per-Row Breakdown]{Part C SMART/RGAIG Framework Alignment: full per-row breakdown. Supporting evidence behind the 4-row decision-level summary in Table~\ref{tab:ch3_partc_smart}.}\label{tab:appc_partc_smart_full} \\
\toprule
\thead{Framework} & \thead{Part C Alignment} \\
\midrule
\endfirsthead
\endhead
\bottomrule
\endlastfoot
\multicolumn{2}{@{}l}{\cellcolor{currentheader!15}\textbf{SMART Goal Alignment}} \\
\addlinespace
Specific & Deploy ASD ensemble in hospital: 14-ch EEG, HL7 FHIR, clinician trust $\geq 3.5/5$, turnaround $\leq 30$\,min \\
Measurable & Accuracy, AUC, latency, turnaround, trust score, $\Delta$ pre--post, ROI, TCO \\
Achievable & Part~B pipeline validated at 97.67\%; clinician sample ($n \geq 30$) recruitable \\
Relevant & B2C accuracy (Part~B) must be validated in hospital context with clinician human-in-the-loop \\
Time-bound & Hospital pilot and clinician surveys within Q1--Q2 2026 \\
\midrule
\multicolumn{2}{@{}l}{\cellcolor{currentheader!15}\textbf{RGAIG Five-Layer Governance Alignment}} \\
\addlinespace
L1--L2 & Input validation: confirm Part~A and Part~B data quality before integration \\
L3: Trust & Triangulate technical trust (accuracy CI) with human trust (survey scores) \\
L4: Decision & Apply threshold logic: if accuracy $\geq 85\%$ AND trust $\geq 3.5/5$ AND governance pass $\geq 95\%$ $\to$ Adopt \\
L5: Governance & Map to NIST AI RMF, ISO/IEC 42001, CMM maturity assessment \\
\midrule
\multicolumn{2}{@{}l}{\cellcolor{currentheader!15}\textbf{Quality, KPI, and Decision Dimensions}} \\
\addlinespace
Quality KPIs & Technical: accuracy, AUC, latency. Human: trust, satisfaction, adoption. Operational: time saved, throughput \\
Illustrative financial scenario Model & TCO analysis, 3-year projected ROI, break-even timeline, sensitivity analysis (3 scenarios) \\
Compliance & HIPAA, FDA SaMD (Class~II), EU AI Act (Article~6), GDPR, HL7 FHIR \\
Decision Output & Adopt (all thresholds met) / Pilot (partial) / Defer (below threshold) per B2B and B2C \\
\end{xltabular}}
\vspace{0.4em}
\noindent\textit{Note.} Thirteen rows aggregated into four chapter-level framework-lens summaries in Table~\ref{tab:ch3_partc_smart}. FHIR = Fast Healthcare Interoperability Resources; HITL = human-in-the-loop; SaMD = Software as a Medical Device; TCO = total cost of ownership.

%% ======================================================================
%% C-S13. Master Bootstrap Resampling Protocol — Full Per-Parameter Breakdown
%% Condensation companion to Chapter 3 tab:ch3_partd_bootstrap (4-row
%% protocol-lens summary). Full 15-row 5-section decomposition migrated
%% 2026-05-13 per main-chapter table-row reduction discipline.
%% ======================================================================
\addcontentsline{toc}{section}{C-S13. Master Bootstrap Resampling Protocol --- Full Per-Parameter Breakdown} % [TOC-appendix-section-insertion]
\section*{C-S13. Master Bootstrap Resampling Protocol --- Full Per-Parameter Breakdown}
\addcontentsline{toc}{section}{C-S13. Master Bootstrap Resampling Protocol}
\label{sec:appc_partd_bootstrap_full}

\noindent This section provides the full per-parameter breakdown summarised in Chapter~3, Table~\ref{tab:ch3_partd_bootstrap}. The four chapter-level protocol-lens rows (Resampling Parameters, Multiple-Comparison Control, Stability \& Decision Mapping, Boundary Conditions) decompose into the original 15-row 5-section enumeration.

\noindent Table~\ref{tab:appc_partd_bootstrap_full} summarises master bootstrap resampling protocol for appendix review.

\clearpage
{\tablestripes\tablefontsize

\noindent Table~\ref{tab:appc_partd_bootstrap_full} below presents Master Bootstrap Resampling Protocol: Full Per-Parameter Breakdown, laying out each row in structured form to help examiners trace the source, applicability, and downstream use at a glance. % [auto-intro-strict inserted]
\paragraph{Master Bootstrap Resampling Protocol Full.} % [auto-heading inserted]
\noindent Table~\ref{tab:appc_partd_bootstrap_full} below presents Master Bootstrap Resampling Protocol: Full Per-Parameter Breakdown, laying out each row in structured form to help examiners trace the source, applicability, and downstream use at a glance. % [auto-intro-after-heading]



\needspace{20\baselineskip}
\begin{xltabular}{\textwidth}{@{} >{\raggedright\arraybackslash}p{3.0cm} L @{}}
\caption[Master Bootstrap Resampling Protocol: Full Per-Parameter Breakdown]{Master Bootstrap Resampling Protocol: full per-parameter breakdown. Supporting evidence behind the 4-row decision-level summary in Table~\ref{tab:ch3_partd_bootstrap}.}\label{tab:appc_partd_bootstrap_full} \\
\toprule
\thead{Parameter} & \thead{Specification} \\
\midrule
\endfirsthead
\endhead
\bottomrule
\endlastfoot
\multicolumn{2}{@{}l}{\cellcolor{currentheader!15}\textbf{Global Resampling Parameters}} \\
\addlinespace
Number of resamples & 1,000 (all Parts); sufficient for 95\% CI estimation per \cite{efron1993bootstrap} \\
Confidence level & 95\% (two-tailed); 99\% for Bonferroni-corrected comparisons \\
Random seed & 42 (fixed across all analyses for exact reproducibility) \\
Resampling method & Non-parametric bootstrap with replacement; bias-corrected and accelerated (BCa) CI \\
\midrule
\multicolumn{2}{@{}l}{\cellcolor{currentheader!15}\textbf{Correction for Multiple Comparisons}} \\
\addlinespace
Within-Part & Bonferroni correction: $\alpha' = 0.05 / k$ where $k$ = number of pairwise tests \\
Across-Parts & Family-wise error rate controlled at $\alpha = 0.05$ across all five hypotheses \\
\midrule
\multicolumn{2}{@{}l}{\cellcolor{currentheader!15}\textbf{Stability Interpretation Guide}} \\
\addlinespace
Stable & CI width $\leq$ threshold AND metric exceeds acceptance criterion in $\geq 95\%$ of resamples \\
Marginal & CI width within threshold BUT metric fails criterion in 5--20\% of resamples \\
Unstable & CI width exceeds threshold OR metric fails criterion in $> 20\%$ of resamples \\
\midrule
\multicolumn{2}{@{}l}{\cellcolor{currentheader!15}\textbf{Decision Mapping}} \\
\addlinespace
All metrics Stable & $\to$ \textbf{Adopt}: evaluation recommendation confirmed \\
$\geq 1$ metric Marginal & $\to$ \textbf{Pilot}: deploy with enhanced monitoring and re-evaluation \\
$\geq 1$ metric Unstable & $\to$ \textbf{Defer}: insufficient evidence; collect more data before evaluation \\
\midrule
\multicolumn{2}{@{}l}{\cellcolor{currentheader!15}\textbf{What Bootstrap Does NOT Replace}} \\
\addlinespace
External validation & Bootstrap tests internal stability, not generalisability to new populations \\
Prospective trial & Bootstrap on retrospective data cannot substitute for a clinical trial \\
Causal inference & Bootstrap CI does not establish causation; only association stability \\
\end{xltabular}}
\vspace{0.4em}
\noindent\textit{Note.} Fifteen rows aggregated into four chapter-level protocol-lens summaries in Table~\ref{tab:ch3_partd_bootstrap}. Group mapping: Global Resampling Parameters $\to$ Resampling Parameters lens; Correction for Multiple Comparisons $\to$ Multiple-Comparison Control lens; Stability Interpretation Guide + Decision Mapping $\to$ Stability \& Decision Mapping lens; What Bootstrap Does NOT Replace $\to$ Boundary Conditions lens. BCa = bias-corrected and accelerated; CI = confidence interval.
